\documentclass{ieeeaccess}
\usepackage{cite}
\usepackage{amsmath,amssymb,amsfonts}
\usepackage{algorithmic}
\usepackage{algorithm}
\usepackage{graphicx}
\usepackage{textcomp}
\usepackage{booktabs}
\usepackage{multirow}
\usepackage{subcaption}
\usepackage{bm}
\usepackage{float}
\usepackage{longtable}
\makeatletter
\AtBeginDocument{\DeclareMathVersion{bold}
\SetSymbolFont{operators}{bold}{T1}{times}{b}{n}
\SetSymbolFont{NewLetters}{bold}{T1}{times}{b}{it}
\SetMathAlphabet{\mathrm}{bold}{T1}{times}{b}{n}
\SetMathAlphabet{\mathit}{bold}{T1}{times}{b}{it}
\SetMathAlphabet{\mathbf}{bold}{T1}{times}{b}{n}
\SetMathAlphabet{\mathtt}{bold}{OT1}{pcr}{b}{n}
\SetSymbolFont{symbols}{bold}{OMS}{cmsy}{b}{n}
\renewcommand\boldmath{\@nomath\boldmath\mathversion{bold}}}
\makeatother

\def\BibTeX{{\rm B\kern-.05em{\sc i\kern-.025em b}\kern-.08em
    T\kern-.1667em\lower.7ex\hbox{E}\kern-.125emX}}

%Your document starts from here ___________________________________________________
\begin{document}
\doi{xxx/ACCESS.xxxx.xxxxxxx}

\title{Design and Evaluation of Multi-Agent Architectures
for Stock Price Prediction: A Vietnam Case Study}
\author{
    \uppercase{Duong Nguyen Minh}\authorrefmark{1}, 
    \uppercase{Dat Nguyen Huu}\authorrefmark{2}
}

\address[1]{School of Computing, Phenikaa University, Vietnam\\
(e-mail: 23010441@st.phenikaa-uni.edu.vn)}

\address[2]{School of Computing, Phenikaa University, Vietnam\\
(e-mail: dat.nguyenhuu@phenikaa-uni.edu.vn)}


\tfootnote{This paragraph of the first footnote will contain support
information, including sponsor and financial support acknowledgment. For
example, ``This work was supported in part by the U.S. Department of
Commerce under Grant BS123456.''}

\markboth
{Author \headeretal: Preparation of Papers for IEEE TRANSACTIONS and JOURNALS}
{Author \headeretal: Preparation of Papers for IEEE TRANSACTIONS and JOURNALS}

\corresp{Corresponding author: Dat Nguyen Huu (e-mail: dat.nguyenhuu@phenikaa-uni.edu.vn). Code and data are available at: \url{https://github.com/nminduo2k5/Multi-agentVNstock.git}}


\begin{abstract}
Traditional stock price prediction methods struggle in emerging markets characterized by retail investor dominance, high volatility, and sentiment-driven trading. We address this gap by developing a multi-agent framework tailored to Vietnam's unique market structure, where 85-90\% of participants are retail investors operating under a T+2.5 settlement cycle. Our system integrates six specialized agents that analyze technical indicators, fundamental metrics, risk profiles, news sentiment, macroeconomic signals, and market microstructure through three coordination mechanisms: Hierarchical, Round Robin, and Ensemble Voting. Empirical evaluation reveals strong horizon-dependent architectural effectiveness. For medium-term (14-day) forecasting, Ensemble Voting consistently achieves superior performance across multiple sectors, effectively reconciling heterogeneous technical and fundamental signals. Short-term (3-day) predictions exhibit stock-specific preferences, with Hierarchical architecture excelling in highly volatile sectors while Ensemble Voting proves optimal for stable growth-oriented stocks. Long-term (60-day) forecasting shows sector-dependent architectural advantages: Hierarchical coordination outperforms in highly regulated industries such as banking, while Round Robin demonstrates superiority in cyclical sectors like industrial construction. Across all horizons, the Multi-Agent framework significantly outperforms traditional LSTM baselines and single-agent approaches, offering a robust and interpretable decision-support tool for Vietnam's financial ecosystem.
\end{abstract}

\begin{keywords}
Architectures, Large Language Models, Multi-agent, Price Prediction, Vietnam Stock
\end{keywords}

\titlepgskip=-21pt

\maketitle

\section{Introduction}
\label{sec:introduction}
The rapid growth of digital technology in equity markets and the increasing availability of real-time market data  have created new opportunities for applying artificial intelligence (AI) in trading and investment analysis. In Vietnam’s stock market, this transformation is especially relevant. By the end of 2024, the market had experienced substantial growth in securities trading accounts, with the vast majority owned by domestic investors. Retail investors accounted for more than 80\% of total trading value \cite{vietnamnews2025}. These figures demonstrate that Vietnam’s market is heavily driven by retail investors and exhibits high volatility, structural inefficiencies, and an evolving regulatory framework.

Traditional methods of equity-market analysis, such as technical analysis based on price patterns and fundamental analysis rooted in corporate disclosures, often struggle to capture the complex, non-linear relationships found in emerging markets. For instance, empirical evidence from Vietnam has revealed sluggish information dissemination, intraday trading frictions, and pronounced influences from investor sentiment and speculative behavior \cite{ha2025}. These characteristics underscore the limitations of conventional modeling approaches and reinforce the necessity for more adaptive and robust analytical frameworks.

Integrating deep learning with multi-agent systems presents a promising approach to addressing these challenges. Deep learning models can automatically extract informative representations from large and complex datasets, thereby uncovering latent patterns and improving predictive accuracy. Simultaneously, multi-agent system designs enable the simulation of diverse market participants and their interactions, offering a more realistic representation of market dynamics and investor behavior.

To effectively deploy AI-based strategies within Vietnam's stock market, careful customization is essential to align with its specific operational and behavioral characteristics. However, the application of these advanced systems to Vietnam's unique market context remains underexplored. This gap motivates three critical research questions that guide this study:

\begin{itemize}
    
    \item \textbf{RQ1:} Under the constraints imposed by the $T+2.5$ settlement cycle, which coordination architecture (\textit{Hierarchical}, \textit{Round Robin}, or \textit{Ensemble Voting}) is most effective across short-term, medium-term, and long-term forecasting horizons?
    
    \item \textbf{RQ2:} How does the choice of the underlying Large Language Model (LLM), with a particular focus on \textit{Gemini 2.0 Flash} compared against alternative LLMs, affect the predictive stability and accuracy of the proposed agent-based system?

    \item \textbf{RQ3:} Does the proposed \textit{Multi-Agent} system empirically outperform traditional \textit{Single-Agent} baselines and standard deep learning models in an emerging market context?

\end{itemize}

Accordingly, we propose a tailored simulation framework that combines specialized agent modules and architectural designs to model Vietnamese stock market dynamics. The focus is on testing trading strategies, liquidity and volatility responses, and inter-agent communication under different coordination architectures. Through this approach, the research contributes both to the academic understanding of agent-to-agent interactions in emerging markets and to the practical design of decision-support tools for Vietnam’s financial ecosystem.



\section{Literature Review}
\label{sec:Literature Review}

\subsection{Stock Market Data}
Data from global stock exchanges have long been recognized as essential for academic research and institutional decision-making. Major trading platforms such as the NYSE, NASDAQ, LSE, and TSE operate at substantial scale, with combined market capitalizations exceeding USD 40–50 trillion and daily intraday trading volumes frequently surpassing hundreds of billions \cite{WorldFederation2024}. These markets provide millisecond-level order book updates, comprehensive limit order details (Level 3), and standardized high-frequency data streams suitable for automated trading systems and microsecond-level execution studies \cite{Jones2019HFT}. Such infrastructure has enabled detailed research on market microstructure, latency-based arbitrage strategies, and systemic risk propagation in high-frequency environments.

In comparison, Vietnam's stock market, comprising HOSE, HNX, and UPCoM, remains structurally smaller and data-constrained. The total market capitalization of listed companies in 2024 was approximately USD 250–300 billion \cite{HOSE2024Report}, with average daily trading volumes ranging between USD 800 million and USD 1.2 billion, significantly lower than those of regional peers such as Thailand and Malaysia \cite{ADB2024}. Moreover, available data are limited to daily high, low, opening, and closing prices, order flow summaries, and, more recently, limited order book visibility, without microsecond timestamps or complete order depth information \cite{Nguyen2022Microstructure}. 

A key behavioral divergence arises from the dominance of retail investors, who account for more than 85\% of total trading volume in Vietnam \cite{SSI2024Retail}. In contrast to institutionally driven markets, where prices typically respond to news and earnings reports, price movements in Vietnam are more heavily influenced by government policy announcements, informal communication channels, and sentiment-based speculation \cite{StateBank2023}. While global markets increasingly incorporate AI-curated news and alternative data into trading systems \cite{BloombergAI2023}, the Vietnamese market remains predominantly influenced by local social media sentiment and expectations regarding major government policy shifts.


\subsection{AI in Stock Price Prediction}

Over the past decades, AI-based methods have been widely applied to stock price prediction. Classical machine learning models such as SVMs, decision trees, random forests, and ANNs can capture non-linear relationships in financial data and often outperform traditional statistical approaches\cite{kim2003,huang2005,patel2015}. Deep learning models such as recurrent neural networks and long short-term memory networks further improved temporal analysis of financial time series \cite{fischer2018,chen2015}.

More recently, research has shifted toward \textit{Multi-Agent Systems} (MAS), in which agents interact in simulated markets. Unlike traditional models that treat prices as isolated data streams, MAS capture dynamic relationships among traders, brokers, and regulators, enabling the study of phenomena such as herd behavior, bubbles, and market irregularities \cite{lebaron2006,chen2012}. At a methodological level, coordination among distributed agents is strongly connected to consensus theory in networked multi-agent systems, which provides formal guarantees for information aggregation and agreement under dynamic interaction graphs \cite{olfati2007consensus,ren2005consensus}. Combining reinforcement learning with agent-based simulations also deepens the analysis of trading strategies, market stability, and inter-agent communication \cite{tesfatsion2006,wang2021}.

While the majority of MAS-based financial research has focused on mature markets such as the U.S. and China, empirical applications to emerging markets remain notably scarce. In the case of Vietnam, the AI-driven research currently available focuses mainly on price prediction trends over time via machine learning or deep learning techniques \cite{tran2020forecasting,pham2023lstm}, without considering how actions driven by individual investors, governmental policy shifts, or sudden shifts in market sentiment can shape market dynamics. This highlights a significant gap: it remains unproven whether Multi-Agent frameworks truly outperform traditional single-agent models (RQ3) in such volatile environments. Furthermore, while MAS is evolving, the specific impact of modern Large Language Models as the cognitive core of these agents, and how different models such as Gemini or GPT influence system performance (RQ2), have not been rigorously quantified in financial contexts.

Overall, this shift from single-model learning to multi-agent simulation reflects a growing recognition of stock markets as complex ecosystems shaped by interactions among diverse actors. This approach provides a useful means to examine regulatory effects and investor behavior in experimental settings, which is particularly relevant for developing economies such as Vietnam.

\subsubsection{Machine Learning for Stock Price Prediction}
The utilization of machine learning methodologies has been a long-standing practice in stock price prediction, particularly for capturing non-linear trends within financial time series data. An illustrative example is seen in Kim \cite{kim2003}, where support vector machines (SVM) were employed to predict stock market fluctuations, achieving a level of effectiveness that surpassed that of conventional statistical modeling techniques. Similarly, Huang et al. \cite{huang2005} utilized the capabilities of SVM to forecast the directional movement of stock indices, whereas Patel et al. \cite{patel2015} employed a hybrid approach combining artificial neural networks (ANN), random forests, and SVM to achieve heightened predictive precision. Progressing into more contemporary research, deep learning techniques have been increasingly adopted, notably long short-term memory (LSTM) models, which possess the inherent capacity to discern temporal interrelationships within price sequences \cite{fischer2018,chen2015}.

However, these pure data-driven models often fail to capture the complex interactions and behavioral dynamics of market participants. To address these challenges, researchers have proposed moving toward \textit{Multi-Agent Systems} (MAS). Multi-Agent Systems, in contrast to traditional machine learning, explicitly represent multiple agents, such as investors, brokers, and regulators, who interact under market rules and adapt their strategies over time. This enables the analysis of complex, emergent behaviors beyond prediction accuracy, offering a more comprehensive understanding of market stability, systemic risks, and the significance of agent communication \cite{lebaron2006,chen2012,tesfatsion2006,wang2021}. From this perspective, Multi-Agent Systems do not merely complement traditional machine learning, but extend it by embedding predictive intelligence into an interactive and behaviorally realistic market environment, a setup particularly well-suited to Vietnam’s retail-dominated and sentiment-driven stock market structure.

\subsubsection{Multi-Agent for Stock Price Prediction}
Initial applications of Multi-Agent Systems in finance were primarily aimed at replicating market behavior rather than predicting prices. Groundbreaking research by LeBaron \cite{lebaron2006} and Tesfatsion \cite{tesfatsion2006} demonstrated that interacting agents can recreate typical market phenomena, including high short-term volatility and speculative bubbles, which standard statistical analysis and single-model learning methods fail to adequately explain.
Subsequent research extended MAS into domains where agents could learn and improve their trading strategies through reinforcement learning. For example, Chen et al. \cite{chen2012} demonstrated how emergent market behaviors can arise from agent interactions, and Wang et al. \cite{wang2021} developed investment portfolios that outperformed those based on fixed, rule-based strategies.

Recent advancements in financial modeling have witnessed a paradigm shift toward hybrid learning-based architectures. In these sophisticated environments, automated agents do not merely function in isolation; they engage in complex, asynchronous communication within artificial market settings to simulate trading dynamics. Consequently, the research frontier has moved beyond merely validating the feasibility of MAS. The critical imperative now lies in identifying the optimal coordination topology. Specifically, we aim to determine which structural framework, whether the stratified logic of Hierarchical models, the cyclical consensus of Round Robin, or the democratic aggregation of Ensemble methods, is best equipped to navigate the idiosyncratic volatility of the Vietnamese stock market (RQ1). This framing is also consistent with consensus-oriented coordination principles in distributed multi-agent control \cite{olfati2007consensus,ren2005consensus}. Addressing this question is vital for balancing the conflicting signals of short-term retail speculation against long-term fundamental trends.

\section{Methodology}
\subsection{Architectural Overview}
The proposed Multi-Agent System (MAS) represents a sophisticated computational framework designed to address the unique challenges of stock price prediction in Vietnam's emerging market environment. The system integrates six specialized autonomous agents, each equipped with domain-specific analytical capabilities, operating within three distinct coordination architectures to deliver robust and interpretable financial forecasts. At its core, the system employs a foundation layer built upon \textbf{Long Short-Term Memory (LSTM)} networks to capture temporal dependencies in historical price data, enhanced through the coordinated intelligence of six specialized agents: \textbf{PricePredictor} (technical analysis), \textbf{InvestmentExpert} (fundamental evaluation), \textbf{RiskExpert} (risk assessment), \textbf{TickerNews} (company-specific sentiment), \textbf{MarketNews} (macroeconomic analysis), and \textbf{StockInfo} (market microstructure).

The cognitive intelligence of the system is powered by \textbf{Google Gemini 2.0 Flash} as the core Large Language Model (LLM), selected for its exceptional scalability, low-latency execution, and expansive context window capabilities essential for processing complex financial narratives and multi-dimensional market data. The system's architectural flexibility is demonstrated through three coordination architectures, namely \textbf{Hierarchical, Round Robin, and Ensemble Voting}, each specifically calibrated to handle Vietnam's market characteristics, including the T+2.5 settlement period, 85-90\% retail investor dominance, and pronounced sentiment-driven volatility, delivering predictions across short-term (3 days), medium-term (14 days), and long-term (60 days) investment horizons.

\begin{figure}[H]
    \centering
    \includegraphics[width=0.9\linewidth]{anh/overallmain.png}
    \caption{
System overview diagram}
    \label{fig:
System overview diagram}
\end{figure}

\subsection{The Specialized Agent Team}
\subsubsection{Agents Overview}
\noindent\textbf{Agent I/O specification and LLM usage.}  
In each prediction cycle, every agent receives structured JSON input and returns formatted JSON output for downstream aggregation. The LLM serves as a constrained reasoning layer (prompted with role, objective, and output schema), while numerical indicators are precomputed by deterministic modules.

\noindent\textit{Generic input schema:}
\begin{align}
X_t=\{ticker, horizon, date, price, \notag\\
fundamental, news, macro, micro\}.
\end{align}

\noindent\textit{Generic output schema of agent $i$:}
\begin{align}
Y_t^{(i)}=\{name, signal, confidence\in[0,1], rationale, flags\}.
\end{align}

\begin{itemize}
    \item \textbf{PricePredictor:} 
    Employs a two-stage forecasting pipeline integrating LSTM- and Transformer-based temporal modeling with technical indicator calibration. The agent leverages RSI for momentum, MACD for trend shifts, and Bollinger Bands for volatility envelopes. By synthesizing these technical signals with the neural baselines, it refines predictions across 3-day, 14-day, and 60-day horizons, delivering price forecasts accompanied by confidence scores and actionable trading parameters (entry zones, take-profit, stop-loss).
    \textit{Input:} price (OHLCV, LSTM/Transformer output, RSI, MACD, Bollinger, volatility).  
    \textit{Output:} pred\_price\_3d/14d/60d, technical\_confidence, signal, entry\_zone, take\_profit, stop\_loss, rationale.

    \begin{breakablealgorithm}
    \caption{PricePredictor pipeline}
    \begin{algorithmic}[1]
    \STATE \textbf{Stage 1 (Neural Temporal Modeling):} 
    \STATE \quad Extract historical sequence $P_{t-L:t}$ and compute baseline $\hat{y}_{base}(t,h)$ for $h \in \{3,14,60\}$ using Stacked LSTM or Transformer Encoder.
    \STATE \textbf{Stage 2 (Technical Calibration):}
    \STATE \quad Calculate technical adjustment $\Delta_{tech}$ by analyzing indicators:
    \STATE \quad \hspace{0.5em} $\delta_{RSI} \gets \pm 0.02$ based on overbought/oversold levels.
    \STATE \quad \hspace{0.5em} $\delta_{MACD} \gets \pm 0.01$ based on momentum divergence.
    \STATE \quad \hspace{0.5em} $\delta_{BB} \gets \pm 0.01$ based on volatility envelope position.
    \STATE \quad Aggregate: $\Delta_{tech} \gets \text{clamp}(\sum \delta_i, -0.05, 0.05)$.
    \STATE \quad Refine: $\hat{y}_{ref}(t,h) \gets \hat{y}_{base}(t,h) \times (1 + \alpha_h \Delta_{tech})$ where $\alpha_h \in \{1.0, 0.7, 0.4\}$.
    \STATE \textbf{Trading Parameter Synthesis:}
    \STATE \quad Determine $\{Signal, Entry, TP, SL\}$ based on $\hat{y}_{ref}$ and $BB$ levels.
    \end{algorithmic}
    \end{breakablealgorithm}
    \vspace{-1.5em}
    \begin{breakablealgorithm}
    \caption{Neural foundation modules}
    \begin{algorithmic}[1]
    \STATE \textbf{Stacked LSTM Architecture (Stage 1):}
    \STATE \quad Model: 2 $\times$ LSTM(units=50) + Dropout(0.2).
    \STATE \quad $\mathbf{h}_t \gets \text{LSTM}(\text{Scaler}(P_{t-L:t}); \Theta_{LSTM})$.
    \STATE \quad $\hat{y}_{t+h} \gets \text{Dense}(\mathbf{h}_t)$ via autoregressive iteration.
    \STATE \textbf{Transformer Encoder Architecture (Stage 1):}
    \STATE \quad Embedding: $X \gets \text{Proj}(P) + \text{PosEncoding}(L, d_{model})$.
    \STATE \quad Encoder ($N=2$): [Pre-LN MHA ($H=4$) + FFN].
    \STATE \quad Output: $\hat{y}_{t+h} \gets \text{Dense}(\text{SelectLastToken}(X))$.
    \end{algorithmic}
    \end{breakablealgorithm}

    \item \textbf{InvestmentExpert:} 
    Conducts comprehensive fundamental analysis by evaluating critical financial metrics including the Price-to-Earnings ratio (P/E), Price-to-Book ratio (P/B), Return on Equity (ROE), Earnings Per Share (EPS) growth trajectory, and Debt-to-Equity leverage ratios. Through systematic assessment of valuation multiples, profitability indicators, capital structure, and growth dynamics, the agent determines whether a stock is undervalued, fairly valued, or overvalued relative to its intrinsic worth and sector peers. The agent generates investment recommendations (BUY, SELL, or HOLD) accompanied by target price bias estimates and confidence scores. Importantly, the BUY/SELL recommendation is encoded into \texttt{target\_bias} as a directional fundamental stance and is used by the Master Agent as an auxiliary feature/consistency filter in the final price fusion, rather than being treated as an independent price forecast.  
    \textit{Input:} fundamental (valuation, profitability, leverage, growth).  
    \textit{Output:} signal (BUY/HOLD/SELL), target\_bias, confidence, rationale.

    \begin{breakablealgorithm}
    \caption{InvestmentExpert fundamental scoring}
    \begin{algorithmic}[1]
    \STATE \textbf{Component 1: Financial Health Assessment ($S_{fin}$):}
    \STATE \quad Evaluate $\{P/E, P/B, EPS, DivYield\}$ against adaptive sector benchmarks.
    \STATE \quad Compute $S_{fin}$ as a weighted average of individual metric scores.
    \STATE \textbf{Component 2: Valuation \& Growth Analysis ($S_{val}$):}
    \STATE \quad Analyze forward $P/E$ trends and dividend payout sustainability.
    \STATE \quad Determine intrinsic value bias based on mean-reversion expectations.
    \STATE \textbf{Component 3: Market Context Integration ($S_{ctx}$):}
    \STATE \quad Assess price position relative to 52-week bounds and volume dynamics.
    \STATE \textbf{Decision Synthesis:}
    \STATE \quad Aggregate: $S_{total} \gets 0.4 S_{fin} + 0.3 S_{val} + 0.3 S_{ctx}$.
    \STATE \quad Map $S_{total}$ to \texttt{target\_bias} $\in [-2, 2]$ and categorical \texttt{signal}.
    \end{algorithmic}
    \end{breakablealgorithm}

    \item \textbf{RiskExpert:} 
    Performs multi-dimensional risk assessment by analyzing realized volatility patterns, market correlation matrices, and historical drawdown profiles to quantify portfolio exposure and downside risk. The agent employs volatility modeling techniques to classify current market regimes (low/medium/high risk) and evaluates systematic risk through beta coefficients and correlation dynamics with broader market indices. By integrating Value-at-Risk (VaR) proxies and maximum drawdown metrics, the agent provides quantified risk levels accompanied by position sizing recommendations and risk flags that alert to elevated exposure conditions, enabling prudent capital allocation and risk-adjusted portfolio construction.  
    Importantly, the risk output is not interpreted as an independent price forecast; instead, \texttt{risk\_level} and \texttt{risk\_flags} modulate the final ensemble's confidence and allowable stance during master aggregation.
    \textit{Input:} risk (volatility, beta, drawdown, VaR).  
    \textit{Output:} risk\_level (LOW/MEDIUM/HIGH), position\_size, confidence, risk\_flags.

    \begin{breakablealgorithm}
    \caption{RiskExpert multi-dimensional assessment}
    \begin{algorithmic}[1]
    \STATE \textbf{Statistical Risk Profiling:}
    \STATE \quad Compute annualized volatility $\sigma$ and maximum drawdown $MDD$ from $P_{t-L:t}$.
    \STATE \quad Calculate Value-at-Risk ($VaR_{95}$) and market correlation $\rho_{mkt}$.
    \STATE \textbf{Systematic Risk Modeling:}
    \STATE \quad Determine sensitivity $\beta$ relative to the $VNIndex$ benchmark.
    \STATE \quad Classify risk level $\mathcal{R} \in \{L, M, H\}$ based on volatility thresholds.
    \STATE \textbf{Portfolio Constraint Synthesis:}
    \STATE \quad Map user risk tolerance $\tau$ to max position size $W_{max} \in [0.05, 0.20]$.
    \STATE \quad Derive dynamic stop-loss $\theta$ adjusted by investment horizon $h$.
    \STATE \textbf{Output Generation:}
    \STATE \quad Return $\{risk\_level, position\_size, confidence, risk\_flags\}$.
    \end{algorithmic}
    \end{breakablealgorithm}

    \item \textbf{TickerNews:} 
    Processes company-specific news streams sourced from leading Vietnamese financial platforms including CafeF, VietStock, and official earnings releases to extract sentiment signals and assess event-driven market impact. Employing finance-oriented sentiment classification algorithms, the agent analyzes headline sentiment, article content, and event tags while applying recency weighting to prioritize recent developments. The agent outputs sentiment polarity ratings (positive/negative/neutral) coupled with quantified impact scores that estimate the magnitude and direction of potential price movements, alongside identification of key events such as earnings surprises, management changes, or regulatory announcements that may drive near-term volatility.  
    Importantly, the sentiment/impact output is used as an auxiliary driver for the final price fusion (e.g., confidence calibration and consistency filtering), rather than being treated as a standalone price prediction.
    \textit{Input:} news (headlines, snippets, event tags, recency weights).  
    \textit{Output:} sentiment (POS/NEG/NEU), impact\_score, confidence, key\_events.

    \begin{breakablealgorithm}
    \caption{TickerNews sentiment and event analysis}
    \begin{algorithmic}[1]
    \STATE \textbf{Sentiment Scoring ($S_{sent}$):}
    \STATE \quad Tokenize news headlines/summaries and map to weight matrix.
    \STATE \quad $S_{sent} \gets \text{clamp}(\sum \text{weight}(\text{token}), -10, 10)$.
    \STATE \textbf{Recency Weighting:}
    \STATE \quad Apply temporal decay $w_j = e^{-age_j/30}$ to news item $j$.
    \STATE \quad Compute weighted average: $\bar{S} = (\sum S_{sent,j} w_j) / \sum w_j$.
    \STATE \textbf{Event Impact Assessment:}
    \STATE \quad Detect events $E \in \{Earnings, M\&A, Crisis\}$ via keyword matching.
    \STATE \quad $Impact \gets \bar{S} + \text{EventBoost}(E)$.
    \STATE \textbf{Output Generation:}
    \STATE \quad Return $\{sentiment, impact\_score, confidence, key\_events\}$.
    \end{algorithmic}
    \end{breakablealgorithm}

    \item \textbf{MarketNews:} 
    Examines macroeconomic updates, market sentiment indicators, and central bank policy communications to classify prevailing macro regimes and assess systemic risk conditions. The agent monitors interest rate signals, inflation trends (CPI), economic growth indicators (GDP), policy announcements from the State Bank of Vietnam, and broader index sentiment to determine whether markets are operating in risk-on (growth-seeking), risk-off (defensive), or neutral regimes. Through systemic risk evaluation and regime classification, the agent generates contextual market ratings that reflect the directional pressure from macroeconomic forces, enabling portfolio positioning aligned with the broader economic cycle and monetary policy stance.  
    Importantly, macro outputs (\texttt{macro\_regime} and \texttt{market\_pressure}) function as regime-level auxiliary context that modulates the confidence and weighting of other expert signals during aggregation, rather than acting as an independent price forecast.
    \textit{Input:} macro (interest rates, CPI/GDP, policy headlines, index sentiment).  
    \textit{Output:} macro\_regime (RISK-ON/RISK-OFF/NEUTRAL), market\_pressure, confidence.

    \begin{breakablealgorithm}
    \caption{MarketNews macro regime classification}
    \begin{algorithmic}[1]
    \STATE \textbf{Regime Classification ($\mathcal{R}$):}
    \STATE \quad Analyze macro-corpus for $\{RISK\text{-}ON, RISK\text{-}OFF\}$ keyword clusters.
    \STATE \quad $S_{net} \gets \text{Count}(K_{ON}) - \text{Count}(K_{OFF})$.
    \STATE \quad $\mathcal{R} \gets \text{sign}(S_{net})$ mapped to categorical regimes.
    \STATE \textbf{Market Pressure Mapping ($\Psi$):}
    \STATE \quad Compute $\Psi \in [-1.0, 1.0]$ based on sentiment intensity and policy signals.
    \STATE \textbf{Systemic Signal Extraction:}
    \STATE \quad Identify $\{NHNN\_stance, CPI, GDP\}$ from policy reports.
    \STATE \quad Detect sector rotation (leading vs. lagging industries).
    \STATE \textbf{Output Generation:}
    \STATE \quad Return $\{macro\_regime, market\_pressure, confidence, macro\_factors\}$.
    \end{algorithmic}
    \end{breakablealgorithm}

    \item \textbf{StockInfo:} 
    Analyzes market microstructure by examining trading volume patterns, liquidity trends, bid-ask spread dynamics, and order flow imbalances to assess execution feasibility and transaction cost implications. The agent evaluates abnormal volume spikes, turnover ratios, and order book depth to determine whether current market conditions support efficient trade execution or pose elevated slippage risk. By integrating these microstructure signals, the agent presents a liquidity-adjusted feasibility score accompanied by tactical execution recommendations (aggressive/passive/wait) that guide optimal trade deployment timing and order placement strategies, ensuring that forecasted opportunities can be practically captured within Vietnam's T+2.5 settlement framework.  
    Importantly, the feasibility score and execution advice are used to filter the final decision notes and risk-adjusted confidence (i.e., practicality of execution), rather than serving as an independent price forecast.
    \textit{Input:} micro (turnover, spread, volume, imbalance).  
    \textit{Output:} feasibility\_score, execution\_advice (AGGRESSIVE/PASSIVE/WAIT), confidence.

    \begin{breakablealgorithm}
    \caption{StockInfo microstructure analysis}
    \begin{algorithmic}[1]
    \STATE \textbf{Composite Liquidity Scoring ($S_{liq}$):}
    \STATE \quad $S_{vol} \gets$ Relative volume vs. 20-day/60-day moving averages.
    \STATE \quad $S_{spr} \gets$ Bid-ask spread percentage and transaction cost proxies.
    \STATE \quad $S_{dep} \gets$ Depth-based imbalance and total order book value.
    \STATE \quad $S_{turn} \gets$ Turnover ratio relative to market capitalization.
    \STATE \textbf{Execution Time Estimation ($S_{exec}$):}
    \STATE \quad Estimate days to execute target order size at participation rate $\gamma$.
    \STATE \textbf{Practical Feasibility Mapping:}
    \STATE \quad $S_{liq} \gets w_{vol}S_{vol} + w_{spr}S_{spr} + w_{dep}S_{dep} + w_{turn}S_{turn} + w_{exec}S_{exec}$.
    \STATE \quad Map $S_{liq}$ to \texttt{execution\_advice} $\in \{AGGRESSIVE, PASSIVE, WAIT\}$.
    \STATE \textbf{Output Generation:}
    \STATE \quad Return $\{feasibility\_score, execution\_advice, confidence\}$.
    \end{algorithmic}
    \end{breakablealgorithm}
\end{itemize}

\noindent\textbf{Illustrative example (14-day, VCB).}  
At time $t$, the orchestrator builds $X_t$ for ticker=VCB, horizon=14. Six agents produce $\{Y_t^{(i)}\}_{i=1}^{6}$ (e.g., technical: mildly bullish, fundamentals: neutral-positive, risk: medium, ticker news: neutral, macro: risk-off, microstructure: liquid). The coordination layer fuses outputs (Hierarchical/Round Robin/Ensemble) to generate final\_pred\_price, final\_signal, overall\_confidence, and decision\_notes for backtesting.

\subsubsection{LLMs Overview}
\textbf{Google Gemini 2.0 Flash} is selected as the core LLM due to its high scalability, low-latency cloud execution, and exceptionally large context window, making it well-suited for production-grade AI pipelines. To provide a broader comparative perspective, we additionally include \textbf{OpenAI GPT-4o} and \textbf{Llama-3.1-8B}. GPT-4o is recognized for its strong reasoning quality and balanced multimodal performance, whereas Llama-3.1 offers significant advantages for privacy-focused and offline deployments. The operational characteristics of these models are summarized in Table~\ref{tab:llm_comparison_summary_en}.

\begin{table}[h!]
\centering
\caption{Operational Characteristics and Execution Behavior of Leading Large Language Models.}
\label{tab:llm_comparison_summary_en}
\resizebox{\columnwidth}{!}{
\begin{tabular}{@{}lccc@{}}
\toprule
\textbf{Criterion} & \textbf{Gemini 2.0 Flash} & \textbf{GPT-4o} & \textbf{Llama-3.1-8B} \\ \midrule
Deployment & Cloud API & Cloud API & Local / Hosted \\
Latency & Low & Medium & Low / Variable \\
Throughput & Very High & High & Infra-dependent \\
Offline Capability & None & None & Yes \\
Context Window & Up to 1M tokens & Up to 128k tokens & Model-dependent \\
Cost Model & SaaS (usage-based) & SaaS (usage-based) & CapEx / SaaS \\
Optimized for & Cloud scale & Reasoning quality & Privacy / Control \\ 
\bottomrule
\end{tabular}
}
\end{table}

\subsection{Coordination Architectures}
\subsubsection{Hierarchical Architecture (Design 1)}

\paragraph{Operational Mechanism}
The hierarchical architecture adopts a Master--Subordinate design. A centralized \textit{Master Agent} (i) orchestrates the workflow, (ii) dispatches a common context to six specialized agents, and (iii) aggregates their outputs into a single coherent market forecast.

Concretely, the Master Agent coordinates six agents: \textit{PricePredictor} (technical and LSTM-based forecasting), \textit{InvestmentExpert} (fundamental valuation signals), \textit{RiskExpert} (risk regime and portfolio constraints), \textit{TickerNews} (company-level sentiment and event impact), \textit{MarketNews} (macro regime and systemic pressure), and \textit{StockInfo} (liquidity and execution feasibility). The computationally expensive LSTM baseline is computed once as the shared backbone output and is included in the context of each agent (as part of \texttt{price\_feat}), while each specialized agent transforms domain-relevant evidence into a structured output package \{$Y_t^{(i)}$\} using its own reasoning logic.

All agents operate independently and can run in parallel. Communication follows an asynchronous send--receive pattern: the Master Agent dispatches inputs, waits for returned outputs, and continues aggregation only after validation succeeds. This improves efficiency while preserving modular expertise and robustness to partial failures.

After receiving all agent outputs, the Master Agent executes (1) schema/range validation, (2) strict consistency checks across directional signals and risk flags, (3) conflict resolution via reliability-aware down-weighting, and (4) a weighted Bayesian-style aggregation that interprets the dynamic weights as posterior model probabilities. The final unified prediction is produced together with an overall confidence score and a decision record, ensuring traceability of how heterogeneous signals jointly determine the market view.

This hierarchical coordination provides a computationally efficient yet explainable fusion layer: specialization improves each agent’s local signal quality, while centralized aggregation reduces cross-domain inconsistencies and stabilizes the final decision under market regime changes~\cite{tesfatsion2006,lebaron2006}.

\begin{figure}[H]
    \centering
    \includegraphics[width=0.9\linewidth]{anh/hiemain.png}
    \caption{Hierarchical architecture where a Master Agent coordinates six specialized agents; agents process domain-specific data in parallel, and the Master Agent aggregates outputs for final prediction.}
    \label{fig:hierarchical_architecture}
\end{figure}


\paragraph{Workflow}
The hierarchical architecture operates through a structured sequence of stages that makes the parallelism and aggregation mechanism explicit:

\begin{enumerate}
    \item \textbf{Data Acquisition and Preprocessing:} 
    The system collects and preprocesses multi-source financial data, including technical indicators, fundamental ratios, and market sentiment signals.
    
    \item \textbf{LSTM Foundation Modeling:} 
    An LSTM-based model generates a baseline forecast $\{P_{\text{LSTM}}(t,h)\}_{h\in\{3,14,60\}}$, which is included in the shared context.
    \textit{Clarification:} The LSTM serves as a foundation module shared by all specialized agents; it is not treated as an independent agent in the hierarchical coordination layer.
    
    \item \textbf{Parallel Agent Processing:} 
    Six specialized agents are executed in parallel. Each agent receives the same context packet $X_t$ (ticker, horizon, indicators, and the shared LSTM baseline) and returns its structured output $Y_t^{(i)}$.
    
    \item \textbf{Validation, Conflict Resolution, and Aggregation:} 
    The Master Agent validates each output package, detects inconsistencies (e.g., directional conflict between technical/fundamental/sentiment signals or risk flags incompatible with the forecast stance), and resolves them by reliability-aware reweighting. In this process, all non-price agents (\textit{InvestmentExpert} via \texttt{target\_bias}, \textit{RiskExpert} via \texttt{risk\_level}/\texttt{risk\_flags}, \textit{TickerNews} via \texttt{sentiment}/\texttt{impact\_score}, \textit{MarketNews} via \texttt{macro\_regime}/\texttt{market\_pressure}, and \textit{StockInfo} via \texttt{feasibility\_score}) are treated as auxiliary context signals that modulate confidence and consistency checks during the final price fusion, rather than being interpreted as independent price forecasts. The master then applies weighted fusion of the expert-adjusted forecasts using dynamic weights to obtain a unified prediction and an overall confidence score.
    
    \item \textbf{Final Output Generation:} 
    The system produces the unified result, including final price forecasts, confidence levels, and investment recommendations.
\end{enumerate}

    \begin{breakablealgorithm}
    \caption{Master Agent aggregation}
    \begin{algorithmic}[1]
    \STATE \textbf{Input:} context packet $X_t$ and recent performance cache for dynamic weights.
    \STATE \textbf{Compute:} baseline forecasts $\{P_{\text{LSTM}}(t,h)\}_{h\in\{3,14,60\}}$.
    \STATE \textbf{Dispatch:} send $X_t$ (including $P_{\text{LSTM}}$) to all specialized agents in parallel.
    \STATE \textbf{Collect:} receive $Y_t^{(i)}$ from each agent $i=1,\dots,6$.
    \STATE \textbf{Validate:} check schema validity, confidence range, and missing-field handling.
    \STATE \textbf{Detect conflicts:}
      \STATE \hspace{0.8em}compare directional signals and risk flags across agents; mark conflict indicators.
    \STATE \textbf{Resolve conflicts:}
      \STATE \hspace{0.8em}compute dynamic weights $w_i(t)$ from recent reliability; down-weight inconsistent experts.
    \STATE \textbf{Fuse:} aggregate expert-adjusted forecasts to obtain $\hat y_{\text{ens}}(t+\Delta|t)$ and $\sigma_{\text{ens}}^2(t+\Delta|t)$.
    \STATE \textbf{Output:} final prediction, overall confidence, and decision notes for backtesting.
    \end{algorithmic}
    \end{breakablealgorithm}

\subsubsection{Round Robin Architecture (Design 2)}
\paragraph{Operational Mechanism}
The Round Robin architecture follows a sequential refinement process, in which each agent receives the output from the previous stage, improves it through its specialized analysis, and passes it forward. This iterative chain allows progressive development of predictions via multi-perspective evaluation \cite{wang2021,tesfatsion2006,lebaron2006,chen2012}. 


\begin{figure}[H]
    \centering
    \includegraphics[width=0.9\linewidth]{anh/rrmain.png}
    \caption{Round Robin architecture with six specialized agents operating in a sequential refinement cycle.}
    \label{fig:hierarchical_architecture}
\end{figure}



\paragraph{Workflow} The Round Robin architecture operates through an organized sequence of distinct stages:

\begin{enumerate}
    \item \textbf{Initialization.} The approach begins with the  \textit{LSTM Foundation Layer}, which constructs the initial predictive state by capturing temporal dependencies in the historical data. This state serves as the baseline representation of the market dynamics before any domain-specific refinement is applied.  

    \item \textbf{Sequential Agent Flow.} The predictive state is propagated through six specialized agents in a fixed order: Round~1 $\rightarrow$ Round~2 $\rightarrow$ Round~3 $\rightarrow$ Round~4 $\rightarrow$ Round~5 $\rightarrow$ Round~6. Each agent operates as a refinement node, sequentially receiving the previous output, performing its own analysis, and forwarding an improved result to the subsequent stage.  

    \item \textbf{Incremental Enhancement.} Within this Round Robin structure, each of the six agents contributes a unique perspective: \textit{PricePredictor} (technical analysis), \textit{InvestmentExpert} (fundamental metrics), \textit{RiskExpert} (risk assessment), \textit{TickerNews} (company-specific sentiment), \textit{MarketNews} (macroeconomic signals), and \textit{StockInfo} (market microstructure). This iterative enhancement ensures that both micro-level indicators and macro-level contexts are systematically embedded into the prediction pipeline.  

    \item \textbf{Decision Traceability.} Throughout the refinement sequence, all intermediate states and analytical decisions are logged to preserve a transparent audit trail. This traceability allows researchers and practitioners to examine how individual agents influenced the outcome, thereby supporting interpretability and trust in model-driven recommendations.  

    \item \textbf{Cumulative Refinement Output.} Upon completion of the six-round iteration, the framework consolidates the results into a final state representing the cumulative learning and multi-perspective reasoning of the entire agent ensemble. The resulting output thus reflects not only predictive accuracy but also contextual awareness and adaptive decision integration across the full refinement cycle.  
\end{enumerate}

\subsubsection{Ensemble Voting Architecture (Design 3)}
\paragraph{Operational Mechanism}
The Ensemble Voting Architecture employs an intelligent voting system in which six specialized agents operate concurrently and independently. Each agent conducts focused analysis on different data domains, ensuring that a wide range of predictive perspectives are considered. Subsequently, the system employs a statistical inference engine to combine these findings. This engine utilizes dynamic weighting based on recent performance, execution time, and confidence scores. This versatile strategy enhances robustness and reduces bias, facilitating more reliable and interpretable collective predictions in financial forecasting~\cite{li2022}.

\begin{figure}[H]
    \centering
    \includegraphics[width=0.9\linewidth,height=6cm]{anh/emmain.png}
    \caption{Ensemble Voting Architecture with six specialized agents employing an intelligent voting system}
    \label{fig:hierarchical_architecture}
\end{figure}


\paragraph{Workflow}
The operational flow of the proposed Ensemble Voting Architecture is summarized as follows:

\begin{enumerate}
    \item \textbf{Data Distribution:} Input data streams are concurrently distributed to all six agents.
    \item \textbf{Parallel Processing:} Each agent operates independently in parallel, ensuring full task isolation.
    \item \textbf{Performance Evaluation:} Real-time assessment of individual agent performance is continuously conducted.
    \item \textbf{Dynamic Weight Adjustment:} Model weights are adaptively updated based on the latest performance metrics.
    \item \textbf{Statistical Inference Integration:} Final aggregation is performed through statistical inference with uncertainty quantification to enhance predictive reliability.
\end{enumerate}




\paragraph{Dynamic Weight Adjustment Process}
To maintain adaptability in volatile market regimes, we formalize dynamic weighting over a rolling window of length $H$. Let $i \in \{1,\dots,N\}$ denote an agent, and let $e_i(\tau)$, $c_i(\tau)$, and $\ell_i(\tau)$ be its error, confidence-calibration loss, and execution latency at time $\tau$. The recent performance score is defined as
\begin{align}
S_i(t) &= \frac{1}{H}\sum_{\tau=t-H+1}^{t}
\Big(
\alpha\,\tilde e_i(\tau) + \beta\,\tilde c_i(\tau) + \gamma\,\tilde \ell_i(\tau)
\Big), \nonumber\\
&\quad \alpha+\beta+\gamma=1,
\label{eq:agent_score}
\end{align}
where $\tilde e_i,\tilde c_i,\tilde \ell_i \in [0,1]$ are min-max normalized losses (lower is better), and $\alpha,\beta,\gamma$ control the contribution of predictive error, calibration quality, and latency, respectively.

We then convert scores into unnormalized utilities via a temperature-controlled exponential mapping:
\begin{align}
u_i(t) = \exp\!\big(-\eta\,S_i(t)\big), \quad \eta>0,
\label{eq:agent_utility}
\end{align}
and obtain normalized dynamic weights
\begin{align}
w_i(t)=\frac{u_i(t)}{\sum_{j=1}^{N}u_j(t)}, \qquad \sum_{i=1}^{N}w_i(t)=1.
\label{eq:dynamic_weight}
\end{align}

To reduce high-frequency oscillation, a real-time smoothing update is applied:
\begin{align}
w_i^{\star}(t)=\lambda\,w_i^{\star}(t-1)+(1-\lambda)\,w_i(t), \quad \lambda\in[0,1).
\label{eq:smoothed_weight}
\end{align}
Equations \eqref{eq:agent_score}--\eqref{eq:smoothed_weight} ensure that agents with better recent efficiency receive larger influence while preserving ensemble stability under regime shifts\cite{McAndrew2021,Lu2023}.
This formulation is a study-specific mathematical instantiation of adaptive ensemble weighting, designed for our multi-agent financial setting and grounded in prior adaptive weighting and ensemble forecasting literature\cite{McAndrew2021,Lu2023,li2022}.

\paragraph{Statistical Inference Process}
The ensemble layer integrates all agent outputs through a statistically grounded inference pipeline. Let $\hat y_i(t+\Delta|t)$ be the $\Delta$-step-ahead forecast of agent $i$. The aggregated forecast is
\begin{align}
\hat y_{\text{ens}}(t+\Delta|t)=\sum_{i=1}^{N}w_i^{\star}(t)\,\hat y_i(t+\Delta|t).
\label{eq:ensemble_forecast}
\end{align}

Uncertainty is estimated by weighted dispersion around the ensemble mean:
\begin{align}
\sigma_{\text{ens}}^{2}(t+\Delta|t)=
\sum_{i=1}^{N}w_i^{\star}(t)\Big(\hat y_i(t+\Delta|t)-\hat y_{\text{ens}}(t+\Delta|t)\Big)^2.
\label{eq:ensemble_variance}
\end{align}

Inter-agent consensus is quantified as
\begin{align}
\mathcal{C}(t)=1-\frac{1}{N(N-1)}\!\sum_{i\neq j}
\frac{\left|\hat y_i(t+\Delta|t)-\hat y_j(t+\Delta|t)\right|}
{\left|\hat y_{\text{ens}}(t+\Delta|t)\right|+\varepsilon},
\label{eq:consensus}
\end{align}
where $\varepsilon>0$ avoids division instability; larger $\mathcal{C}(t)$ indicates stronger agreement.

Finally, statistical validation is conducted by monitoring calibration and residual diagnostics of $r(t+\Delta)=y(t+\Delta)-\hat y_{\text{ens}}(t+\Delta|t)$ over the rolling horizon, ensuring internal coherence and predictive robustness in empirical settings\cite{wang2021,li2022}. In combination, \eqref{eq:dynamic_weight}--\eqref{eq:consensus} act as an adaptive filter that down-weights temporarily underperforming agents and preserves stability under abnormal market fluctuations.
Accordingly, Eqs. \eqref{eq:ensemble_forecast}--\eqref{eq:consensus} should be interpreted as a formalized integration layer tailored to this paper, rather than a verbatim reproduction of any single prior model\cite{wang2021,li2022}.






\section{Empirical Evaluation}

This section evaluates the proposed Multi-Agent forecasting framework under Vietnam's $T+2.5$ settlement cycle constraints, systematically addressing research questions RQ1–RQ3 as formulated in the Introduction.

\subsection{RQ1: Architectural Effectiveness Across Horizons}

\textbf{Research Question:} Under the constraints imposed by the $T+2.5$ settlement cycle, which coordination architecture (Hierarchical, Round Robin, or Ensemble Voting) is most effective across short-term, medium-term, and long-term forecasting horizons?


\subsection{RQ2: Impact of Foundation LLMs}

\textbf{Research Question:} How does the choice of underlying LLM, with particular focus on Gemini 2.0 Flash compared against alternative LLMs, affect predictive stability and accuracy?

\subsection{RQ3: Multi-Agent vs. Traditional Baselines}

\textbf{Research Question:} Does the proposed Multi-Agent system empirically outperform traditional Single-Agent baselines and standard deep learning models in an emerging market context?



\subsection{Dataset \& Settings}
A fundamental structural characteristic of the Vietnamese stock market is its \textbf{T+2.5 settlement period}. Given that \textit{domestic individual investors overwhelmingly dominate the market, accounting for approximately 85-90\% of daily trading value} \cite{TinNhanhChungkhoan2025}, this configuration strongly incentivizes rapid, short-term speculative trading. This retail-heavy structure, often characterized by significant \textit{herd behavior} rather than fundamental analysis, is identified by analysts as a primary driver of the market's notoriously high price volatility \cite{WorldBank2024, Yuanta2025}.



Inspired by the substantial effects of this behavioral pattern, the present study conducts a sector-specific analysis by selecting representative stocks from four key sectors of the Vietnamese economy: \textbf{Banking (VCB)}, \textbf{Technology (FPT)}, \textbf{Real Estate (DXG)}, and \textbf{Industrial Construction (VCS)}. These sectors are chosen to capture the structural diversity of the Vietnamese stock market in terms of regulatory intensity, growth potential, and sensitivity to macroeconomic conditions.

VCB represents the banking sector due to its large market capitalization, high liquidity, and systemic importance within Vietnam’s financial system, making it highly sensitive to monetary policy and regulatory changes. FPT is selected as a proxy for the technology sector, where price movements are largely driven by innovation, earnings growth, and investor sentiment. DXG reflects the real estate sector, which is characterized by strong cyclicality and elevated volatility arising from credit conditions and macroeconomic fluctuations. Finally, VCS represents the industrial construction sector, whose stock performance is closely linked to infrastructure investment, export demand, and long-term development cycles, typically exhibiting more stable yet trend-driven behavior.

This diversified sectoral selection enables a comprehensive evaluation of forecasting performance across heterogeneous market dynamics.

Historical price data and trading volumes are sourced from both the Ho Chi Minh Stock Exchange (HOSE) and the Hanoi Stock Exchange (HNX) via multiple data providers including VNDirect Securities Corporation's public API, vnstock Python library, and SSI iBoard platform. Daily closing prices, intraday trading volumes, and order book data are collected to support technical analysis. Macroeconomic indicators, monetary policy announcements, and company-specific news are aggregated from official financial news platforms including CafeF, VietStock, SSI Research, and Vndirect Research. Fundamental data such as financial statements, earnings reports, and valuation metrics are obtained from company disclosures and third-party financial databases.

To rigorously evaluate predictive performance across heterogeneous investment horizons, the 
analysis is structured into three temporal regimes:
\begin{itemize}
    \item \textbf{Short-term (3 days):} aligned with the T+2.5 settlement window, capturing immediate speculative turnover.
    \item \textbf{Medium-term (14 days):} representing a conventional tactical positioning horizon.
    \item \textbf{Long-term (60 days):} proxying strategic buy-and-hold investment behavior.
\end{itemize}

To ensure robustness and reduce temporal bias, the evaluation at each forecasting horizon (\textbf{Short-term}, \textbf{Medium-term}, \textbf{Long-term}) is conducted across four distinct market intervals:
\begin{enumerate}
    \item \textbf{Late August to early September}: capturing post-summer liquidity shifts;
    \item \textbf{16/10/2025}: representing a period of heightened trading activity;
    \item \textbf{21/10/2025}: corresponding to a stabilization phase after short-term volatility;
    \item \textbf{22/10/2025}: reflecting market behavior during an immediate rebound window.
\end{enumerate}
Assessing performance across these four checkpoints enables a comprehensive examination of model behavior under varying liquidity, volatility, and sentiment conditions, thereby strengthening the validity of comparative results across all three time horizons.


\subsection{Evaluation Metrics.} 
To rigorously evaluate the predictive performance of the proposed framework, we utilize three established statistical indicators: Mean Absolute Error (MAE), Root Mean Square Error (RMSE), and Mean Absolute Percentage Error (MAPE). MAE quantifies the average magnitude of errors in a set of forecasts without considering their direction, providing a linear score that weights all individual differences equally~\cite{willmott2005advantages}. Conversely, RMSE operates as a quadratic scoring rule; by squaring residuals prior to averaging, it assigns greater weight to larger deviations, thereby increasing sensitivity to outliers~\cite{chai2014root}. Finally, MAPE measures accuracy as a relative percentage, facilitating the comparison of forecasting performance across datasets with differing scales~\cite{hyndman2006another}.

\subsection{Practical Trading Simulation (T+2.5 Constraints)}
To address the practical feasibility concern, we complement forecast-error evaluation with a simplified trading backtest that explicitly accounts for (i) transaction costs, (ii) slippage driven by liquidity conditions, and (iii) the Vietnam settlement cycle $T+2.5$~\cite{TinNhanhChungkhoan2025}.

\paragraph{Trading policy.}
For each forecasting origin $t$ and horizon $h \in \{3,14,60\}$, the system outputs a final price forecast and an associated trading stance (BUY/HOLD/SELL) from the forecasting module. We implement a long-only policy for practicality: when the stance is BUY, the system opens/maintains a long position; when it is SELL, it closes the position at the next executable time step; HOLD leaves the position unchanged. Position sizing uses the risk-oriented output \texttt{position\_size} from \textit{RiskExpert} when available, otherwise defaults to a fixed fraction of capital.

\paragraph{Execution price with costs and slippage.}
Let $P_{\text{exec}}(t)$ denote the daily execution reference price. Transaction costs are modeled with a proportional fee rate $f$ (brokerage + statutory charges). Slippage is approximated as a liquidity-dependent perturbation $s(t)$ derived from the feasibility score output of \textit{StockInfo}. The effective buy/sell execution prices are modeled as:
\begin{align}
P_{\text{buy}}(t) &= P_{\text{exec}}(t)\big(1+f+s(t)\big),\\
P_{\text{sell}}(t) &= P_{\text{exec}}(t)\big(1-f-s(t)\big),
\end{align}
where $s(t)$ is consistent with execution-cost modeling principles in the optimal execution literature~\cite{almgren2001optimal}.

\paragraph{T+2.5 settlement constraint.}
Orders submitted at day $t$ are executed at the next trading day, while cash and positions become reusable only after settlement. In the discrete daily setting, we approximate the $T+2.5$ cycle with a settlement delay of $d_s$ trading days. The simulation maintains a cash-lock mechanism: after a BUY, the corresponding notional plus estimated fees are locked until settlement at $t+d_s$; similarly, a SELL closing action releases cash only after settlement. This constraint prevents unrealistic rapid turnover in the backtest.

\paragraph{Backtest metrics.}
For each episode $(\text{stock},\text{horizon},\text{interval})$, we compute net trade returns and portfolio-level performance under costs:
\begin{align}
r_{\text{trade}} &= \frac{P_{\text{sell}}(t+h)-P_{\text{buy}}(t)}{P_{\text{buy}}(t)},\\
\text{DD} &= \max_{\tau}\Big(\max_{u\le \tau} V(u) - V(\tau)\Big), 
\end{align}
where $V(\tau)$ is the simulated portfolio value. Max drawdown (DD) is used as a standard risk measure capturing the largest peak-to-trough decline in portfolio value over the backtest window~\cite{magdon2004maxdrawdown}. We report cumulative net return, max drawdown (DD), and cost-adjusted directional accuracy (the fraction of forecast-driven stances yielding non-negative net returns after frictions).

\paragraph{Simulation procedure (pseudo-code).}
\begin{algorithmic}[1]
\STATE \textbf{Input:} historical prices/volumes, feasibility and risk outputs, fee rate $f$, settlement delay $d_s$, horizons $h\in\{3,14,60\}$.
\STATE \textbf{For each interval and each origin } $t$ \textbf{in the test window:}
\STATE \hspace{0.8em}Run forecasting pipeline to obtain final forecast and stance signal(t).
\STATE \hspace{0.8em}Determine execution day $t_{\text{exec}}$ (next trading day) and compute $P_{\text{buy}}(t)$ and $P_{\text{sell}}(t)$.
\STATE \hspace{0.8em}If signal(t)=BUY: open/maintain position; lock cash until $t+d_s$.
\STATE \hspace{0.8em}If signal(t)=SELL: close position; release cash after settlement.
\STATE \hspace{0.8em}Update portfolio value $V(\tau)$ including transaction costs and slippage.
\STATE \textbf{Output:} cumulative net return, max drawdown, and cost-adjusted directional accuracy.
\end{algorithmic}

\noindent This practical simulation complements MAE/RMSE/MAPE by evaluating whether forecast-driven decisions remain beneficial under Vietnamese settlement constraints and realistic execution frictions.

\section{RQ1: Architectural Effectiveness Across Horizons}
This section compares the coordination effectiveness of three multi-agent architectures, namely Hierarchical, Round Robin, and Ensemble Voting, across short-term, medium-term, and long-term forecasting horizons in Vietnam’s volatile stock market. The Hierarchical architecture employs centralized aggregation to enhance stability. Round Robin enables sequential refinement but may suffer from error propagation. Ensemble Voting combines parallel predictions to reduce individual model bias through collective decision-making.
\clearpage
\onecolumn

{
\setlength{\tabcolsep}{3pt}
\renewcommand{\arraystretch}{1.05}
\scriptsize
\begin{longtable}{llcc|cccc|cccc|cccc}
\caption{Comparison of Short-term, Medium-term, and Long-term Forecasts Across Architectures and Stocks (LLM: Google Gemini 2.0 Flash, LSTM Baseline)}
\label{tab:full_forecast} \\
\toprule
\multirow{2}{*}{\textbf{Stock}} &
\multirow{2}{*}{\textbf{Architecture}} &
\multicolumn{2}{c|}{\textbf{Start}} &
\multicolumn{4}{c|}{\textbf{Short-term}} &
\multicolumn{4}{c|}{\textbf{Medium-term}} &
\multicolumn{4}{c}{\textbf{Long-term}} \\
\cmidrule(lr){3-4}
\cmidrule(lr){5-8}
\cmidrule(lr){9-12}
\cmidrule(lr){13-16}
& &
\textbf{Date} & \textbf{Price} &
\textbf{End Date} & \textbf{Pred} & \textbf{Act} & \textbf{Gap} &
\textbf{End Date} & \textbf{Pred} & \textbf{Act} & \textbf{Gap} &
\textbf{End Date} & \textbf{Pred} & \textbf{Act} & \textbf{Gap} \\
\midrule
\endfirsthead

\multicolumn{16}{c}%
{{\tablename\ \thetable{} -- continued from previous page}} \\
\toprule
\multirow{2}{*}{\textbf{Stock}} &
\multirow{2}{*}{\textbf{Architecture}} &
\multicolumn{2}{c|}{\textbf{Start}} &
\multicolumn{4}{c|}{\textbf{Short-term}} &
\multicolumn{4}{c|}{\textbf{Medium-term}} &
\multicolumn{4}{c}{\textbf{Long-term}} \\
\cmidrule(lr){3-4}
\cmidrule(lr){5-8}
\cmidrule(lr){9-12}
\cmidrule(lr){13-16}
& &
\textbf{Date} & \textbf{Price} &
\textbf{End Date} & \textbf{Pred} & \textbf{Act} & \textbf{Gap} &
\textbf{End Date} & \textbf{Pred} & \textbf{Act} & \textbf{Gap} &
\textbf{End Date} & \textbf{Pred} & \textbf{Act} & \textbf{Gap} \\
\midrule
\endhead

\midrule \multicolumn{16}{r}{{Continued on next page}} \\
\endfoot

\endlastfoot


%%%%%%%%%%%%%%%%%%%%%%%%%%%%%%%%%%%%%%%%%%%%%%%%%%%%%%%%%%%%%%%%%%%%%%%%%%%%%%%%%%%%
% ===================== DXG ======================
%%%%%%%%%%%%%%%%%%%%%%%%%%%%%%%%%%%%%%%%%%%%%%%%%%%%%%%%%%%%%%%%%%%%%%%%%%%%%%%%%%%%
\multirow{4}{*}{\textbf{DXG}} 
& \textbf{Hierarchical}
& 31/08 & 22.8 & 03/09 & 22.25 & 24 & \textcolor{red}{\(\downarrow\)-1.75} & 14/09 & 23.38 & 24.05 & \textcolor{green}{\(\uparrow\)+0.67} & 30/10 & 21.83 & 21.2 & \textcolor{green}{\(\uparrow\)+0.63} \\

& 
& 16/10 & 22.45 & 19/10 & 22.10 & 22.6 & \textcolor{red}{\(\downarrow\)-0.50} & 30/10 & 23.34 & 21.2 & \textcolor{green}{\(\uparrow\)+2.14} & 15/12 & 24.46 & 16.3 & \textcolor{red}{\(\downarrow\)-8.16} \\

& 
& 21/10 & 20.0 & 24/10 & 20.51 & 20.9 & \textcolor{green}{\(\uparrow\)+0.39} & 04/11 & 20.52 & 20.15 & \textcolor{green}{\(\uparrow\)+0.37} & 20/12 & 20.81 & 17.8 & \textcolor{red}{\(\downarrow\)-3.01} \\

&
& 22/10 & 20.5 & 25/10 & 21.10 & 20.9 & \textcolor{green}{\(\uparrow\)+0.20} & 05/11 & 28.88 & 20.1 & \textcolor{red}{\(\downarrow\)-8.78} & 21/12 & 20.42 & 17.8 & \textcolor{green}{\(\uparrow\)+2.62} \\
\midrule


\multirow{4}{*}{\textbf{DXG}} 
& \textbf{Round Robin}
& 31/08 & 22.8 & 03/09 & 22.25 & 24 & \textcolor{red}{\(\downarrow\)-1.75} & 14/09 & 21.98 & 24.05 & \textcolor{red}{\(\downarrow\)-2.07} & 30/10 & 24.47 & 21.2 & \textcolor{red}{\(\downarrow\)-3.27} \\

&
& 16/10 & 22.45 & 19/10 & 22.15 & 22.6 & \textcolor{red}{\(\downarrow\)-0.45} & 30/10 & 22.28 & 21.2 & \textcolor{green}{\(\uparrow\)+1.08} & 15/12 & 25.63 & 16.3 & \textcolor{red}{\(\downarrow\)-9.33} \\

&
& 21/10 & 20.0 & 24/10 & 20.34 & 20.9 & \textcolor{green}{\(\uparrow\)+0.56} & 04/11 & 20.79 & 20.15 & \textcolor{green}{\(\uparrow\)+0.64} & 20/12 & 20.96 & 17.8 & \textcolor{red}{\(\downarrow\)-3.16} \\

&
& 22/10 & 20.5 & 25/10 & 20.48 & 20.9 & \textcolor{red}{\(\downarrow\)-0.42} & 05/11 & 23.68 & 20.1 & \textcolor{red}{\(\downarrow\)-3.58} & 21/12 & 20.96 & 17.8 & \textcolor{red}{\(\downarrow\)-3.16} \\
\midrule


\multirow{4}{*}{\textbf{DXG}} 
& \textbf{Ensemble}
& 31/08 & 22.8 & 03/09 & 22.11 & 24 & \textcolor{red}{\(\downarrow\)-1.89} & 14/09 & 22.98 & 24.05 & \textcolor{green}{\(\uparrow\)+1.07} & 30/10 & 21.47 & 21.2 & \textcolor{red}{\(\downarrow\)-0.27} \\

&
& 16/10 & 22.45 & 19/10 & 21.92 & 22.6 & \textcolor{red}{\(\downarrow\)-0.68} & 30/10 & 21.75 & 21.2 & \textcolor{green}{\(\uparrow\)+0.55} & 15/12 & 23.36 & 16.3 & \textcolor{red}{\(\downarrow\)-7.06} \\

&
& 21/10 & 20.0 & 24/10 & 20.07 & 20.9 & \textcolor{green}{\(\uparrow\)+0.83} & 04/11 & 20.52 & 20.15 & \textcolor{green}{\(\uparrow\)+0.37} & 20/12 & 20.81 & 17.8 & \textcolor{red}{\(\downarrow\)-3.01} \\

&
& 22/10 & 20.5 & 25/10 & 21.05 & 20.9 & \textcolor{green}{\(\uparrow\)+0.15} & 05/11 & 26.35 & 20.1 & \textcolor{red}{\(\downarrow\)-6.25} & 21/12 & 20.45 & 17.8 & \textcolor{green}{\(\uparrow\)+2.65} \\
\midrule



%%%%%%%%%%%%%%%%%%%%%%%%%%%%%%%%%%%%%%%%%%%%%%%%%%%%%%%%%%%%%%%%%%%%%%%%%%%%%%%%%%%%
% ===================== FPT ======================
%%%%%%%%%%%%%%%%%%%%%%%%%%%%%%%%%%%%%%%%%%%%%%%%%%%%%%%%%%%%%%%%%%%%%%%%%%%%%%%%%%%%
\multirow{4}{*}{\textbf{FPT}} 
& \textbf{Hierarchical}
& 01/09 & 101.6 & 04/09 & 103.60 & 105 & \textcolor{green}{\(\uparrow\)+1.40} & 15/09 & 106.07 & 101.9 & \textcolor{red}{\(\downarrow\)-4.17} & 31/10 & 105.57 & 103.9 & \textcolor{green}{\(\uparrow\)+1.67} \\

&
& 16/10 & 89.8 & 19/10 & 90.83 & 88.1 & \textcolor{red}{\(\downarrow\)-2.73} & 30/10 & 89.32 & 102.7 & \textcolor{red}{\(\downarrow\)-13.38} & 15/12 & 92.13 & 93.8 & \textcolor{green}{\(\uparrow\)+1.67} \\

&
& 21/10 & 93 & 24/10 & 95.37 & 97.7 & \textcolor{green}{\(\uparrow\)+2.33} & 04/11 & 95.80 & 103.3 & \textcolor{green}{\(\uparrow\)+7.50} & 20/12 & 97.56 & 93.9 &\textcolor{green}{\(\uparrow\)+3.66}\\

&
& 22/10 & 97 & 25/10 & 93.5 & 97.7 & \textcolor{red}{\(\downarrow\)-4.20} & 05/11 & 97.60 & 100.9 & \textcolor{green}{\(\uparrow\)+3.30} & 21/12 & 108.92 & 93.9 & \textcolor{red}{\(\downarrow\)-15.02} \\
\midrule


\multirow{4}{*}{\textbf{FPT}} 
& \textbf{Round Robin}
& 01/09 & 101.6 & 04/09 & 104.42 & 105 & \textcolor{green}{\(\uparrow\)+0.58} & 15/09 & 105.66 & 101.9 & \textcolor{red}{\(\downarrow\)-3.76} & 31/10 & 104.27 & 103.9 & \textcolor{green}{\(\uparrow\)+0.37} \\

&
& 16/10 & 89.8 & 19/10 & 90.46 & 88.1 & \textcolor{red}{\(\downarrow\)-2.36} & 30/10 & 91.28 & 102.7 & \textcolor{red}{\(\downarrow\)-11.42} & 15/12 & 90.70 & 93.8 & \textcolor{green}{\(\uparrow\)+3.1} \\

&
& 21/10 & 93 & 24/10 & 96.00 & 97.7 & \textcolor{green}{\(\uparrow\)+1.70} & 04/11 & 95.25 & 103.3 & \textcolor{green}{\(\uparrow\)+8.05} & 20/12 & 97.48 & 93.9 & \textcolor{green}{\(\uparrow\)+3.58} \\

&
& 22/10 & 97 & 25/10 & 93.68 & 97.7 & \textcolor{red}{\(\downarrow\)-4.02} & 05/11 & 93.60 & 100.9 & \textcolor{green}{\(\uparrow\)+7.30} & 21/12 & 108.92 & 93.9 & \textcolor{red}{\(\downarrow\)-15.02} \\
\midrule



\multirow{4}{*}{\textbf{FPT}} 
& \textbf{Ensemble}
& 01/09 & 101.6 & 04/09 & 103.60 & 105 & \textcolor{green}{\(\uparrow\)+1.40} & 15/09 & 103.6 & 101.9 & \textcolor{red}{\(\downarrow\)-1.70} & 31/10 & 102.77 & 103.9 & \textcolor{green}{\(\uparrow\)+1.13} \\

&
& 16/10 & 89.8 & 19/10 & 88.55 & 88.1 & \textcolor{green}{\(\uparrow\)-0.45} & 30/10 & 91.17 & 102.7 & \textcolor{red}{\(\downarrow\)-11.53} & 15/12 & 92.24 & 93.8 & \textcolor{green}{\(\uparrow\)+1.56} \\

&
& 21/10 & 93 & 24/10 & 95.14 & 97.7 & \textcolor{green}{\(\uparrow\)+2.56} & 04/11 & 96.26 & 103.3 & \textcolor{green}{\(\uparrow\)+7.04} & 20/12 & 97.22 & 93.9 & \textcolor{green}{\(\uparrow\)+3.32} \\

&
& 22/10 & 97 & 25/10 & 96.59 & 97.7 & \textcolor{green}{\(\uparrow\)+1.11} & 05/11 & 97.93 & 100.9 & \textcolor{green}{\(\uparrow\)+2.97} & 21/11 & 105.33 & 93.9 & \textcolor{red}{\(\downarrow\)-11.43} \\
\midrule




%%%%%%%%%%%%%%%%%%%%%%%%%%%%%%%%%%%%%%%%%%%%%%%%%%%%%%%%%%%%%%%%%%%%%%%%%%%%%%%%%%%%
% ===================== VCB ======================
%%%%%%%%%%%%%%%%%%%%%%%%%%%%%%%%%%%%%%%%%%%%%%%%%%%%%%%%%%%%%%%%%%%%%%%%%%%%%%%%%%%%
\multirow{4}{*}{\textbf{VCB}} 
& \textbf{Hierarchical}
& 30/08 & 68.6 & 02/09 & 68.66 & 68.6 & \textcolor{yellow}{0.06} & 13/09 & 68.52 & 65.8 & \textcolor{red}{\(\downarrow\)-2.72} & 29/10 & 67.6 & 60.7 & \textcolor{green}{\(\uparrow\)+6.9} \\

&
& 16/10 & 62.9 & 19/10 & 63.29 & 61.9 & \textcolor{red}{\(\downarrow\)-1.39} & 30/10 & 64.29 & 60.6 & \textcolor{red}{\(\downarrow\)-3.69} & 15/12 & 64.47 & 56.8 & \textcolor{red}{\(\downarrow\)-7.67} \\

&
& 21/10 & 59.3 & 24/10 & 61.68 & 59.5 & \textcolor{green}{\(\uparrow\)-2.18} & 04/11 & 62.37 & 60.1 & \textcolor{green}{\(\uparrow\)+2.27} & 20/12 & 58.5 & 57.5 & \textcolor{green}{\(\uparrow\)+1.0} \\

&
& 22/10 & 59.6 & 25/10 & 61.21 & 59.5 & \textcolor{red}{\(\downarrow\)-1.71} & 05/11 & 62.37 & 60.8 & \textcolor{green}{\(\uparrow\)+1.57} & 21/12 & 62.76 & 57.5 & \textcolor{red}{\(\downarrow\)-5.26} \\
\midrule


\multirow{4}{*}{\textbf{VCB}} 
& \textbf{Round Robin}
& 30/08 & 68.6 & 02/09 & 68.46 & 68.6 & \textcolor{yellow}{0.14} & 13/09 & 68.9 & 65.8 & \textcolor{red}{\(\downarrow\)-3.10} & 29/10 & 77.81 & 60.7 & \textcolor{red}{\(\downarrow\)-17.11} \\

&
& 16/10 & 62.9 & 19/10 & 62.51 & 61.9 & \textcolor{green}{\(\uparrow\)+0.61} & 30/10 & 64.03 & 60.6 & \textcolor{red}{\(\downarrow\)-3.43} & 15/12 & 64.57 & 56.8 & \textcolor{red}{\(\downarrow\)-7.77} \\

&
& 21/10 & 59.3 & 24/10 & 61.15 & 59.5 & \textcolor{green}{\(\uparrow\)+1.65} & 04/11 & 62.37 & 60.1 & \textcolor{green}{\(\uparrow\)+2.27} & 20/12 & 63.00 & 57.5 & \textcolor{red}{\(\downarrow\)-5.5} \\

&
& 22/10 & 59.6 & 25/10 & 61.46 & 59.5 & \textcolor{red}{\(\downarrow\)-1.96} & 05/11 & 61.96 & 60.8 & \textcolor{green}{\(\uparrow\)+1.16} & 21/12 & 62.45 & 57.5 & \textcolor{red}{\(\downarrow\)-4.95} \\
\midrule

\multirow{4}{*}{\textbf{VCB}} 
& \textbf{Ensemble}
& 30/08 & 68.6 & 02/09 & 67.64 & 68.6 & \textcolor{yellow}{0.96} & 13/09 & 68.85 & 65.8 & \textcolor{red}{\(\downarrow\)-3.05} & 29/10 & 67.9 & 60.7 & \textcolor{green}{\(\uparrow\)+7.2} \\

&
& 16/10 & 62.9 & 19/10 & 63.65 & 61.9 & \textcolor{red}{\(\downarrow\)-1.75} & 30/10 & 64.34 & 60.6 & \textcolor{red}{\(\downarrow\)-3.74} & 15/12 & 61.94 & 56.8 & \textcolor{green}{\(\uparrow\)+5.14} \\

&
& 21/10 & 59.3 & 24/10 & 62.58 & 59.5 & \textcolor{green}{\(\uparrow\)+3.08} & 04/11 & 62.36 & 60.1 & \textcolor{green}{\(\uparrow\)+2.26} & 20/12 & 63.01 & 57.5 & \textcolor{red}{\(\downarrow\)-5.51} \\

&
& 22/10 & 59.6 & 25/10 & 60.88 & 59.5 & \textcolor{red}{\(\downarrow\)-1.38} & 05/11 & 61.78 & 60.8 & \textcolor{green}{\(\uparrow\)+0.98} & 21/12 & 62.29 & 57.5 & \textcolor{red}{\(\downarrow\)-4.79} \\
\midrule



%%%%%%%%%%%%%%%%%%%%%%%%%%%%%%%%%%%%%%%%%%%%%%%%%%%%%%%%%%%%%%%%%%%%%%%%%%%%%%%%%%%%
% ===================== VCS ======================
%%%%%%%%%%%%%%%%%%%%%%%%%%%%%%%%%%%%%%%%%%%%%%%%%%%%%%%%%%%%%%%%%%%%%%%%%%%%%%%%%%%%
\multirow{4}{*}{\textbf{VCS}} 
& \textbf{Hierarchical}
& 30/08 & 48.8 & 02/09 & 47.97 & 48.8 & \textcolor{yellow}{0.83} & 13/09 & 50.20 & 50 & \textcolor{green}{\(\uparrow\)+0.20} & 29/10 & 50.51 & 47.6 & \textcolor{red}{\(\downarrow\)-2.91} \\

&
& 16/10 & 47.2 & 19/10 & 47.86 & 46.8 & \textcolor{red}{\(\downarrow\)-1.06} & 30/10 & 49.74 & 47.4 & \textcolor{green}{\(\uparrow\)+2.34} & 15/12 & 42.13 & 46.5 & \textcolor{green}{\(\uparrow\)+4.37} \\

&
& 21/10 & 45.5 & 24/10 & 47.12 & 46.2 & \textcolor{green}{\(\uparrow\)+0.92} & 04/11 & 47.06 & 47.1 & \textcolor{green}{\(\uparrow\)+0.04} & 20/12 & 47.66 & 43 & \textcolor{red}{\(\downarrow\)-4.66} \\

&
& 22/10 & 45.5 & 25/10 & 46.38 & 46.2 & \textcolor{green}{\(\uparrow\)+0.18} & 05/11 & 47.27 & 46.9 & \textcolor{green}{\(\uparrow\)+0.37} & 21/12 & 47.56 & 43 & \textcolor{red}{\(\downarrow\)-4.56} \\
\midrule


\multirow{4}{*}{\textbf{VCS}} 
& \textbf{Round Robin}
& 30/08 & 48.8 & 02/09 & 48.62 & 48.8 & \textcolor{yellow}{0.18} & 13/09 & 46.2 & 50 & \textcolor{red}{\(\downarrow\)-3.80} & 29/10 & 50.51 & 47.6 & \textcolor{red}{\(\downarrow\)-2.91} \\

&
& 16/10 & 47.2 & 19/10 & 47.55 & 46.8 & \textcolor{red}{\(\downarrow\)-0.75} & 30/10 & 46.94 & 47.4 & \textcolor{green}{\(\uparrow\)+0.46} & 15/12 & 45.5 & 46.5 & \textcolor{green}{\(\uparrow\)+1} \\

&
& 21/10 & 45.5 & 24/10 & 47.43 & 46.2 & \textcolor{green}{\(\uparrow\)+1.23} & 04/11 & 46.12 & 47.1 & \textcolor{green}{\(\uparrow\)+0.98} & 20/12 & 44.16 & 43 & \textcolor{green}{\(\uparrow\)+1.16} \\

&
& 22/10 & 45.5 & 25/10 & 46.92 & 46.2 & \textcolor{green}{\(\uparrow\)+0.72} & 05/11 & 46.6 & 46.9 & \textcolor{green}{\(\uparrow\)+0.30} & 21/12 & 44.15 & 43 & \textcolor{green}{\(\uparrow\)+1.15} \\
\midrule



\multirow{4}{*}{\textbf{VCS}} 
& \textbf{Ensemble}
& 30/08 & 48.8 & 02/09 & 48.66 & 48.8 & \textcolor{yellow}{0.14} & 13/09 & 49.20 & 50 & \textcolor{green}{\(\uparrow\)+0.80} & 29/10 & 49.65 & 47.6 & \textcolor{red}{\(\downarrow\)-2.05} \\

&
& 16/10 & 47.2 & 19/10 & 45.48 & 46.8 & \textcolor{green}{\(\uparrow\)+1.32} & 30/10 & 45.74 & 47.4 & \textcolor{green}{\(\uparrow\)+1.66} & 15/12 & 41.33 & 46.5 & \textcolor{green}{\(\uparrow\)-5.17} \\

&
& 21/10 & 45.5 & 24/10 & 46.95 & 46.2 & \textcolor{green}{\(\uparrow\)+0.75} & 04/11 & 46.78 & 47.1 & \textcolor{green}{\(\uparrow\)+0.32} & 20/12 & 47.43 & 43 & \textcolor{red}{\(\downarrow\)-4.43} \\

&
& 22/10 & 45.5 & 25/10 & 45.70 & 46.2 & \textcolor{green}{\(\uparrow\)+0.50} & 05/11 & 46.66 & 46.9 & \textcolor{green}{\(\uparrow\)+0.24} & 21/12 & 46.53 & 43 & \textcolor{red}{\(\downarrow\)-3.53}  \\


\multirow{4}{*}{\textbf{HPG}} 
& \textbf{Hierarchical}
& 01/09 & 27.5 & 04/09 & 28.29 & 29.85 & \textcolor{green}{\(\uparrow\)+1.56} & 15/09 & 28 & 30.35 & \textcolor{green}{\(\uparrow\)+2.35} & 31/10 & 28.1 & 26.7 & \textcolor{red}{\(\downarrow\)-1.2} \\

& 
& 16/10 & 28.25 & 19/10 & 28.60 & 28 & \textcolor{green}{\(\uparrow\)+0.6} & 30/10 & 29.01 & 26.9 & \textcolor{red}{\(\downarrow\)-2.11} & 15/12 & 25.11 & 26.25 & \textcolor{green}{\(\uparrow\)+1.14} \\

& 
& 21/10 & 26.8 & 24/10 & 28.02 & 26.4 & \textcolor{red}{\(\downarrow\)-1.62} & 04/11 & 27.2 & 26.75 & \textcolor{red}{\(\downarrow\)-0.45} & 20/12 & 28.03 & 26.7 & \textcolor{red}{\(\downarrow\)-1.33} \\

&
& 22/10 & 26.7 & 25/10 & 27.61 & 26.4 & \textcolor{red}{\(\downarrow\)-1.21} & 05/11 & 27 & 26.3 & \textcolor{red}{\(\downarrow\)-0.7} & 21/12 & 27.33 & 26.7 & \textcolor{yellow}{0.63} \\
\midrule


\multirow{4}{*}{\textbf{HPG}} 
& \textbf{Round Robin}
& 01/09 & 27.5 & 04/09 & 28.84 & 29.85 & \textcolor{green}{\(\uparrow\)+1.01} & 15/09 & 28.84 & 30.35 & \textcolor{green}{\(\uparrow\)+1.51} & 31/10 & 26.41 & 26.7 & \textcolor{green}{\(\uparrow\)+0.49} \\

& 
& 16/10 & 28.25 & 19/10 & 29.63 & 28 & \textcolor{red}{\(\downarrow\)-1.63} & 30/10 & 29.01 & 26.9 & \textcolor{red}{\(\downarrow\)-2.11} & 15/12 & 25.11 & 26.25 & \textcolor{green}{\(\uparrow\)+1.14} \\

& 
& 21/10 & 26.8 & 24/10 & 26.6 & 26.4 & \textcolor{green}{\(\uparrow\)+0.2} & 04/11 & 27.2 & 26.75 & \textcolor{red}{\(\downarrow\)-0.45} & 20/12 & 28.03 & 26.7 & \textcolor{red}{\(\downarrow\)-1.33} \\

&
& 22/10 & 26.7 & 25/10 & 28 & 26.4 & \textcolor{red}{\(\downarrow\)-1.6} & 05/11 & 27 & 26.3 & \textcolor{red}{\(\downarrow\)-0.7} & 21/12 & 27.33 & 26.7 & \textcolor{yellow}{0.63} \\
\midrule


\multirow{4}{*}{\textbf{HPG}} 
& \textbf{Ensemble}
& 01/09 & 27.5 & 04/09 & 28.29 & 29.85& \textcolor{green}{\(\uparrow\)+1.02} & 15/09 & 28 & 30.35 & \textcolor{green}{\(\uparrow\)+1.8} & 31/10 & 28.1 & 26.7 & \textcolor{red}{\(\downarrow\)-1.88} \\

& 
& 16/10 & 28.25 & 19/10 & 28.60 & 28 & \textcolor{green}{\(\uparrow\)+0.6} & 30/10 & 29.01 & 26.9 & \textcolor{red}{\(\downarrow\)-2.11} & 15/12 & 25.11 & 26.25 & \textcolor{green}{\(\uparrow\)+1.14} \\

& 
& 21/10 & 26.8 & 24/10 & 28.02 & 26.4 & \textcolor{red}{\(\downarrow\)-1.62} & 04/11 & 27.2 & 26.75 & \textcolor{red}{\(\downarrow\)-0.45} & 20/12 & 28.03 & 26.7 & \textcolor{red}{\(\downarrow\)-1.33} \\

&
& 22/10 & 26.7 & 25/10 & 27.61 & 26.4 & \textcolor{red}{\(\downarrow\)-1.21} & 05/11 & 27 & 26.3 & \textcolor{red}{\(\downarrow\)-0.7} & 21/12 & 27.33 & 26.7 & \textcolor{yellow}{0.63} \\
\midrule



%%%%%%%%%%%%%%%%%%%%%%%%%%%%%%%%%%%%%%%%%%%%%%%%%%%%%%%%%%%%%%%%%%%%%%%%%%%%%%%%%%%%
% ===================== TCB ======================
%%%%%%%%%%%%%%%%%%%%%%%%%%%%%%%%%%%%%%%%%%%%%%%%%%%%%%%%%%%%%%%%%%%%%%%%%%%%%%%%%%%%
\multirow{4}{*}{\textbf{TCB}} 
& \textbf{Ensemble}
& 01/09 & 39.6 & 04/09 & 35.28 & 39.6 & \textcolor{yellow}{4.32} & 15/09 & 35.10 & 39 & \textcolor{green}{\(\uparrow\)+3.9} & 31/10 & 35.42 & 35.1 & \textcolor{green}{\(\uparrow\)+0.32} \\

&
& 16/10 & 41.25 & 19/10 & 37.9 & 40.65 & \textcolor{green}{\(\uparrow\)+2,75} & 30/10 & 37.94 & 35.7 & \textcolor{green}{\(\uparrow\)+2.24} & 15/12 & 37.97 & 32 & \textcolor{green}{\(\uparrow\)+5.97} \\

&
& 21/10 & 37.65 & 24/10 & 34.78 & 36.1 & \textcolor{green}{\(\uparrow\)+1.32} & 04/11 & 34.67 & 35 & \textcolor{green}{\(\uparrow\)+0.33} & 20/12 & 34.76 & 33.6 &\textcolor{green}{\(\uparrow\)+1.16}\\

&
& 22/10 & 38 & 25/10 & 35.1 & 36.1 & \textcolor{green}{\(\uparrow\)+1.0} & 05/11 & 35.14 & 34.1 & \textcolor{green}{\(\uparrow\)+1.04} & 21/12 & 35.10 & 33.6 & \textcolor{green}{\(\uparrow\)+1.50} \\
\midrule


\multirow{4}{*}{\textbf{TCB}} 
& \textbf{Hierarchical}
& 01/09 & 39.6 & 04/09 & 40.02 & 39.6 & \textcolor{yellow}{0.42} & 15/09 & 39.82 & 39 & \textcolor{red}{\(\downarrow\)-0.82} & 31/10 & 40.03 & 35.1 & \textcolor{red}{\(\downarrow\)-4.93} \\

&
& 16/10 & 41.25 & 19/10 & 42 & 40.65 & \textcolor{red}{\(\downarrow\)-1.35} & 30/10 & 42.27 & 35.7 & \textcolor{red}{\(\downarrow\)-6.57} & 15/12 & 41.2 & 32 & \textcolor{green}{\(\uparrow\)+9.2} \\

&
& 21/10 & 37.65 & 24/10 & 37.5 & 36.1 & \textcolor{green}{\(\uparrow\)+1.4} & 04/11 & 38.82 & 35 & \textcolor{red}{\(\downarrow\)-3.82} & 20/12 & 38 & 33.6 &\textcolor{red}{\(\downarrow\)-4.4}\\

&
& 22/10 & 38 & 25/10 & 39.01 & 36.1 & \textcolor{red}{\(\downarrow\)-2.91} & 05/11 & 38.5 & 34.1 & \textcolor{red}{\(\downarrow\)-4.4} & 21/12 & 37.03 & 33.6 & \textcolor{green}{\(\uparrow\)+3.43} \\
\midrule



\multirow{4}{*}{\textbf{TCB}} 
& \textbf{RoundRobin}
& 01/09 & 39.6 & 04/09 & 37.81 & 39.6 & \textcolor{yellow}{1.79} & 15/09 & 38.25 & 39 & \textcolor{green}{\(\uparrow\)+0.75} & 31/10 & 36.91 & 35.1 & \textcolor{green}{\(\uparrow\)+1.81} \\

&
& 16/10 & 41.25 & 19/10 & 42.81 & 40.65 & \textcolor{red}{\(\downarrow\)-2.16} & 30/10 & 36.07 & 35.7 & \textcolor{red}{\(\downarrow\)-0.37} & 15/12 & 36.03 & 32 & \textcolor{green}{\(\uparrow\)+4.03} \\

&
& 21/10 & 37.65 & 24/10 & 36.73 & 36.1 & \textcolor{green}{\(\uparrow\)+0.63} & 04/11 & 36.08 & 35 & \textcolor{green}{\(\uparrow\)+1.08} & 20/12 & 33.59 & 33.6 &\textcolor{green}{\(\uparrow\)+0.01}\\

&
& 22/10 & 38 & 25/10 & 37.22 & 36.1 & \textcolor{green}{\(\uparrow\)+1.12} & 05/11 & 38.62 & 34.1 & \textcolor{red}{\(\downarrow\)-4.52} & 21/12 & 36.93 & 33.6 & \textcolor{green}{\(\uparrow\)+3.33} \\
\midrule




%%%%%%%%%%%%%%%%%%%%%%%%%%%%%%%%%%%%%%%%%%%%%%%%%%%%%%%%%%%%%%%%%%%%%%%%%%%%%%%%%%%%
% ===================== VHM ======================
%%%%%%%%%%%%%%%%%%%%%%%%%%%%%%%%%%%%%%%%%%%%%%%%%%%%%%%%%%%%%%%%%%%%%%%%%%%%%%%%%%%%
\multirow{4}{*}{\textbf{VHM}} 
& \textbf{Hierarchical}
& 01/09 & 104.5 & 04/09 & 105.5 & 102.3 & \textcolor{red}{\(\downarrow\)-3.2} & 15/09 & 107.54 & 104.5 & \textcolor{yellow}{3.04} & 31/10 & 106.01 & 99.2 & \textcolor{red}{\(\downarrow\)-6.81} \\

&
& 16/10 & 122 & 19/10 & 124.69 & 116 & \textcolor{red}{\(\downarrow\)-8.69} & 30/10 & 115.50 & 104 & \textcolor{green}{\(\uparrow\)+11.5} & 15/12 & 105.04 & 93.9 & \textcolor{green}{\(\uparrow\)+11.14} \\

&
& 21/10 & 110.9 & 24/10 & 114.33 & 114.5 & \textcolor{green}{\(\uparrow\)+0.17} & 04/11 & 113.84 & 100.2 & \textcolor{red}{\(\downarrow\)-13.64} & 20/12 & 102.01 & 101.5 & \textcolor{green}{\(\uparrow\)+0.51} \\

&
& 22/10 & 113 & 25/10 & 114.64 & 114.5 & \textcolor{green}{\(\uparrow\)+0.14} & 05/11 & 105.49 & 99.6 & \textcolor{green}{\(\uparrow\)+5.89} & 21/12 & 106.45 & 101.5 & \textcolor{green}{\(\uparrow\)+4.95} \\
\midrule


\multirow{4}{*}{\textbf{VHM}} 
& \textbf{Round Robin}
& 01/09 & 104.5 & 04/09 & 105.13 & 102.3 & \textcolor{red}{\(\downarrow\)-2.83} & 15/09 & 107.47 & 104.5 & \textcolor{yellow}{2.97} & 31/10 & 107.83 & 99.2 & \textcolor{red}{\(\downarrow\)-8.63} \\

&
& 16/10 & 122 & 19/10 & 122.73 & 116 & \textcolor{red}{\(\downarrow\)-6.73} & 30/10 & 105.23 & 104 & \textcolor{green}{\(\uparrow\)+1.23} & 15/12 & 100.56 & 93.9 & \textcolor{green}{\(\uparrow\)+6.66} \\

&
& 21/10 & 110.9 & 24/10 & 111.57 & 114.5 & \textcolor{green}{\(\uparrow\)+2.93} & 04/11 & 104.05 & 100.2 & \textcolor{green}{\(\uparrow\)+3.85} & 20/12 & 105.05 & 101.5 & \textcolor{green}{\(\uparrow\)+3.55} \\

&
& 22/10 & 113 & 25/10 & 113.68 & 114.5 & \textcolor{green}{\(\uparrow\)+0.82} & 05/11 & 106.21 & 99.6 & \textcolor{green}{\(\uparrow\)+6.61} & 21/12 & 105.71 & 101.5 & \textcolor{green}{\(\uparrow\)+4.21} \\
\midrule


\multirow{4}{*}{\textbf{VHM}} 
& \textbf{Ensemble}
& 01/09 & 104.5 & 04/09 & 105.18 & 102.3 & \textcolor{red}{\(\downarrow\)-2.88} & 15/09 & 105.31 & 104.5 & \textcolor{yellow}{0.81} & 31/10 & 95.79 & 99.2 & \textcolor{green}{\(\uparrow\)+3.41} \\

&
& 16/10 & 122 & 19/10 & 111.7 & 116 & \textcolor{green}{\(\uparrow\)+4.3} & 30/10 & 112 & 104 & \textcolor{green}{\(\uparrow\)+8} & 15/12 & 102.3 & 93.9 & \textcolor{green}{\(\uparrow\)+8.4} \\

&
& 21/10 & 110.9 & 24/10 & 111.66 & 114.5 & \textcolor{green}{\(\uparrow\)+2.84} & 04/11 & 101.29 & 100.2 & \textcolor{green}{\(\uparrow\)+1.09} & 20/12 & 99.62 & 101.5 & \textcolor{green}{\(\uparrow\)+1.88} \\

&
& 22/10 & 113 & 25/10 & 117.23 & 114.5 & \textcolor{green}{\(\uparrow\)+2.73} & 05/11 & 95.86 & 99.6 & \textcolor{green}{\(\uparrow\)+3.74} & 21/12 & 94.42 & 101.5 & \textcolor{green}{\(\uparrow\)+7.08} \\
\midrule




%%%%%%%%%%%%%%%%%%%%%%%%%%%%%%%%%%%%%%%%%%%%%%%%%%%%%%%%%%%%%%%%%%%%%%%%%%%%%%%%%%%%
% ===================== MWG ======================
%%%%%%%%%%%%%%%%%%%%%%%%%%%%%%%%%%%%%%%%%%%%%%%%%%%%%%%%%%%%%%%%%%%%%%%%%%%%%%%%%%%%
\multirow{4}{*}{\textbf{MWG}} 
& \textbf{Hierarchical}
& 01/09 & 78 & 04/09 & 73.12 & 77.5 & \textcolor{green}{\(\uparrow\)+4.38} & 15/09 & 73.42 & 79.5 & \textcolor{red}{\(\downarrow\)-6.08} & 31/10 & 78.64 & 82.6 & \textcolor{green}{\(\uparrow\)+3.96} \\

&
& 16/10 & 84.5 & 19/10 & 87.88 & 84.5 & \textcolor{yellow}{3.38} & 30/10 & 88.18 & 83.9 & \textcolor{red}{\(\downarrow\)-4.28} & 15/12 & 81.35 & 77.7 & \textcolor{green}{\(\uparrow\)+3.65} \\

&
& 21/10 & 83 & 24/10 & 86.4 & 85.7 & \textcolor{green}{\(\uparrow\)+0.7} & 04/11 & 86.08 & 81.8 & \textcolor{red}{\(\downarrow\)-4.28} & 20/12 & 81.23 & 82.9 & \textcolor{green}{\(\uparrow\)+1.67} \\

&
& 22/10 & 84.5 & 25/10 & 87.6 & 85.7 & \textcolor{green}{\(\uparrow\)+1.9} & 05/11 & 83.14 & 80.2 & \textcolor{green}{\(\uparrow\)+2.94} & 21/12 & 87.42 & 82.9 & \textcolor{red}{\(\downarrow\)-4.52} \\
\midrule


\multirow{4}{*}{\textbf{MWG}} 
& \textbf{Round Robin}
& 01/09 & 78 & 04/09 & 71.61 & 77.5 & \textcolor{green}{\(\uparrow\)+5.89} & 15/09 & 71.86 & 79.5 & \textcolor{red}{\(\downarrow\)-7.64} & 31/10 & 78.62 & 82.6 & \textcolor{green}{\(\uparrow\)+3.98} \\

&
& 16/10 & 84.5 & 19/10 & 87.75 & 84.5 & \textcolor{yellow}{3.25} & 30/10 & 77.72 & 83.9 & \textcolor{green}{\(\uparrow\)+6.18} & 15/12 & 77.81 & 77.7 & \textcolor{green}{\(\uparrow\)+0.11} \\

&
& 21/10 & 83 & 24/10 & 85.63 & 85.7 & \textcolor{green}{\(\uparrow\)+0.07} & 04/11 & 80.45 & 81.8 & \textcolor{green}{\(\uparrow\)+1.35} & 20/12 & 82.69 & 82.9 & \textcolor{green}{\(\uparrow\)+0.21} \\

&
& 22/10 & 84.5 & 25/10 & 87.2 & 85.7 & \textcolor{green}{\(\uparrow\)+1.5} & 05/11 & 79.27 & 80.2 & \textcolor{green}{\(\uparrow\)+0.93} & 21/12 & 78.35 & 82.9 & \textcolor{green}{\(\uparrow\)+4.55} \\
\midrule




\multirow{4}{*}{\textbf{MWG}} 
& \textbf{Ensemble}
& 01/09 & 78 & 04/09 & 77.31 & 77.5 & \textcolor{green}{\(\uparrow\)+0.19} & 15/09 & 75.71 & 79.5 & \textcolor{red}{\(\downarrow\)-3.79} & 31/10 & 76.21 & 82.6 & \textcolor{red}{\(\downarrow\)-6.39} \\

&
& 16/10 & 84.5 & 19/10 & 77.04 & 84.5 & \textcolor{red}{\(\downarrow\)-7.46} & 30/10 & 82.79 & 83.9 & \textcolor{green}{\(\uparrow\)+1.11} & 15/12 & 77.09 & 77.7 & \textcolor{green}{\(\uparrow\)+0.61} \\

&
& 21/10 & 83 & 24/10 & 79.49 & 85.7 & \textcolor{red}{\(\downarrow\)-6.21} & 04/11 & 80.18 & 81.8 & \textcolor{green}{\(\uparrow\)+1.62} & 20/12 & 81.82 & 82.9 & \textcolor{green}{\(\uparrow\)+1.08} \\

&
& 22/10 & 84.5 & 25/10 & 87.54 & 85.7 & \textcolor{green}{\(\uparrow\)+1.84} & 05/11 & 81.39 & 80.2 & \textcolor{green}{\(\uparrow\)+1.19} & 21/12 & 83.53 & 82.9 & \textcolor{green}{\(\uparrow\)+0.63} \\






\multirow{4}{*}{\textbf{CMG}} 
& \textbf{Hierarchical}
& 01/09 & 40.3 & 04/09 & 39.13 & 41.85 & \textcolor{red}{\(\downarrow\)-2.72} & 15/09 & 33.28 & 40.8 & \textcolor{red}{\(\downarrow\)-7.52} & 31/10 & 33.24 & 41.9 & \textcolor{red}{\(\downarrow\)-8.66} \\

& 
& 16/10 & 38 & 19/10 & 37.52 & 37.3 & \textcolor{green}{\(\uparrow\)+0.22} & 30/10 & 38.14 & 41.3 & \textcolor{green}{\(\uparrow\)+3.16} & 15/12 & 33.63 & 35.1 & \textcolor{green}{\(\uparrow\)+1.58} \\

& 
& 21/10 & 38.4 & 24/10 & 38.12 & 38.75 & \textcolor{red}{\(\downarrow\)-0.63} & 04/11 & 35.82 & 39.6 & \textcolor{red}{\(\downarrow\)-3.78} & 20/12 & 31.82 & 35.2 & \textcolor{green}{\(\uparrow\)+3.38} \\

&
& 22/10 & 38.5 & 25/10 & 35.33 & 38.75 & \textcolor{red}{\(\downarrow\)-3.42} & 05/11 & 37.44 & 39.3 & \textcolor{red}{\(\downarrow\)-1.86} & 21/12 & 31.21 & 35.2 & \textcolor{green}{\(\uparrow\)+3.99} \\
\midrule


\multirow{4}{*}{\textbf{CMG}} 
& \textbf{Round Robin}
& 01/09 & 40.3 & 04/09 & 37.27 & 41.85 & \textcolor{red}{\(\downarrow\)-4.58} & 15/09 & 36.84 & 40.8 & \textcolor{red}{\(\downarrow\)-3.96} & 31/10 & 36.88 & 41.9 & \textcolor{red}{\(\downarrow\)-5.02} \\

& 
& 16/10 & 38 & 19/10 & 39.28 & 37.3 & \textcolor{red}{\(\downarrow\)-1.98} & 30/10 & 36.58 & 41.3 & \textcolor{red}{\(\downarrow\)-4.72} & 15/12 & 32.52 & 35.1 & \textcolor{green}{\(\uparrow\)+2.58} \\

& 
& 21/10 & 38.4 & 24/10 & 36.22 & 38.75 & \textcolor{red}{\(\downarrow\)-2.53} & 04/11 & 36.72 & 39.6 & \textcolor{red}{\(\downarrow\)-2.88} & 20/12 & 31.29 & 35.2 & \textcolor{red}{\(\downarrow\)-3.91} \\

&
& 22/10 & 38.5 & 25/10 & 35.82 & 38.75 & \textcolor{red}{\(\downarrow\)-2.93} & 05/11 & 38.83 & 39.3 & \textcolor{green}{\(\uparrow\)+0.47} & 21/12 & 33.11 & 35.2 & \textcolor{green}{\(\uparrow\)+2.09} \\
\midrule


\multirow{4}{*}{\textbf{CMG}} 
& \textbf{Ensemble}
& 01/09 & 40.3 & 04/09 & 42.15 & 41.85 & \textcolor{green}{\(\uparrow\)+0.3} & 15/09 & 42.21 & 40.8 & \textcolor{green}{\(\uparrow\)+1.41} & 31/10 & 42.88 & 41.9 & \textcolor{green}{\(\uparrow\)+0.98} \\

& 
& 16/10 & 38 & 19/10 & 39.85 & 37.3 & \textcolor{red}{\(\downarrow\)-2.55} & 30/10 & 39.25 & 41.3 & \textcolor{green}{\(\uparrow\)+2.05} & 15/12 & 35.52 & 35.1 & \textcolor{green}{\(\uparrow\)+0.42} \\

& 
& 21/10 & 38.4 & 24/10 & 38.66 & 38.75 & \textcolor{green}{\(\uparrow\)+0.09} & 04/11 & 38.87 & 39.6 & \textcolor{green}{\(\uparrow\)+0.73} & 20/12 & 31.21 & 35.2 & \textcolor{green}{\(\uparrow\)+3.99} \\

&
& 22/10 & 38.5 & 25/10 & 39.62 & 38.75 & \textcolor{green}{\(\uparrow\)+0.87} & 05/11 & 41.26 & 39.3 & \textcolor{green}{\(\uparrow\)+1.96} & 21/12 & 39.41 & 35.2 & \textcolor{red}{\(\downarrow\)-4.21} \\
\midrule



%%%%%%%%%%%%%%%%%%%%%%%%%%%%%%%%%%%%%%%%%%%%%%%%%%%%%%%%%%%%%%%%%%%%%%%%%%%%%%%%%%%%
% ===================== KDH======================
%%%%%%%%%%%%%%%%%%%%%%%%%%%%%%%%%%%%%%%%%%%%%%%%%%%%%%%%%%%%%%%%%%%%%%%%%%%%%%%%%%%%
\multirow{4}{*}{\textbf{KDH}} 
& \textbf{Hierarchical}
& 01/09 & 36.5 & 04/09 & 33.18 & 36.75 & \textcolor{red}{\(\downarrow\)-3.57} & 15/09 & 33.38 & 34.85 & \textcolor{green}{\(\uparrow\)+1.47} & 31/10 & 33.47 & 35.85 & \textcolor{green}{\(\uparrow\)+2.38} \\

&
& 16/10 & 34.2 & 19/10 & 31.31 & 33.9 & \textcolor{green}{\(\uparrow\)+2.59} & 30/10 & 31.44 & 35.85 & \textcolor{red}{\(\downarrow\)-4.41} & 15/12 & 31.44 & 29.6 & \textcolor{green}{\(\uparrow\)+1.84} \\

&
& 21/10 & 31.75 & 24/10 & 28.48 & 33.8 & \textcolor{red}{\(\downarrow\)-9.32} & 04/11 & 29.07 & 35.6 & \textcolor{red}{\(\downarrow\)-6.53} & 20/12 & 28.95 & 32.65 &\textcolor{red}{\(\downarrow\)-3.7}\\

&
& 22/10 & 32.2 & 25/10 & 27.76 & 33.8 & \textcolor{red}{\(\downarrow\)-6.04} & 05/11 & 25.76 & 34.95 & \textcolor{red}{\(\downarrow\)-9.19} & 21/12 & 26.64 & 32.65 & \textcolor{red}{\(\downarrow\)-6.01} \\
\midrule


\multirow{4}{*}{\textbf{KDH}} 
& \textbf{Round Robin}
& 01/09 & 36.5 & 04/09 & 32.25 & 36.75 & \textcolor{green}{\(\uparrow\)+1.40} & 15/09 & 34.61 & 34.85 & \textcolor{red}{\(\downarrow\)-4.17} & 31/10 & 35.21 & 35.85 & \textcolor{green}{\(\uparrow\)+0.64} \\

&
& 16/10 & 34.2 & 19/10 & 34.89 & 33.9 & \textcolor{red}{\(\downarrow\)-0.99} & 30/10 & 34.61 & 35.85 & \textcolor{green}{\(\uparrow\)+1.24} & 15/12 & 29.52 & 29.6 & \textcolor{green}{\(\uparrow\)+0.08} \\

&
& 21/10 & 31.75 & 24/10 & 33.45 & 33.8 & \textcolor{green}{\(\uparrow\)+0.35} & 04/11 & 32.14 & 35.6 & \textcolor{green}{\(\uparrow\)+3.46} & 20/12 & 29.56 & 32.65 &\textcolor{red}{\(\downarrow\)-3.09}\\

&
& 22/10 & 32.2 & 25/10 & 33.24 & 33.8 & \textcolor{green}{\(\uparrow\)+0.56} & 05/11 & 33.45 & 34.95 & \textcolor{green}{\(\uparrow\)+1.5} & 21/12 & 30.54 & 32.65 & \textcolor{red}{\(\downarrow\)-2.11} \\
\midrule



\multirow{4}{*}{\textbf{KDH}} 
& \textbf{Ensemble}
& 01/09 & 36.5 & 04/09 & 39.25 & 36.75 & \textcolor{green}{\(\uparrow\)+2.5} & 15/09 & 35.71 & 34.85 & \textcolor{green}{\(\uparrow\)+0.86} & 31/10 & 34.97 & 35.85 & \textcolor{green}{\(\uparrow\)+0.88} \\

&
& 16/10 & 34.2 & 19/10 & 32.39 & 33.9 & \textcolor{red}{\(\downarrow\)-1.51} & 30/10 & 37.51 & 35.85 & \textcolor{green}{\(\uparrow\)+1.66} & 15/12 & 28.92 & 29.6 & \textcolor{green}{\(\uparrow\)+0.68} \\

&
& 21/10 & 31.75 & 24/10 & 35.55 & 33.8 & \textcolor{green}{\(\uparrow\)+1.75} & 04/11 & 37.04 & 35.6 & \textcolor{green}{\(\uparrow\)+1.44} & 20/12 & 34.16 & 32.65 &\textcolor{green}{\(\uparrow\)+1.51}\\

&
& 22/10 & 32.2 & 25/10 & 32.59 & 33.8 & \textcolor{green}{\(\uparrow\)+1.21} & 05/11 & 33.40 & 34.95 & \textcolor{green}{\(\uparrow\)+1.55} & 21/12 & 27.99 & 32.65 & \textcolor{red}{\(\downarrow\)-4.66} \\
\midrule




%%%%%%%%%%%%%%%%%%%%%%%%%%%%%%%%%%%%%%%%%%%%%%%%%%%%%%%%%%%%%%%%%%%%%%%%%%%%%%%%%%%%
% ===================== DGC ======================
%%%%%%%%%%%%%%%%%%%%%%%%%%%%%%%%%%%%%%%%%%%%%%%%%%%%%%%%%%%%%%%%%%%%%%%%%%%%%%%%%%%%
\multirow{4}{*}{\textbf{DGC}} 
& \textbf{Hierarchical}
& 01/09 & 98.2 & 04/09 & 96.01 & 99.3 & \textcolor{red}{\(\downarrow\)-3.29} & 15/09 & 94.5 & 99.5 & \textcolor{red}{\(\downarrow\)-5} & 31/10 & 95 & 96 & \textcolor{green}{\(\uparrow\)+1} \\

&
& 16/10 & 95 & 19/10 & 93.52 & 93.2 & \textcolor{green}{\(\uparrow\)+0.32} & 30/10 & 92.82 & 93.6 & \textcolor{green}{\(\uparrow\)+0.78} & 15/12 & 92.01 & 93 & \textcolor{green}{\(\uparrow\)+0.99} \\

&
& 21/10 & 89.2 & 24/10 & 87.01 & 92.4 & \textcolor{red}{\(\downarrow\)-5.39} & 04/11 & 86 & 94.6 & \textcolor{red}{\(\downarrow\)-8.6} & 20/12 & 76 & 70.2 & \textcolor{green}{\(\uparrow\)+5.8} \\

&
& 22/10 & 91 & 25/10 & 90.3 & 92.4 & \textcolor{red}{\(\downarrow\)-2.1} & 05/11 & 89.02 & 95 & \textcolor{red}{\(\downarrow\)-5.98} & 21/12 & 77.8 & 70.2 & \textcolor{green}{\(\uparrow\)+7.6} \\
\midrule


\multirow{4}{*}{\textbf{DGC}} 
& \textbf{Round Robin}
& 01/09 & 98.2 & 04/09 & 95.27 & 99.3 & \textcolor{red}{\(\downarrow\)-4.03} & 15/09 & 99.77 & 99.5 & \textcolor{green}{\(\uparrow\)+0.27} & 31/10 & 95.82 & 96 & \textcolor{green}{\(\uparrow\)+0.18} \\

&
& 16/10 & 95 & 19/10 & 92.82 & 93.2 & \textcolor{red}{\(\downarrow\)-0.38} & 30/10 & 95.12 & 93.6 & \textcolor{red}{\(\downarrow\)-1.52} & 15/12 & 94.36 & 93 & \textcolor{red}{\(\downarrow\)-1.36} \\

&
& 21/10 & 89.2 & 24/10 & 98.37 & 92.4 & \textcolor{green}{\(\uparrow\)+5.97} & 04/11 & 91.06 & 94.6 & \textcolor{green}{\(\uparrow\)+3.54} & 20/12 & 78.52 & 70.2 & \textcolor{green}{\(\uparrow\)+8.32} \\

&
& 22/10 & 91 & 25/10 & 97.22 & 92.4 & \textcolor{green}{\(\uparrow\)+4.82} & 05/11 & 96.4 & 95 & \textcolor{green}{\(\uparrow\)+1.4} & 21/12 & 77.92 & 70.2 & \textcolor{green}{\(\downarrow\)-7.72} \\
\midrule


\multirow{4}{*}{\textbf{DGC}} 
& \textbf{Ensemble}
& 01/09 & 98.2 & 04/09 & 94.36 & 99.3 & \textcolor{red}{\(\downarrow\)-4.94} & 15/09 & 100.97 & 99.5 & \textcolor{green}{\(\uparrow\)+1.47} & 31/10 & 95.31 & 96 & \textcolor{green}{\(\uparrow\)+0.69} \\

&
& 16/10 & 95 & 19/10 & 93.35 & 93.2 & \textcolor{green}{\(\uparrow\)+0.15} & 30/10 & 93.84 & 93.6 & \textcolor{green}{\(\uparrow\)+0.24} & 15/12 & 93.6 & 93 & \textcolor{green}{\(\uparrow\)+0.6} \\

&
& 21/10 & 89.2 & 24/10 & 94.82 & 92.4 & \textcolor{green}{\(\uparrow\)+2.42} & 04/11 & 95.11 & 94.6 & \textcolor{green}{\(\uparrow\)+0.51} & 20/12 & 77.25 & 70.2 & \textcolor{green}{\(\uparrow\)+7.05} \\

&
& 22/10 & 91 & 25/10 & 93.78 & 92.4 & \textcolor{green}{\(\uparrow\)+1.38} & 05/11 & 95.13 & 95 & \textcolor{green}{\(\uparrow\)+0.13} & 21/12 & 72.95 & 70.2 & \textcolor{green}{\(\uparrow\)+2.75} \\
\midrule



%%%%%%%%%%%%%%%%%%%%%%%%%%%%%%%%%%%%%%%%%%%%%%%%%%%%%%%%%%%%%%%%%%%%%%%%%%%%%%%%%%%%
% ===================== MBB ======================
%%%%%%%%%%%%%%%%%%%%%%%%%%%%%%%%%%%%%%%%%%%%%%%%%%%%%%%%%%%%%%%%%%%%%%%%%%%%%%%%%%%%
\multirow{4}{*}{\textbf{MBB}} 
& \textbf{Hierarchical}
& 01/09 & 27.75 & 04/09 & 26.77 & 28.25 & \textcolor{red}{\(\downarrow\)-1.48} & 15/09 & 29.54 & 26.85 & \textcolor{red}{\(\downarrow\)-2.69} & 31/10 & 22.21 & 23.6 & \textcolor{green}{\(\uparrow\)+1.39} \\

&
& 16/10 & 27.2 & 19/10 & 26.81 & 27.1 & \textcolor{green}{\(\uparrow\)+0.29} & 30/10 & 21.58 & 23.95 & \textcolor{green}{\(\uparrow\)+2.37} & 15/12 & 26.95 & 23.75 & \textcolor{green}{\(\uparrow\)+3.2} \\

&
& 21/10 & 25.3 & 24/10 & 24.78 & 24.4 & \textcolor{green}{\(\uparrow\)+0.38} & 04/11 & 25.18 & 24 & \textcolor{green}{\(\uparrow\)+1.18} & 20/12 & 24.96 & 24.7 & \textcolor{green}{\(\uparrow\)+0.26} \\

&
& 22/10 & 25.45 & 25/10 & 23.54 & 24.4 & \textcolor{green}{\(\uparrow\)+0.86} & 05/11 & 24.58 & 23.9 & \textcolor{green}{\(\uparrow\)+0.68} & 21/12 & 26.44 & 24.7 & \textcolor{red}{\(\downarrow\)-1.74} \\
\midrule


\multirow{4}{*}{\textbf{MBB}} 
& \textbf{Round Robin}
& 01/09 & 27.75 & 04/09 & 28.2 & 28.25 & \textcolor{green}{\(\uparrow\)+0.05} & 15/09 & 25.51 & 26.85 & \textcolor{green}{\(\uparrow\)+1.34} & 31/10 & 25.93 & 23.6 & \textcolor{green}{\(\uparrow\)+2.33} \\

&
& 16/10 & 27.2 & 19/10 & 27.7 & 27.1 & \textcolor{green}{\(\uparrow\)+0.6} & 30/10 & 26.78 & 23.95 & \textcolor{green}{\(\uparrow\)+2.83} & 15/12 & 22.67 & 23.75 & \textcolor{green}{\(\uparrow\)+1.08} \\

&
& 21/10 & 25.3 & 24/10 & 24.7 & 24.4 & \textcolor{green}{\(\uparrow\)+0.3} & 04/11 & 24.41 & 24 & \textcolor{green}{\(\uparrow\)+0.41} & 20/12 & 24.07 & 24.7 & \textcolor{green}{\(\uparrow\)+0.63} \\

&
& 22/10 & 25.45 & 25/10 & 23.13 & 24.4 & \textcolor{green}{\(\uparrow\)+1.27} & 05/11 & 22.99 & 23.9 & \textcolor{green}{\(\uparrow\)+0.91} & 21/12 & 22.94 & 24.7 & \textcolor{green}{\(\uparrow\)+1.76} \\
\midrule



\multirow{4}{*}{\textbf{MBB}} 
& \textbf{Ensemble}
& 01/09 & 27.75 & 04/09 & 29.65 & 28.25 & \textcolor{green}{\(\uparrow\)+1.4} & 15/09 & 25.61 & 26.85 & \textcolor{green}{\(\uparrow\)+1.24} & 31/10 & 24.48 & 23.6 & \textcolor{green}{\(\uparrow\)+0.88} \\

&
& 16/10 & 27.2 & 19/10 & 26.97 & 27.1 & \textcolor{green}{\(\uparrow\)+0.13} & 30/10 & 23.1 & 23.95 & \textcolor{green}{\(\uparrow\)+0.85} & 15/12 & 23.02 & 23.75 & \textcolor{green}{\(\uparrow\)+0.73} \\

&
& 21/10 & 25.3 & 24/10 & 24.97 & 24.4 & \textcolor{green}{\(\uparrow\)+0.57} & 04/11 & 24.35 & 24 & \textcolor{green}{\(\uparrow\)+0.35} & 20/12 & 24.73 & 24.7 & \textcolor{green}{\(\uparrow\)+0.03} \\

&
& 22/10 & 25.45 & 25/10 & 24.94 & 24.4 & \textcolor{green}{\(\uparrow\)+0.54} & 05/11 & 23.55 & 23.9 & \textcolor{green}{\(\uparrow\)+0.35} & 21/12 & 24.34 & 24.7 & \textcolor{green}{\(\uparrow\)+0.36} \\



\bottomrule
\end{longtable}
}







\clearpage
\onecolumn

{
\setlength{\tabcolsep}{3pt}
\renewcommand{\arraystretch}{1.05}
\scriptsize
\begin{longtable}{llcc|cccc|cccc|cccc}
\caption{Comparison of Short-term, Medium-term, and Long-term Forecasts Across Architectures and Stocks (LLM: Google Gemini 2.0 Flash, Transformer Baseline)}
\label{tab:full_forecast_2} \\
\toprule
\multirow{2}{*}{\textbf{Stock}} &
\multirow{2}{*}{\textbf{Architecture}} &
\multicolumn{2}{c|}{\textbf{Start}} &
\multicolumn{4}{c|}{\textbf{Short-term}} &
\multicolumn{4}{c|}{\textbf{Medium-term}} &
\multicolumn{4}{c}{\textbf{Long-term}} \\
\cmidrule(lr){3-4}
\cmidrule(lr){5-8}
\cmidrule(lr){9-12}
\cmidrule(lr){13-16}
& &
\textbf{Date} & \textbf{Price} &
\textbf{End Date} & \textbf{Pred} & \textbf{Act} & \textbf{Gap} &
\textbf{End Date} & \textbf{Pred} & \textbf{Act} & \textbf{Gap} &
\textbf{End Date} & \textbf{Pred} & \textbf{Act} & \textbf{Gap} \\
\midrule
\endfirsthead

\multicolumn{16}{c}%
{{\tablename\ \thetable{} -- continued from previous page}} \\
\toprule
\multirow{2}{*}{\textbf{Stock}} &
\multirow{2}{*}{\textbf{Architecture}} &
\multicolumn{2}{c|}{\textbf{Start}} &
\multicolumn{4}{c|}{\textbf{Short-term}} &
\multicolumn{4}{c|}{\textbf{Medium-term}} &
\multicolumn{4}{c}{\textbf{Long-term}} \\
\cmidrule(lr){3-4}
\cmidrule(lr){5-8}
\cmidrule(lr){9-12}
\cmidrule(lr){13-16}
& &
\textbf{Date} & \textbf{Price} &
\textbf{End Date} & \textbf{Pred} & \textbf{Act} & \textbf{Gap} &
\textbf{End Date} & \textbf{Pred} & \textbf{Act} & \textbf{Gap} &
\textbf{End Date} & \textbf{Pred} & \textbf{Act} & \textbf{Gap} \\
\midrule
\endhead

\midrule \multicolumn{16}{r}{{Continued on next page}} \\
\endfoot

\endlastfoot


%%%%%%%%%%%%%%%%%%%%%%%%%%%%%%%%%%%%%%%%%%%%%%%%%%%%%%%%%%%%%%%%%%%%%%%%%%%%%%%%%%%%
% ===================== DXG ======================
%%%%%%%%%%%%%%%%%%%%%%%%%%%%%%%%%%%%%%%%%%%%%%%%%%%%%%%%%%%%%%%%%%%%%%%%%%%%%%%%%%%%
\multirow{4}{*}{\textbf{DXG}} 
& \textbf{Hierarchical}
& 31/08 & 22.8 & 03/09 & 23 & 24 & \textcolor{green}{\(\uparrow\)+1} & 14/09 & 23.31 & 24.05 & \textcolor{green}{\(\uparrow\)+0.74} & 30/10 & 23.86 & 21.2 & \textcolor{red}{\(\downarrow\)-2.66} \\

& 
& 16/10 & 22.45 & 19/10 & 21.69 & 22.6 & \textcolor{red}{\(\downarrow\)-0.91} & 30/10 & 20.94 & 21.2 & \textcolor{red}{\(\downarrow\)-0.26} & 15/12 & 23.61 & 16.3 & \textcolor{red}{\(\downarrow\)-7.31} \\

& 
& 21/10 & 20.0 & 24/10 & 19.89 & 20.9 & \textcolor{red}{\(\downarrow\)-1.01} & 04/11 & 18.13 & 20.15 & \textcolor{red}{\(\downarrow\)-2.02} & 20/12 & 20.92 & 17.8 & \textcolor{red}{\(\downarrow\)-3.12} \\

&
& 22/10 & 20.5 & 25/10 & 17.66 & 20.9 & \textcolor{red}{\(\downarrow\)-3.24} & 05/11 & 20.94 & 20.1 & \textcolor{red}{\(\downarrow\)-0.84} & 21/12 & 21.35 & 17.8 & \textcolor{red}{\(\downarrow\)+3.55} \\
\midrule


\multirow{4}{*}{\textbf{DXG}} 
& \textbf{Round Robin}
& 31/08 & 22.8 & 03/09 & 23.07 & 24 & \textcolor{green}{\(\uparrow\)+0.93} & 14/09 & 19.38 & 24.05 & \textcolor{red}{\(\downarrow\)-4.67} & 30/10 & 24.01 & 21.2 & \textcolor{red}{\(\downarrow\)-2.81} \\

&
& 16/10 & 22.45 & 19/10 & 22.77 & 22.6 & \textcolor{green}{\(\uparrow\)+0.17} & 30/10 & 23.10 & 21.2 & \textcolor{red}{\(\downarrow\)-1.90} & 15/12 & 23.92 & 16.3 & \textcolor{red}{\(\downarrow\)-7.62} \\

&
& 21/10 & 20.0 & 24/10 & 20.17 & 20.9 & \textcolor{green}{\(\uparrow\)+0.73} & 04/11 & 18.34 & 20.15 & \textcolor{red}{\(\downarrow\)+1.81} & 20/12 & 20.74 & 17.8 & \textcolor{red}{\(\downarrow\)-2.94} \\

&
& 22/10 & 20.5 & 25/10 & 18.35 & 20.9 & \textcolor{red}{\(\downarrow\)-2.55} & 05/11 & 20.89 & 20.1 & \textcolor{red}{\(\downarrow\)-0.79} & 21/12 & 21.40 & 17.8 & \textcolor{red}{\(\downarrow\)-3.6} \\
\midrule


\multirow{4}{*}{\textbf{DXG}} 
& \textbf{Ensemble}
& 31/08 & 22.8 & 03/09 & 23.97 & 24 & \textcolor{green}{\(\uparrow\)+1.03} & 14/09 & 23.18 & 24.05 & \textcolor{green}{\(\uparrow\)+0.87} & 30/10 & 23.58 & 21.2 & \textcolor{red}{\(\downarrow\)-2.38} \\

&
& 16/10 & 22.45 & 19/10 & 22.65 & 22.6 & \textcolor{green}{\(\uparrow\)+0.05} & 30/10 & 22.91 & 21.2 & \textcolor{red}{\(\downarrow\)-1.71} & 15/12 & 23.33 & 16.3 & \textcolor{red}{\(\downarrow\)-7.03} \\

&
& 21/10 & 20.0 & 24/10 & 20.15 & 20.9 & \textcolor{green}{\(\uparrow\)+0.75} & 04/11 & 20.36 & 20.15 & \textcolor{green}{\(\uparrow\)+0.21} & 20/12 & 20.66 & 17.8 & \textcolor{red}{\(\downarrow\)-2.86} \\

&
& 22/10 & 20.5 & 25/10 & 20.65 & 20.9 & \textcolor{green}{\(\uparrow\)+0.25} & 05/11 & 20.79 & 20.1 & \textcolor{red}{\(\downarrow\)-0.69} & 21/12 & 21.19 & 17.8 & \textcolor{red}{\(\downarrow\)-3.39} \\
\midrule



%%%%%%%%%%%%%%%%%%%%%%%%%%%%%%%%%%%%%%%%%%%%%%%%%%%%%%%%%%%%%%%%%%%%%%%%%%%%%%%%%%%%
% ===================== FPT ======================
%%%%%%%%%%%%%%%%%%%%%%%%%%%%%%%%%%%%%%%%%%%%%%%%%%%%%%%%%%%%%%%%%%%%%%%%%%%%%%%%%%%%
\multirow{4}{*}{\textbf{FPT}} 
& \textbf{Hierarchical}
& 01/09 & 101.6 & 04/09 & 101.09 & 105 & \textcolor{red}{\(\downarrow\)-3.91} & 15/09 & 101.91 & 101.9 & \textcolor{green}{\(\uparrow\)+0.09} & 31/10 & 103.71 & 103.9 & \textcolor{green}{\(\uparrow\)+0.19} \\

&
& 16/10 & 89.8 & 19/10 & 89.47 & 88.1 & \textcolor{green}{\(\uparrow\)+0.91} & 30/10 & 90.01 & 102.7 & \textcolor{green}{\(\uparrow\)+12.69} & 15/12 & 91.69 & 93.8 & \textcolor{green}{\(\uparrow\)+2.11} \\

&
& 21/10 & 93 & 24/10 & 92.61 & 97.7 & \textcolor{red}{\(\downarrow\)-5.09} & 04/11 & 93.36 & 103.3 & \textcolor{green}{\(\uparrow\)+9.94} & 20/12 & 94.59 & 93.9 &\textcolor{green}{\(\uparrow\)+0.69}\\

&
& 22/10 & 97 & 25/10 & 96.6 & 97.7 & \textcolor{red}{\(\downarrow\)-1.1} & 05/11 & 97.22 & 100.9 & \textcolor{green}{\(\uparrow\)-3.68} & 21/12 & 98.65 & 93.9 & \textcolor{red}{\(\downarrow\)-4.75} \\
\midrule


\multirow{4}{*}{\textbf{FPT}} 
& \textbf{Round Robin}
& 01/09 & 101.6 & 04/09 & 101.16 & 105 & \textcolor{red}{\(\downarrow\)-3.84} & 15/09 & 102.01 & 101.9 & \textcolor{green}{\(\uparrow\)+0.01} & 31/10 & 104.16 & 103.9 & \textcolor{green}{\(\uparrow\)+0.26} \\

&
& 16/10 & 89.8 & 19/10 & 89.35 & 88.1 & \textcolor{green}{\(\uparrow\)+1.25} & 30/10 & 89.19 & 102.7 & \textcolor{red}{\(\downarrow\)-13.51} & 15/12 & 90.96 & 93.8 & \textcolor{green}{\(\uparrow\)+2.84} \\

&
& 21/10 & 93 & 24/10 & 92.59 & 97.7 & \textcolor{red}{\(\downarrow\)-5.11} & 04/11 & 93.29 & 103.3 & \textcolor{green}{\(\uparrow\)+10.01} & 20/12 & 94.52 & 93.9 & \textcolor{green}{\(\uparrow\)+0.62} \\

&
& 22/10 & 97 & 25/10 & 93.48 & 97.7 & \textcolor{red}{\(\downarrow\)-4.22} & 05/11 & 97.28 & 100.9 & \textcolor{green}{\(\uparrow\)+3.62} & 21/12 & 98.65 & 93.9 & \textcolor{red}{\(\downarrow\)-4.75} \\
\midrule



\multirow{4}{*}{\textbf{FPT}} 
& \textbf{Ensemble}
& 01/09 & 101.6 & 04/09 & 101.08 & 105 & \textcolor{green}{\(\uparrow\)+3.92} & 15/09 & 101.63 & 101.9 & \textcolor{green}{\(\uparrow\)+0.27} & 31/10 & 103.14 & 103.9 & \textcolor{green}{\(\uparrow\)+0.76} \\

&
& 16/10 & 89.8 & 19/10 & 89.37 & 88.1 & \textcolor{green}{\(\uparrow\)+1.29} & 30/10 & 101.88 & 102.7 & \textcolor{green}{\(\uparrow\)+0.82} & 15/12 & 93.09 & 93.8 & \textcolor{green}{\(\uparrow\)+0.71} \\

&
& 21/10 & 93 & 24/10 & 92.53 & 97.7 & \textcolor{red}{\(\downarrow\)-5.17} & 04/11 & 103.09 & 103.3 & \textcolor{green}{\(\uparrow\)+0.21} & 20/12 & 94.23 & 93.9 & \textcolor{green}{\(\uparrow\)+0.33} \\

&
& 22/10 & 97 & 25/10 & 96.48 & 97.7 & \textcolor{red}{\(\downarrow\)-1.22} & 05/11 & 96.52 & 100.9 & \textcolor{red}{\(\downarrow\)-4.38} & 21/11 & 98.15 & 93.9 & \textcolor{red}{\(\downarrow\)-4.25} \\
\midrule




%%%%%%%%%%%%%%%%%%%%%%%%%%%%%%%%%%%%%%%%%%%%%%%%%%%%%%%%%%%%%%%%%%%%%%%%%%%%%%%%%%%%
% ===================== VCB ======================
%%%%%%%%%%%%%%%%%%%%%%%%%%%%%%%%%%%%%%%%%%%%%%%%%%%%%%%%%%%%%%%%%%%%%%%%%%%%%%%%%%%%
\multirow{4}{*}{\textbf{VCB}} 
& \textbf{Hierarchical}
& 30/08 & 68.6 & 02/09 & 68.37 & 68.6 & \textcolor{yellow}{0.23} & 13/09 & 68.70 & 65.8 & \textcolor{red}{\(\downarrow\)+2.9} & 29/10 & 69.41 & 60.7 & \textcolor{red}{\(\downarrow\)-8.71} \\

&
& 16/10 & 62.9 & 19/10 & 63.19 & 61.9 & \textcolor{red}{\(\downarrow\)-1.29} & 30/10 & 63.43 & 60.6 & \textcolor{red}{\(\downarrow\)-2.83} & 15/12 & 64.05 & 56.8 & \textcolor{red}{\(\downarrow\)-7.25} \\

&
& 21/10 & 59.3 & 24/10 & 59.56 & 59.5 & \textcolor{green}{\(\uparrow\)+0.04} & 04/11 & 59.87 & 60.1 & \textcolor{green}{\(\uparrow\)+0.23} & 20/12 & 60.44 & 57.5 & \textcolor{red}{\(\downarrow\)-2.94} \\

&
& 22/10 & 59.6 & 25/10 & 61.88 & 59.5 & \textcolor{red}{\(\downarrow\)-2.38} & 05/11 & 60.18 & 60.8 & \textcolor{green}{\(\uparrow\)+0.62} & 21/12 & 58.79 & 57.5 & \textcolor{green}{\(\uparrow\)+1.29} \\
\midrule


\multirow{4}{*}{\textbf{VCB}} 
& \textbf{Round Robin}
& 30/08 & 68.6 & 02/09 & 68.44 & 68.6 & \textcolor{yellow}{0.16} & 13/09 & 68.89 & 65.8 & \textcolor{red}{\(\downarrow\)-3.10} & 29/10 & 69.99 & 60.7 & \textcolor{red}{\(\downarrow\)-9.29} \\

&
& 16/10 & 62.9 & 19/10 & 63.19 & 61.9 & \textcolor{green}{\(\uparrow\)+1.29} & 30/10 & 63.60 & 60.6 & \textcolor{red}{\(\downarrow\)-3} & 15/12 & 54.31 & 56.8 & \textcolor{green}{\(\uparrow\)-2.49} \\

&
& 21/10 & 59.3 & 24/10 & 59.53 & 59.5 & \textcolor{red}{\(\downarrow\)-0.03} & 04/11 & 59.79 & 60.1 & \textcolor{red}{\(\downarrow\)-0.31} & 20/12 & 60.37 & 57.5 & \textcolor{red}{\(\downarrow\)-2.87} \\

&
& 22/10 & 59.6 & 25/10 & 59.82 & 59.5 & \textcolor{red}{\(\downarrow\)-0.32} & 05/11 & 60.08 & 60.8 & \textcolor{green}{\(\uparrow\)+0.72} & 21/12 & 60.55 & 57.5 & \textcolor{red}{\(\downarrow\)-3.05} \\
\midrule

\multirow{4}{*}{\textbf{VCB}} 
& \textbf{Ensemble}
& 30/08 & 68.6 & 02/09 & 68.36 & 68.6 & \textcolor{yellow}{0.24} & 13/09 & 68.66 & 65.8 & \textcolor{red}{\(\downarrow\)-2.86} & 29/10 & 69.33 & 60.7 & \textcolor{red}{\(\downarrow\)-8.63} \\

&
& 16/10 & 62.9 & 19/10 & 63.14 & 61.9 & \textcolor{red}{\(\downarrow\)-1.24} & 30/10 & 63.41 & 60.6 & \textcolor{red}{\(\downarrow\)-2.51} & 15/12 & 55.87 & 56.8 & \textcolor{green}{\(\uparrow\)+0.93} \\

&
& 21/10 & 59.3 & 24/10 & 59.52 & 59.5 & \textcolor{green}{\(\uparrow\)+0.02} & 04/11 & 59.86 & 60.1 & \textcolor{green}{\(\uparrow\)+0.24} & 20/12 & 60.25 & 57.5 & \textcolor{red}{\(\downarrow\)-2.75} \\

&
& 22/10 & 59.6 & 25/10 & 59.83 & 59.5 & \textcolor{red}{\(\downarrow\)-0.33} & 05/11 & 60.11 & 60.8 & \textcolor{green}{\(\uparrow\)+0.69} & 21/12 & 60.69 & 57.5 & \textcolor{red}{\(\downarrow\)-3.19} \\
\midrule



%%%%%%%%%%%%%%%%%%%%%%%%%%%%%%%%%%%%%%%%%%%%%%%%%%%%%%%%%%%%%%%%%%%%%%%%%%%%%%%%%%%%
% ===================== VCS ======================
%%%%%%%%%%%%%%%%%%%%%%%%%%%%%%%%%%%%%%%%%%%%%%%%%%%%%%%%%%%%%%%%%%%%%%%%%%%%%%%%%%%%
\multirow{4}{*}{\textbf{VCS}} 
& \textbf{Hierarchical}
& 30/08 & 48.8 & 02/09 & 47.26 & 48.8 & \textcolor{yellow}{1.54} & 13/09 & 47.72 & 50 & \textcolor{red}{\(\downarrow\)-2.28} & 29/10 & 48.32 & 47.6 & \textcolor{green}{\(\uparrow\)+0.72} \\

&
& 16/10 & 47.2 & 19/10 & 45.67 & 46.8 & \textcolor{green}{\(\uparrow\)+1.13} & 30/10 & 46.07 & 47.4 & \textcolor{red}{\(\downarrow\)-1.33} & 15/12 & 46.59 & 46.5 & \textcolor{green}{\(\uparrow\)+0.09} \\

&
& 21/10 & 45.5 & 24/10 & 44.06 & 46.2 & \textcolor{red}{\(\downarrow\)-2.14} & 04/11 & 44.37 & 47.1 & \textcolor{red}{\(\downarrow\)-2.73} & 20/12 & 44.97 & 43 & \textcolor{green}{\(\uparrow\)+1.97} \\

&
& 22/10 & 45.5 & 25/10 & 44.05 & 46.2 & \textcolor{red}{\(\downarrow\)-0.15} & 05/11 & 42.31 & 46.9 & \textcolor{red}{\(\downarrow\)-4.59} & 21/12 & 45.11 & 43 & \textcolor{green}{\(\uparrow\)+2.11} \\
\midrule


\multirow{4}{*}{\textbf{VCS}} 
& \textbf{Round Robin}
& 30/08 & 48.8 & 02/09 & 47.23 & 48.8 & \textcolor{yellow}{1.57} & 13/09 & 47.58 & 50 & \textcolor{red}{\(\downarrow\)-2.42} & 29/10 & 48.05 & 47.6 & \textcolor{green}{\(\uparrow\)+0.45} \\

&
& 16/10 & 47.2 & 19/10 & 45.71 & 46.8 & \textcolor{green}{\(\uparrow\)-1.09} & 30/10 & 47.99 & 47.4 & \textcolor{green}{\(\uparrow\)+0.59} & 15/12 & 47.39 & 46.5 & \textcolor{red}{\(\downarrow\)-0.89} \\

&
& 21/10 & 45.5 & 24/10 & 43.98 & 46.2 & \textcolor{red}{\(\downarrow\)-2.22} & 04/11 & 48.21 & 47.1 & \textcolor{green}{\(\uparrow\)+1.11} & 20/12 & 45.88 & 43 & \textcolor{red}{\(\downarrow\)-2.88} \\

&
& 22/10 & 45.5 & 25/10 & 44.01 & 46.2 & \textcolor{red}{\(\downarrow\)-2.19} & 05/11 & 46.33 & 46.9 & \textcolor{green}{\(\uparrow\)+0.57} & 21/12 & 45.77 & 43 & \textcolor{red}{\(\downarrow\)-2.77} \\
\midrule



\multirow{4}{*}{\textbf{VCS}} 
& \textbf{Ensemble}
& 30/08 & 48.8 & 02/09 & 47.19 & 48.8 & \textcolor{yellow}{1.61} & 13/09 & 52.54 & 50 & \textcolor{green}{\(\uparrow\)+2.54} & 29/10 & 48.86 & 47.6 & \textcolor{red}{\(\downarrow\)-1.26} \\

&
& 16/10 & 47.2 & 19/10 & 45.64 & 46.8 & \textcolor{green}{\(\uparrow\)+1.16} & 30/10 & 47.16 & 47.4 & \textcolor{green}{\(\uparrow\)+0.24} & 15/12 & 46.57 & 46.5 & \textcolor{green}{\(\uparrow\)+0.07} \\

&
& 21/10 & 45.5 & 24/10 & 43.99 & 46.2 & \textcolor{red}{\(\downarrow\)-2.21} & 04/11 & 44.28 & 47.1 & \textcolor{red}{\(\downarrow\)-2.82} & 20/12 & 44.85 & 43 & \textcolor{green}{\(\uparrow\)+1.85} \\

&
& 22/10 & 45.5 & 25/10 & 43.97 & 46.2 & \textcolor{red}{\(\downarrow\)-2.23} & 05/11 & 44.28 & 46.9 & \textcolor{red}{\(\downarrow\)-2.62} & 21/12 & 44.9 & 43 & \textcolor{green}{\(\uparrow\)+1.9}  \\
\midrule


%%%%%%%%%%%%%%%%%%%%%%%%%%%%%%%%%%%%%%%%%%%%%%%%%%%%%%%%%%%%%%%%%%%%%%%%%%%%%%%%%%%%
% ===================== HPG ======================
%%%%%%%%%%%%%%%%%%%%%%%%%%%%%%%%%%%%%%%%%%%%%%%%%%%%%%%%%%%%%%%%%%%%%%%%%%%%%%%%%%%%
\multirow{4}{*}{\textbf{HPG}} 
& \textbf{Hierarchical}
& 01/09 & 27.5 & 04/09 & 27.67 & 29.85 & \textcolor{green}{\(\uparrow\)+2.18} & 15/09 & 27.83 & 30.35 & \textcolor{green}{\(\uparrow\)+2.52} & 31/10 & 28.17 & 26.7 & \textcolor{red}{\(\downarrow\)-1.47} \\

& 
& 16/10 & 28.25 & 19/10 & 28.44 & 28 & \textcolor{red}{\(\downarrow\)-0.44} & 30/10 & 28.68 & 26.9 & \textcolor{red}{\(\downarrow\)-1.78} & 15/12 & 29.12 & 26.25 & \textcolor{red}{\(\downarrow\)-2.87} \\

& 
& 21/10 & 26.8 & 24/10 & 26.94 & 26.4 & \textcolor{red}{\(\downarrow\)-0.54} & 04/11 & 27.1 & 26.75 & \textcolor{red}{\(\downarrow\)-0.35} & 20/12 & 26.42 & 26.7 & \textcolor{green}{\(\uparrow\)+0.28} \\

&
& 22/10 & 26.7 & 25/10 & 26.55 & 26.4 & \textcolor{green}{\(\uparrow\)+0.15} & 05/11 & 26.97 & 26.3 & \textcolor{red}{\(\downarrow\)-0.67} & 21/12 & 27.37 & 26.7 & \textcolor{yellow}{0.67} \\
\midrule


\multirow{4}{*}{\textbf{HPG}} 
& \textbf{Round Robin}
& 01/09 & 27.5 & 04/09 & 28.84 & 29.85 & \textcolor{green}{\(\uparrow\)+1.01} & 15/09 & 28.84 & 30.35 & \textcolor{green}{\(\uparrow\)+1.51} & 31/10 & 26.41 & 26.7 & \textcolor{green}{\(\uparrow\)+0.49} \\

& 
& 16/10 & 28.25 & 19/10 & 29.63 & 28 & \textcolor{red}{\(\downarrow\)-1.63} & 30/10 & 29.01 & 26.9 & \textcolor{red}{\(\downarrow\)-2.11} & 15/12 & 29.30 & 26.25 & \textcolor{red}{\(\downarrow\)+3.05} \\

& 
& 21/10 & 26.8 & 24/10 & 26.29 & 26.4 & \textcolor{green}{\(\uparrow\)+0.11} & 04/11 & 26.06 & 26.75 & \textcolor{green}{\(\uparrow\)+0.69} & 20/12 & 27.41 & 26.7 & \textcolor{red}{\(\downarrow\)-0.71} \\

&
& 22/10 & 26.7 & 25/10 & 26.79 & 26.4 & \textcolor{red}{\(\downarrow\)-0.39} & 05/11 & 26.89 & 26.3 & \textcolor{red}{\(\downarrow\)-0.59} & 21/12 & 27.23 & 26.7 & \textcolor{yellow}{0.53} \\
\midrule


\multirow{4}{*}{\textbf{HPG}} 
& \textbf{Ensemble}
& 01/09 & 27.5 & 04/09 & 27.65 & 29.85& \textcolor{green}{\(\uparrow\)+2.2} & 15/09 & 27.8 & 30.35 & \textcolor{green}{\(\uparrow\)+2.55} & 31/10 & 28.11 & 26.7 & \textcolor{red}{\(\downarrow\)-1.41} \\

& 
& 16/10 & 28.25 & 19/10 & 28.4 & 28 & \textcolor{red}{\(\downarrow\)-0.4} & 30/10 & 28.58 & 26.9 & \textcolor{red}{\(\downarrow\)-1.68} & 15/12 & 27.98 & 26.25 & \textcolor{green}{\(\uparrow\)+1.73} \\

& 
& 21/10 & 26.8 & 24/10 & 26.93 & 26.4 & \textcolor{red}{\(\downarrow\)-0.53} & 04/11 & 27.06 & 26.75 & \textcolor{red}{\(\downarrow\)-0.31} & 20/12 & 26.5 & 26.7 & \textcolor{green}{\(\uparrow\)+0.2} \\

&
& 22/10 & 26.7 & 25/10 & 26.83 & 26.4 & \textcolor{red}{\(\downarrow\)-0.43} & 05/11 & 26.5 & 26.3 & \textcolor{green}{\(\uparrow\)+0.2} & 21/12 & 27.35 & 26.7 & \textcolor{yellow}{0.65} \\
\midrule



%%%%%%%%%%%%%%%%%%%%%%%%%%%%%%%%%%%%%%%%%%%%%%%%%%%%%%%%%%%%%%%%%%%%%%%%%%%%%%%%%%%%
% ===================== TCB ======================
%%%%%%%%%%%%%%%%%%%%%%%%%%%%%%%%%%%%%%%%%%%%%%%%%%%%%%%%%%%%%%%%%%%%%%%%%%%%%%%%%%%%
\multirow{4}{*}{\textbf{TCB}} 
& \textbf{Ensemble}
& 01/09 & 39.6 & 04/09 & 38.78 & 39.6 & \textcolor{yellow}{0.82} & 15/09 & 39.03 & 39 & \textcolor{green}{\(\uparrow\)+0.03} & 31/10 & 39.43 & 35.1 & \textcolor{green}{\(\uparrow\)+4.33} \\

&
& 16/10 & 41.25 & 19/10 & 41.5 & 40.65 & \textcolor{red}{\(\downarrow\)-0.85} & 30/10 & 41.77 & 35.7 & \textcolor{red}{\(\downarrow\)-6.07} & 15/12 & 42.14 & 32 & \textcolor{red}{\(\downarrow\)-10.14} \\

&
& 21/10 & 37.65 & 24/10 & 37.88 & 36.1 & \textcolor{red}{\(\downarrow\)-1.78} & 04/11 & 38.08 & 35 & \textcolor{red}{\(\downarrow\)-3.08} & 20/12 & 38.43 & 33.6 &\textcolor{red}{\(\downarrow\)-4.83}\\

&
& 22/10 & 38 & 25/10 & 35.2 & 36.1 & \textcolor{green}{\(\uparrow\)+0.9} & 05/11 & 34.41 & 34.1 & \textcolor{green}{\(\uparrow\)+0.31} & 21/12 & 33.02 & 33.6 & \textcolor{green}{\(\uparrow\)+0.58} \\
\midrule


\multirow{4}{*}{\textbf{TCB}} 
& \textbf{Hierarchical}
& 01/09 & 39.6 & 04/09 & 38.84 & 39.6 & \textcolor{yellow}{0.76} & 15/09 & 39.09 & 39 & \textcolor{green}{\(\uparrow\)+0.09} & 31/10 & 39.68 & 35.1 & \textcolor{red}{\(\downarrow\)-4.58} \\

&
& 16/10 & 41.25 & 19/10 & 41.52 & 40.65 & \textcolor{red}{\(\downarrow\)-0.87} & 30/10 & 41.85 & 35.7 & \textcolor{red}{\(\downarrow\)-6.15} & 15/12 & 38.43 & 32 & \textcolor{green}{\(\uparrow\)+6.43} \\

&
& 21/10 & 37.65 & 24/10 & 37.88 & 36.1 & \textcolor{red}{\(\downarrow\)-1.78} & 04/11 & 35.23 & 35 & \textcolor{green}{\(\uparrow\)+0.23} & 20/12 & 38.52 & 33.6 &\textcolor{red}{\(\downarrow\)-4.92}\\

&
& 22/10 & 38 & 25/10 & 38.25 & 36.1 & \textcolor{red}{\(\downarrow\)-2.15} & 05/11 & 36.59 & 34.1 & \textcolor{green}{\(\uparrow\)+2.49} & 21/12 & 34.09 & 33.6 & \textcolor{green}{\(\uparrow\)+0.49} \\
\midrule



\multirow{4}{*}{\textbf{TCB}} 
& \textbf{RoundRobin}
& 01/09 & 39.6 & 04/09 & 38.85 & 39.6 & \textcolor{yellow}{0.75} & 15/09 & 39.23 & 39 & \textcolor{green}{\(\uparrow\)+0.23} & 31/10 & 39.88 & 35.1 & \textcolor{red}{\(\downarrow\)-4.78} \\

&
& 16/10 & 41.25 & 19/10 & 41.54 & 40.65 & \textcolor{red}{\(\downarrow\)-0.89} & 30/10 & 41.97 & 35.7 & \textcolor{red}{\(\downarrow\)-6.27} & 15/12 & 32.71 & 32 & \textcolor{green}{\(\uparrow\)+0.71} \\

&
& 21/10 & 37.65 & 24/10 & 36.27 & 36.1 & \textcolor{green}{\(\uparrow\)+0.17} & 04/11 & 38.12 & 35 & \textcolor{red}{\(\downarrow\)-3.12} & 20/12 & 38.6 & 33.6 &\textcolor{red}{\(\downarrow\)-5}\\

&
& 22/10 & 38 & 25/10 & 38.27 & 36.1 & \textcolor{red}{\(\downarrow\)-2.17} & 05/11 & 38.49 & 34.1 & \textcolor{red}{\(\downarrow\)-4.39} & 21/12 & 38.87 & 33.6 & \textcolor{red}{\(\downarrow\)-5.27} \\
\midrule




%%%%%%%%%%%%%%%%%%%%%%%%%%%%%%%%%%%%%%%%%%%%%%%%%%%%%%%%%%%%%%%%%%%%%%%%%%%%%%%%%%%%
% ===================== VHM ======================
%%%%%%%%%%%%%%%%%%%%%%%%%%%%%%%%%%%%%%%%%%%%%%%%%%%%%%%%%%%%%%%%%%%%%%%%%%%%%%%%%%%%
\multirow{4}{*}{\textbf{VHM}} 
& \textbf{Hierarchical}
& 01/09 & 104.5 & 04/09 & 105.28 & 102.3 & \textcolor{red}{\(\downarrow\)-2.98} & 15/09 & 106.31 & 104.5 & \textcolor{yellow}{1.81} & 31/10 & 97.97 & 99.2 & \textcolor{green}{\(\uparrow\)+1.23} \\

&
& 16/10 & 122 & 19/10 & 123.23 & 116 & \textcolor{red}{\(\downarrow\)-7.23} & 30/10 & 104.43 & 104 & \textcolor{green}{\(\uparrow\)+0.43} & 15/12 & 107 & 93.9 & \textcolor{green}{\(\uparrow\)+13.1} \\

&
& 21/10 & 110.9 & 24/10 & 111.77 & 114.5 & \textcolor{green}{\(\uparrow\)+2.73} & 04/11 & 112.53 & 100.2 & \textcolor{red}{\(\downarrow\)-12.33} & 20/12 & 114.76 & 101.5 & \textcolor{red}{\(\downarrow\)-13.26} \\

&
& 22/10 & 113 & 25/10 & 113.8 & 114.5 & \textcolor{green}{\(\uparrow\)+0.7} & 05/11 & 104.52 & 99.6 & \textcolor{green}{\(\uparrow\)+4.92} & 21/12 & 106.61 & 101.5 & \textcolor{green}{\(\uparrow\)+5.11} \\
\midrule


\multirow{4}{*}{\textbf{VHM}} 
& \textbf{Round Robin}
& 01/09 & 104.5 & 04/09 & 105.54 & 102.3 & \textcolor{red}{\(\downarrow\)-3.24} & 15/09 & 106.72 & 104.5 & \textcolor{yellow}{2.22} & 31/10 & 108.85 & 99.2 & \textcolor{red}{\(\downarrow\)-9.65} \\

&
& 16/10 & 122 & 19/10 & 123.89 & 116 & \textcolor{red}{\(\downarrow\)-7.89} & 30/10 & 105.79 & 104 & \textcolor{green}{\(\uparrow\)+1.79} & 15/12 & 100.95 & 93.9 & \textcolor{green}{\(\uparrow\)+7.05} \\

&
& 21/10 & 110.9 & 24/10 & 111.92 & 114.5 & \textcolor{green}{\(\uparrow\)+2.58} & 04/11 & 103 & 100.2 & \textcolor{green}{\(\uparrow\)+2.8} & 20/12 & 105.22 & 101.5 & \textcolor{green}{\(\uparrow\)+3.72} \\

&
& 22/10 & 113 & 25/10 & 113.88 & 114.5 & \textcolor{green}{\(\uparrow\)+0.62} & 05/11 & 104.64 & 99.6 & \textcolor{green}{\(\uparrow\)+5.04} & 21/12 & 106.08 & 101.5 & \textcolor{green}{\(\uparrow\)+4.58} \\
\midrule


\multirow{4}{*}{\textbf{VHM}} 
& \textbf{Ensemble}
& 01/09 & 104.5 & 04/09 & 101.14 & 102.3 & \textcolor{green}{\(\uparrow\)+1.16} & 15/09 & 106.03 & 104.5 & \textcolor{yellow}{1.53} & 31/10 & 97.62 & 99.2 & \textcolor{green}{\(\uparrow\)+1.58} \\

&
& 16/10 & 122 & 19/10 & 122.91 & 116 & \textcolor{red}{\(\downarrow\)-6.91} & 30/10 & 103.8 & 104 & \textcolor{green}{\(\uparrow\)+0.2} & 15/12 & 95.53 & 93.9 & \textcolor{green}{\(\uparrow\)+1.63} \\

&
& 21/10 & 110.9 & 24/10 & 111.53 & 114.5 & \textcolor{green}{\(\uparrow\)+2.97} & 04/11 & 102.32 & 100.2 & \textcolor{green}{\(\uparrow\)+2.12} & 20/12 & 104.03 & 101.5 & \textcolor{green}{\(\uparrow\)+2.53} \\

&
& 22/10 & 113 & 25/10 & 113.67 & 114.5 & \textcolor{green}{\(\uparrow\)+0.83} & 05/11 & 99.5 & 99.6 & \textcolor{green}{\(\uparrow\)+0.1} & 21/12 & 105.83 & 101.5 & \textcolor{green}{\(\uparrow\)+4.33} \\
\midrule




%%%%%%%%%%%%%%%%%%%%%%%%%%%%%%%%%%%%%%%%%%%%%%%%%%%%%%%%%%%%%%%%%%%%%%%%%%%%%%%%%%%%
% ===================== MWG ======================
%%%%%%%%%%%%%%%%%%%%%%%%%%%%%%%%%%%%%%%%%%%%%%%%%%%%%%%%%%%%%%%%%%%%%%%%%%%%%%%%%%%%
\multirow{4}{*}{\textbf{MWG}} 
& \textbf{Hierarchical}
& 01/09 & 78 & 04/09 & 78.54 & 77.5 & \textcolor{red}{\(\downarrow\)-1.04} & 15/09 & 79.14 & 79.5 & \textcolor{green}{\(\uparrow\)+0.36} & 31/10 & 80.53 & 82.6 & \textcolor{green}{\(\uparrow\)+2.07} \\

&
& 16/10 & 84.5 & 19/10 & 85.09 & 84.5 & \textcolor{yellow}{0.59} & 30/10 & 85.93 & 83.9 & \textcolor{red}{\(\downarrow\)-2.03} & 15/12 & 87.25 & 77.7 & \textcolor{red}{\(\downarrow\)-9.55} \\

&
& 21/10 & 83 & 24/10 & 83.53 & 85.7 & \textcolor{green}{\(\uparrow\)+2.17} & 04/11 & 84.26 & 81.8 & \textcolor{red}{\(\downarrow\)-2.46} & 20/12 & 85.63 & 82.9 & \textcolor{red}{\(\downarrow\)-2.73} \\

&
& 22/10 & 84.5 & 25/10 & 85.15 & 85.7 & \textcolor{green}{\(\uparrow\)+0.55} & 05/11 & 85.82 & 80.2 & \textcolor{red}{\(\downarrow\)-5.62} & 21/12 & 87.32 & 82.9 & \textcolor{red}{\(\downarrow\)-4.42} \\
\midrule


\multirow{4}{*}{\textbf{MWG}} 
& \textbf{Round Robin}
& 01/09 & 78 & 04/09 & 78.8 & 77.5 & \textcolor{red}{\(\downarrow\)-1.3} & 15/09 & 79.76 & 79.5 & \textcolor{green}{\(\uparrow\)+0.26} & 31/10 & 81.67 & 82.6 & \textcolor{green}{\(\uparrow\)+0.93} \\

&
& 16/10 & 84.5 & 19/10 & 85.31 & 84.5 & \textcolor{yellow}{0.81} & 30/10 & 86.44 & 83.9 & \textcolor{red}{\(\downarrow\)-2.54} & 15/12 & 85.75 & 77.7 & \textcolor{red}{\(\downarrow\)-8.05} \\

&
& 21/10 & 83 & 24/10 & 83.63 & 85.7 & \textcolor{green}{\(\uparrow\)+2.07s} & 04/11 & 84.31 & 81.8 & \textcolor{red}{\(\downarrow\)-2.51} & 20/12 & 81.22 & 82.9 & \textcolor{green}{\(\uparrow\)+1.68} \\

&
& 22/10 & 84.5 & 25/10 & 85.2 & 85.7 & \textcolor{green}{\(\uparrow\)+0.5} & 05/11 & 85.85 & 80.2 & \textcolor{red}{\(\downarrow\)-5.65} & 21/12 & 87.7 & 82.9 & \textcolor{green}{\(\uparrow\)+4.8} \\
\midrule




\multirow{4}{*}{\textbf{MWG}} 
& \textbf{Ensemble}
& 01/09 & 78 & 04/09 & 78.49 & 77.5 & \textcolor{red}{\(\downarrow\)-0.99} & 15/09 & 79 & 79.5 & \textcolor{green}{\(\uparrow\)+0.5} & 31/10 & 80.38 & 82.6 & \textcolor{green}{\(\uparrow\)+2.22} \\

&
& 16/10 & 84.5 & 19/10 & 85.06 & 84.5 & \textcolor{yellow}{0.56} & 30/10 & 85.64 & 83.9 & \textcolor{red}{\(\downarrow\)-1.74} & 15/12 & 76.79 & 77.7 & \textcolor{green}{\(\uparrow\)+0.91} \\

&
& 21/10 & 83 & 24/10 & 83.48 & 85.7 & \textcolor{green}{\(\uparrow\)+2.22} & 04/11 & 83.92 & 81.8 & \textcolor{green}{\(\uparrow\)+2.12} & 20/12 & 85.31 & 82.9 & \textcolor{red}{\(\downarrow\)-2.41} \\

&
& 22/10 & 84.5 & 25/10 & 84.95 & 85.7 & \textcolor{green}{\(\uparrow\)+0.75} & 05/11 & 85.52 & 80.2 & \textcolor{red}{\(\downarrow\)-5.32} & 21/12 & 86.68 & 82.9 & \textcolor{red}{\(\downarrow\)-3.78} \\




\midrule


%%%%%%%%%%%%%%%%%%%%%%%%%%%%%%%%%%%%%%%%%%%%%%%%%%%%%%%%%%%%%%%%%%%%%%%%%%%%%%%%%%%%
% ===================== CMG ======================
%%%%%%%%%%%%%%%%%%%%%%%%%%%%%%%%%%%%%%%%%%%%%%%%%%%%%%%%%%%%%%%%%%%%%%%%%%%%%%%%%%%%
\multirow{4}{*}{\textbf{CMG}} 
& \textbf{Hierarchical}
& 01/09 & 40.3 & 04/09 & 36.08 & 41.85 & \textcolor{red}{\(\downarrow\)-5.77} & 15/09 & 34.13 & 40.8 & \textcolor{red}{\(\downarrow\)-6.67} & 31/10 & 49.31 & 41.9 & \textcolor{green}{\(\uparrow\)+7.41} \\

& 
& 16/10 & 38 & 19/10 & 34 & 37.3 & \textcolor{red}{\(\downarrow\)-3.3} & 30/10 & 32.06 & 41.3 & \textcolor{red}{\(\downarrow\)-9.24} & 15/12 & 37.64 & 35.1 & \textcolor{green}{\(\uparrow\)+2.54} \\

& 
& 21/10 & 38.4 & 24/10 & 34.38 & 38.75 & \textcolor{red}{\(\downarrow\)-4.37} & 04/11 & 32.52 & 39.6 & \textcolor{red}{\(\downarrow\)-7.08} & 20/12 & 37.93 & 35.2 & \textcolor{green}{\(\uparrow\)+2.73} \\

&
& 22/10 & 38.5 & 25/10 & 39.12 & 38.75 & \textcolor{green}{\(\uparrow\)+0.37} & 05/11 & 32.22 & 39.3 & \textcolor{red}{\(\downarrow\)-7.08} & 21/12 & 38 & 35.2 & \textcolor{green}{\(\uparrow\)+2.8} \\
\midrule


\multirow{4}{*}{\textbf{CMG}} 
& \textbf{Round Robin}
& 01/09 & 40.3 & 04/09 & 35.39 & 41.85 & \textcolor{red}{\(\downarrow\)-6.46} & 15/09 & 41.18 & 40.8 & \textcolor{green}{\(\uparrow\)+0.38} & 31/10 & 39.31 & 41.9 & \textcolor{red}{\(\downarrow\)-2.59} \\

& 
& 16/10 & 38 & 19/10 & 33.39 & 37.3 & \textcolor{green}{\(\uparrow\)+3.91} & 30/10 & 39.46 & 41.3 & \textcolor{green}{\(\uparrow\)+1.84} & 15/12 & 27.64 & 35.1 & \textcolor{green}{\(\uparrow\)+7.46} \\

& 
& 21/10 & 38.4 & 24/10 & 33.52 & 38.75 & \textcolor{red}{\(\downarrow\)-5.23} & 04/11 & 38.88 & 39.6 & \textcolor{green}{\(\uparrow\)+0.72} & 20/12 & 27.93 & 35.2 & \textcolor{red}{\(\downarrow\)-7.27} \\

&
& 22/10 & 38.5 & 25/10 & 33.57 & 38.75 & \textcolor{red}{\(\downarrow\)-5.18} & 05/11 & 38.22 & 39.3 & \textcolor{red}{\(\downarrow\)-1.08} & 21/12 & 36.74 & 35.2 & \textcolor{green}{\(\uparrow\)+1.54} \\
\midrule


\multirow{4}{*}{\textbf{CMG}} 
& \textbf{Ensemble}
& 01/09 & 40.3 & 04/09 & 41.94 & 41.85 & \textcolor{green}{\(\uparrow\)+0.09} & 15/09 & 33.51 & 40.8 & \textcolor{red}{\(\downarrow\)-7.29} & 31/10 & 42.88 & 41.9 & \textcolor{green}{\(\uparrow\)+0.98} \\

& 
& 16/10 & 38 & 19/10 & 33.97 & 37.3 & \textcolor{red}{\(\downarrow\)-3.33} & 30/10 & 41.94 & 41.3 & \textcolor{green}{\(\uparrow\)+0.64} & 15/12 & 36.78 & 35.1 & \textcolor{green}{\(\uparrow\)+1.68} \\

& 
& 21/10 & 38.4 & 24/10 & 34.36 & 38.75 & \textcolor{red}{\(\downarrow\)-4.39} & 04/11 & 39.42 & 39.6 & \textcolor{green}{\(\uparrow\)+0.18} & 20/12 & 39.52 & 35.2 & \textcolor{red}{\(\downarrow\)-4.32} \\

&
& 22/10 & 38.5 & 25/10 & 34.47 & 38.75 & \textcolor{red}{\(\downarrow\)-4.28} & 05/11 & 39.59 & 39.3 & \textcolor{green}{\(\uparrow\)+0.29} & 21/12 & 39.13 & 35.2 & \textcolor{red}{\(\downarrow\)-3.93} \\
\midrule



%%%%%%%%%%%%%%%%%%%%%%%%%%%%%%%%%%%%%%%%%%%%%%%%%%%%%%%%%%%%%%%%%%%%%%%%%%%%%%%%%%%%
% ===================== KDH======================
%%%%%%%%%%%%%%%%%%%%%%%%%%%%%%%%%%%%%%%%%%%%%%%%%%%%%%%%%%%%%%%%%%%%%%%%%%%%%%%%%%%%
\multirow{4}{*}{\textbf{KDH}} 
& \textbf{Hierarchical}
& 01/09 & 36.5 & 04/09 & 35.84 & 36.75 & \textcolor{red}{\(\downarrow\)-0.91} & 15/09 & 33.52 & 34.85 & \textcolor{green}{\(\uparrow\)+1.33} & 31/10 & 39.2 & 35.85 & \textcolor{red}{\(\downarrow\)-3.35} \\

&
& 16/10 & 34.2 & 19/10 & 33.59 & 33.9 & \textcolor{green}{\(\uparrow\)+0.31} & 30/10 & 31.46 & 35.85 & \textcolor{red}{\(\downarrow\)-4.39} & 15/12 & 27.36 & 29.6 & \textcolor{green}{\(\uparrow\)+2.24} \\

&
& 21/10 & 31.75 & 24/10 & 31.19 & 33.8 & \textcolor{red}{\(\downarrow\)-2.61} & 04/11 & 29.22 & 35.6 & \textcolor{red}{\(\downarrow\)-6.38} & 20/12 & 35.21 & 32.65 &\textcolor{green}{\(\uparrow\)+2.56}\\

&
& 22/10 & 32.2 & 25/10 & 31.16 & 33.8 & \textcolor{red}{\(\downarrow\)-2.64} & 05/11 & 34.52 & 34.95 & \textcolor{green}{\(\uparrow\)+0.43} & 21/12 & 35.56 & 32.65 & \textcolor{green}{\(\uparrow\)+2.91} \\
\midrule


\multirow{4}{*}{\textbf{KDH}} 
& \textbf{Round Robin}
& 01/09 & 36.5 & 04/09 & 35.53 & 36.75 & \textcolor{red}{\(\downarrow\)-1.22} & 15/09 & 32.2 & 34.85 & \textcolor{green}{\(\uparrow\)+2.65} & 31/10 & 35.88 & 35.85 & \textcolor{green}{\(\uparrow\)+0.03} \\

&
& 16/10 & 34.2 & 19/10 & 33.12 & 33.9 & \textcolor{green}{\(\uparrow\)+0.78} & 30/10 & 30.45 & 35.85 & \textcolor{red}{\(\downarrow\)-5.4} & 15/12 & 29.12 & 29.6 & \textcolor{green}{\(\uparrow\)+0.48} \\

&
& 21/10 & 31.75 & 24/10 & 33.67 & 33.8 & \textcolor{green}{\(\uparrow\)+0.13} & 04/11 & 37.01 & 35.6 & \textcolor{green}{\(\uparrow\)+1.41} & 20/12 & 25.4 & 32.65 &\textcolor{red}{\(\downarrow\)-7.25}\\

&
& 22/10 & 32.2 & 25/10 & 31 & 33.8 & \textcolor{red}{\(\downarrow\)-2.8} & 05/11 & 36.98 & 34.95 & \textcolor{green}{\(\uparrow\)+2.03} & 21/12 & 35.76 & 32.65 & \textcolor{green}{\(\uparrow\)-3.11} \\
\midrule



\multirow{4}{*}{\textbf{KDH}} 
& \textbf{Ensemble}
& 01/09 & 36.5 & 04/09 & 35.62 & 36.75 & \textcolor{green}{\(\uparrow\)+1.13} & 15/09 & 32.59 & 34.85 & \textcolor{green}{\(\uparrow\)+2.26} & 31/10 & 35.32 & 35.85 & \textcolor{green}{\(\uparrow\)+0.53} \\

&
& 16/10 & 34.2 & 19/10 & 33.48 & 33.9 & \textcolor{green}{\(\uparrow\)+0.42} & 30/10 & 35.96 & 35.85 & \textcolor{green}{\(\uparrow\)+0.11} & 15/12 & 30.02 & 29.6 & \textcolor{green}{\(\uparrow\)+0.42} \\

&
& 21/10 & 31.75 & 24/10 & 31.12 & 33.8 & \textcolor{red}{\(\downarrow\)-2.68} & 04/11 & 36.9 & 35.6 & \textcolor{green}{\(\uparrow\)+1.3} & 20/12 & 34.05 & 32.65 &\textcolor{green}{\(\uparrow\)+1.4}\\

&
& 22/10 & 32.2 & 25/10 & 31.58 & 33.8 & \textcolor{red}{\(\downarrow\)-2.22} & 05/11 & 34.41 & 34.95 & \textcolor{green}{\(\uparrow\)+0.54} & 21/12 & 32.85 & 32.65 & \textcolor{green}{\(\uparrow\)+0.2} \\
\midrule




%%%%%%%%%%%%%%%%%%%%%%%%%%%%%%%%%%%%%%%%%%%%%%%%%%%%%%%%%%%%%%%%%%%%%%%%%%%%%%%%%%%%
% ===================== DGC ======================
%%%%%%%%%%%%%%%%%%%%%%%%%%%%%%%%%%%%%%%%%%%%%%%%%%%%%%%%%%%%%%%%%%%%%%%%%%%%%%%%%%%%
\multirow{4}{*}{\textbf{DGC}} 
& \textbf{Hierarchical}
& 01/09 & 98.2 & 04/09 & 99.9 & 99.3 & \textcolor{green}{\(\uparrow\)+0.6} & 15/09 & 94.33 & 99.5 & \textcolor{red}{\(\downarrow\)-5.17} & 31/10 & 95.27 & 96 & \textcolor{green}{\(\uparrow\)+0.73} \\

&
& 16/10 & 95 & 19/10 & 88.47 & 93.2 & \textcolor{green}{\(\uparrow\)+4.73} & 30/10 & 99.7 & 93.6 & \textcolor{red}{\(\downarrow\)-6.1} & 15/12 & 92.82 & 93 & \textcolor{green}{\(\uparrow\)+0.18} \\

&
& 21/10 & 89.2 & 24/10 & 82.43 & 92.4 & \textcolor{red}{\(\downarrow\)-9.97} & 04/11 & 92.19 & 94.6 & \textcolor{green}{\(\uparrow\)+2.41} & 20/12 & 68.37 & 70.2 & \textcolor{green}{\(\uparrow\)+1.83} \\

&
& 22/10 & 91 & 25/10 & 93.48 & 92.4 & \textcolor{green}{\(\uparrow\)+1.08} & 05/11 & 91.18 & 95 & \textcolor{green}{\(\uparrow\)+3.82} & 21/12 & 69.75 & 70.2 & \textcolor{green}{\(\uparrow\)+0.45} \\
\midrule


\multirow{4}{*}{\textbf{DGC}} 
& \textbf{Round Robin}
& 01/09 & 98.2 & 04/09 & 98.86 & 99.3 & \textcolor{green}{\(\uparrow\)+0.44} & 15/09 & 95.27 & 99.5 & \textcolor{red}{\(\downarrow\)-4.23} & 31/10 & 100.02 & 96 & \textcolor{red}{\(\downarrow\)-4.02} \\

&
& 16/10 & 95 & 19/10 & 86.2 & 93.2 & \textcolor{green}{\(\uparrow\)+7} & 30/10 & 92.82 & 93.6 & \textcolor{green}{\(\uparrow\)+0.78} & 15/12 & 96.33 & 93 & \textcolor{red}{\(\downarrow\)-3.33} \\

&
& 21/10 & 89.2 & 24/10 & 91.14 & 92.4 & \textcolor{green}{\(\uparrow\)+1.26} & 04/11 & 98.37 & 94.6 & \textcolor{green}{\(\uparrow\)+3.77} & 20/12 & 68.37 & 70.2 & \textcolor{green}{\(\uparrow\)+1.83} \\

&
& 22/10 & 91 & 25/10 & 83.09 & 92.4 & \textcolor{red}{\(\downarrow\)-9.31} & 05/11 & 99.75 & 95 & \textcolor{green}{\(\uparrow\)+4.75} & 21/12 & 70.35 & 70.2 & \textcolor{green}{\(\downarrow\)-0.15} \\
\midrule


\multirow{4}{*}{\textbf{DGC}} 
& \textbf{Ensemble}
& 01/09 & 98.2 & 04/09 & 99.12 & 99.3 & \textcolor{red}{\(\downarrow\)-0.18} & 15/09 & 99.71 & 99.5 & \textcolor{green}{\(\uparrow\)+0.21} & 31/10 & 95.74 & 96 & \textcolor{red}{\(\downarrow\)-0.26} \\

&
& 16/10 & 95 & 19/10 & 93.41 & 93.2 & \textcolor{green}{\(\uparrow\)+0.21} & 30/10 & 93.38 & 93.6 & \textcolor{red}{\(\downarrow\)-0.22} & 15/12 & 93.15 & 93 & \textcolor{green}{\(\uparrow\)+0.15} \\

&
& 21/10 & 89.2 & 24/10 & 92.15 & 92.4 & \textcolor{red}{\(\downarrow\)-0.25} & 04/11 & 94.82 & 94.6 & \textcolor{green}{\(\uparrow\)+0.22} & 20/12 & 70.85 & 70.2 & \textcolor{green}{\(\uparrow\)+0.65} \\

&
& 22/10 & 91 & 25/10 & 92.51 & 92.4 & \textcolor{green}{\(\uparrow\)+0.11} & 05/11 & 94.87 & 95 & \textcolor{red}{\(\downarrow\)-0.13} & 21/12 & 69.84 & 70.2 & \textcolor{red}{\(\downarrow\)-0.36} \\
\midrule



%%%%%%%%%%%%%%%%%%%%%%%%%%%%%%%%%%%%%%%%%%%%%%%%%%%%%%%%%%%%%%%%%%%%%%%%%%%%%%%%%%%%
% ===================== MBB ======================
%%%%%%%%%%%%%%%%%%%%%%%%%%%%%%%%%%%%%%%%%%%%%%%%%%%%%%%%%%%%%%%%%%%%%%%%%%%%%%%%%%%%
\multirow{4}{*}{\textbf{MBB}} 
& \textbf{Hierarchical}
& 01/09 & 27.75 & 04/09 & 28.15 & 28.25 & \textcolor{red}{\(\downarrow\)-0.10} & 15/09 & 27.05 & 26.85 & \textcolor{green}{\(\uparrow\)+0.20} & 31/10 & 23.35 & 23.6 & \textcolor{red}{\(\downarrow\)-0.25} \\

&
& 16/10 & 27.2 & 19/10 & 27.35 & 27.1 & \textcolor{green}{\(\uparrow\)+0.25} & 30/10 & 23.68 & 23.95 & \textcolor{red}{\(\downarrow\)-0.27} & 15/12 & 24.02 & 23.75 & \textcolor{green}{\(\uparrow\)+0.27} \\

&
& 21/10 & 25.3 & 24/10 & 24.18 & 24.4 & \textcolor{red}{\(\downarrow\)-0.22} & 04/11 & 24.28 & 24 & \textcolor{green}{\(\uparrow\)+0.28} & 20/12 & 24.48 & 24.7 & \textcolor{red}{\(\downarrow\)-0.22} \\

&
& 22/10 & 25.45 & 25/10 & 24.62 & 24.4 & \textcolor{green}{\(\uparrow\)+0.22} & 05/11 & 23.65 & 23.9 & \textcolor{red}{\(\downarrow\)-0.25} & 21/12 & 24.95 & 24.7 & \textcolor{green}{\(\uparrow\)+0.25} \\
\midrule


\multirow{4}{*}{\textbf{MBB}} 
& \textbf{Round Robin}
& 01/09 & 27.75 & 04/09 & 28.48 & 28.25 & \textcolor{green}{\(\uparrow\)+0.23} & 15/09 & 26.62 & 26.85 & \textcolor{red}{\(\downarrow\)-0.23} & 31/10 & 23.82 & 23.6 & \textcolor{green}{\(\uparrow\)+0.22} \\

&
& 16/10 & 27.2 & 19/10 & 26.88 & 27.1 & \textcolor{red}{\(\downarrow\)-0.22} & 30/10 & 24.18 & 23.95 & \textcolor{green}{\(\uparrow\)+0.23} & 15/12 & 23.52 & 23.75 & \textcolor{red}{\(\downarrow\)-0.23} \\

&
& 21/10 & 25.3 & 24/10 & 24.62 & 24.4 & \textcolor{green}{\(\uparrow\)+0.22} & 04/11 & 23.78 & 24 & \textcolor{red}{\(\downarrow\)-0.22} & 20/12 & 24.92 & 24.7 & \textcolor{green}{\(\uparrow\)+0.22} \\

&
& 22/10 & 25.45 & 25/10 & 24.18 & 24.4 & \textcolor{red}{\(\downarrow\)-0.22} & 05/11 & 24.12 & 23.9 & \textcolor{green}{\(\uparrow\)+0.22} & 21/12 & 24.48 & 24.7 & \textcolor{red}{\(\downarrow\)-0.22} \\
\midrule



\multirow{4}{*}{\textbf{MBB}} 
& \textbf{Ensemble}
& 01/09 & 27.75 & 04/09 & 28.35 & 28.25 & \textcolor{green}{\(\uparrow\)+0.10} & 15/09 & 26.75 & 26.85 & \textcolor{red}{\(\downarrow\)-0.10} & 31/10 & 23.70 & 23.6 & \textcolor{green}{\(\uparrow\)+0.10} \\

&
& 16/10 & 27.2 & 19/10 & 27.00 & 27.1 & \textcolor{red}{\(\downarrow\)-0.10} & 30/10 & 24.05 & 23.95 & \textcolor{green}{\(\uparrow\)+0.10} & 15/12 & 23.65 & 23.75 & \textcolor{red}{\(\downarrow\)-0.10} \\

&
& 21/10 & 25.3 & 24/10 & 24.50 & 24.4 & \textcolor{green}{\(\uparrow\)+0.10} & 04/11 & 23.90 & 24 & \textcolor{red}{\(\downarrow\)-0.10} & 20/12 & 24.80 & 24.7 & \textcolor{green}{\(\uparrow\)+0.10} \\

&
& 22/10 & 25.45 & 25/10 & 24.30 & 24.4 & \textcolor{red}{\(\downarrow\)-0.10} & 05/11 & 24.00 & 23.9 & \textcolor{green}{\(\uparrow\)+0.10} & 21/12 & 24.60 & 24.7 & \textcolor{red}{\(\downarrow\)-0.10} \\


\bottomrule
\end{longtable}
}







\clearpage
\onecolumn
\section{RQ2: Impact of Foundation LLMs}
\begin{table*}[!h]
\centering
\caption{Performance Comparison Across Forecasting Horizons and LLMs}
\label{tab:full_performance_all}
\setlength{\tabcolsep}{3pt}
\renewcommand{\arraystretch}{1.05}
\scriptsize
\begin{tabular}{llcccc|ccc|ccc}
\toprule
\multirow{2}{*}{\textbf{Horizon}} &
\multirow{2}{*}{\textbf{Stock}} &
\multirow{2}{*}{\textbf{Architecture}} &
\multicolumn{3}{c|}{\textbf{Llama-3.1-8B}} &
\multicolumn{3}{c|}{\textbf{GPT-4o}} &
\multicolumn{3}{c}{\textbf{Gemini 2.0}} \\
\cmidrule(lr){4-6} \cmidrule(lr){7-9} \cmidrule(lr){10-12}
& & & \textbf{MAE} & \textbf{RMSE} & \textbf{MAPE} & \textbf{MAE} & \textbf{RMSE} & \textbf{MAPE} & \textbf{MAE} & \textbf{RMSE} & \textbf{MAPE} \\
\midrule

% SHORT
\multirow{12}{*}{\textbf{Short-term}}
& \multirow{3}{*}{\textbf{DXG}} & Hierarchical & 0.54 & 0.68 & 2.4 & 0.50 & 0.62 & 2.3 & 0.71 & 0.94 & 3.1 \\
& & Round Robin & \textbf{0.51} & \textbf{0.61} & \textbf{2.3} & \textbf{0.50} & \textbf{0.60} & \textbf{2.3} & 0.79 & 0.97 & 3.5 \\
& & Ensemble & 0.55 & 0.69 & 2.5 & 0.51 & 0.64 & 2.3 & \textbf{0.68} & \textbf{0.80} & \textbf{1.5} \\
& \multirow{3}{*}{\textbf{FPT}} & Hierarchical & 2.82 & 3.29 & 2.9 & \textbf{2.04} & \textbf{2.83} & \textbf{2.1} & 2.67 & 2.85 & 2.8 \\
& & Round Robin & \textbf{2.79} & \textbf{3.27} & \textbf{2.8} & 2.75 & 3.23 & 2.8 & 2.16 & 2.50 & 2.3 \\
& & Ensemble & 2.90 & 3.38 & 2.9 & 2.73 & 3.20 & 2.8 & \textbf{1.38} & \textbf{1.58} & \textbf{1.4} \\
& \multirow{3}{*}{\textbf{VCB}} & Hierarchical & 1.43 & 2.06 & 2.4 & 0.58 & 0.73 & 0.9 & 1.34 & 1.55 & 2.2 \\
& & Round Robin & \textbf{0.51} & \textbf{0.70} & \textbf{0.8} & \textbf{0.26} & \textbf{0.34} & \textbf{0.4} & \textbf{1.09} & \textbf{1.32} & \textbf{1.8} \\
& & Ensemble & 9.46 & 13.40 & 14.4 & 0.46 & 0.65 & 0.7 & 1.79 & 1.96 & 2.9 \\
& \multirow{3}{*}{\textbf{VCS}} & Hierarchical & \textbf{1.73} & \textbf{1.79} & \textbf{3.7} & 1.76 & 1.81 & 3.8 & 0.75 & 0.82 & 1.6 \\
& & Round Robin & 1.80 & 1.85 & 3.8 & 1.79 & 1.85 & 3.8 & 0.72 & 0.81 & \textbf{1.5} \\
& & Ensemble & 2.26 & 2.40 & 4.8 & \textbf{1.31} & \textbf{1.50} & \textbf{2.8} & \textbf{0.68} & \textbf{0.80} & \textbf{1.5} \\
\midrule

% MEDIUM
\multirow{12}{*}{\textbf{Medium-term}}
& \multirow{3}{*}{\textbf{DXG}} & Hierarchical & 0.84 & 0.99 & 3.9 & 0.86 & 1.03 & 4.0 & 2.99 & 4.53 & 14.6 \\
& & Round Robin & 0.87 & 1.07 & 4.1 & 0.88 & 1.08 & 4.1 & \textbf{1.84} & \textbf{2.16} & \textbf{8.7} \\
& & Ensemble & \textbf{0.82} & \textbf{0.97} & \textbf{3.8} & \textbf{0.76} & \textbf{0.95} & \textbf{3.5} & 2.06 & 3.19 & 10.0 \\
& \multirow{3}{*}{\textbf{FPT}} & Hierarchical & 6.66 & 8.33 & 6.5 & \textbf{6.58} & \textbf{8.31} & \textbf{6.4} & 7.09 & 8.12 & 6.9 \\
& & Round Robin & \textbf{6.64} & \textbf{8.32} & \textbf{6.5} & 6.64 & 8.36 & 6.5 & 7.63 & 8.10 & 7.5 \\
& & Ensemble & 7.15 & 8.48 & 7.0 & 6.72 & 8.36 & 6.5 & \textbf{5.81} & \textbf{6.97} & \textbf{5.7} \\
& \multirow{3}{*}{\textbf{VCB}} & Hierarchical & \textbf{1.16} & \textbf{1.58} & \textbf{1.8} & \textbf{1.61} & \textbf{1.80} & \textbf{2.6} & 2.56 & 2.68 & 4.1 \\
& & Round Robin & 1.67 & 2.23 & 2.7 & 1.88 & 2.32 & 3.0 & \textbf{2.49} & \textbf{2.64} & \textbf{4.0} \\
& & Ensemble & 1.63 & 2.01 & 2.6 & 2.24 & 2.49 & 3.6 & 2.51 & 2.71 & 4.0 \\
& \multirow{3}{*}{\textbf{VCS}} & Hierarchical & 2.23 & 2.30 & 4.7 & 2.20 & 2.26 & 4.6 & \textbf{0.74} & \textbf{1.19} & \textbf{1.6} \\
& & Round Robin & \textbf{2.22} & \textbf{2.27} & \textbf{4.6} & \textbf{1.65} & \textbf{1.95} & \textbf{3.4} & 1.38 & 1.98 & 2.8 \\
& & Ensemble & 2.30 & 2.37 & 4.8 & 1.85 & 2.05 & 3.8 & 0.75 & 0.94 & 1.6 \\
\midrule

% LONG
\multirow{12}{*}{\textbf{Long-term}}
& \multirow{3}{*}{\textbf{DXG}} & Hierarchical & 4.07 & 4.45 & 23.2 & 3.83 & 4.29 & 21.8 & \textbf{3.60} & \textbf{4.55} & \textbf{21.2} \\
& & Round Robin & 4.27 & 4.67 & 24.3 & 4.19 & 4.63 & 23.9 & 4.73 & 5.42 & 27.0 \\
& & Ensemble & \textbf{3.89} & \textbf{4.28} & \textbf{22.2} & \textbf{3.86} & \textbf{4.29} & \textbf{22.1} & 3.25 & 4.06 & 19.1 \\
& \multirow{3}{*}{\textbf{FPT}} & Hierarchical & 3.12 & 3.61 & 3.3 & 2.24 & 2.86 & 2.4 & 5.50 & 7.82 & 5.8 \\
& & Round Robin & 3.70 & 5.88 & 3.9 & \textbf{1.48} & \textbf{1.99} & \textbf{1.6} & 5.52 & 7.88 & 5.9 \\
& & Ensemble & \textbf{2.06} & \textbf{2.73} & \textbf{2.2} & 2.11 & 2.68 & 2.2 & \textbf{4.36} & \textbf{6.03} & \textbf{4.6} \\
& \multirow{3}{*}{\textbf{VCB}} & Hierarchical & 6.34 & 6.80 & 10.8 & 5.26 & 5.95 & 9.0 & \textbf{5.21} & \textbf{5.81} & \textbf{8.9} \\
& & Round Robin & 6.60 & 7.05 & 11.3 & \textbf{5.09} & \textbf{5.75} & \textbf{8.7} & 8.83 & 10.10 & 15.0 \\
& & Ensemble & \textbf{5.95} & \textbf{6.45} & \textbf{10.2} & 4.37 & 5.03 & 7.4 & 5.66 & 5.74 & 9.7 \\
& \multirow{3}{*}{\textbf{VCS}} & Hierarchical & 1.32 & 1.52 & 3.0 & 1.33 & 1.54 & 3.0 & 4.12 & 4.19 & 9.2 \\
& & Round Robin & \textbf{1.04} & \textbf{1.27} & \textbf{2.4} & \textbf{1.07} & \textbf{1.29} & \textbf{2.4} & \textbf{1.55} & \textbf{1.74} & \textbf{3.4} \\
& & Ensemble & 1.19 & 1.43 & 2.7 & 1.20 & 1.42 & 2.7 & 3.79 & 3.97 & 8.5 \\

\bottomrule
\end{tabular}

\vspace{0.3cm}
\footnotesize
\textbf{Note:} Bold values indicate best performance within each LLM group.
\end{table*}

\noindent\textbf{Conclusion:} The experimental results indicate that LLM performance in stock forecasting is highly dependent on the prediction horizon and market structure. While GPT-4o demonstrates strong accuracy in short-term forecasting, particularly for volatile stocks, Google Gemini 2.0 exhibits greater suitability for medium- and long-term predictions, achieving more stable error behavior and consistent performance across different coordination architectures. These characteristics suggest that Gemini 2.0 is better aligned with trend-oriented and strategic forecasting applications, whereas short-term accuracy-driven tasks may benefit more from GPT-4o.
\newpage
\twocolumn


\subsubsection{Short-term Performance}
\begin{table}[!h]
\centering
\caption{Short-term Performance of Prediction Architectures Across Stocks (LLM: Google Gemini 2.0 Flash)}
\resizebox{\columnwidth}{!}{
\begin{tabular}{c c c c c}
\hline
\textbf{Stock} & \textbf{Architecture} & \textbf{MAE} & \textbf{RMSE} & \textbf{MAPE} \\
\hline

\multirow{3}{*}{\textbf{DXG}}
& Hierarchical    & \textbf{0.71} & \textbf{0.94} & \textbf{3.1\%} \\
& RoundRobin      & 0.79 & 0.97 & 3.5\% \\
& Ensemble Voting & 0.89 & 1.09 & 3.9\% \\
\hline

\multirow{3}{*}{\textbf{FPT}}
& Ensemble Voting & \textbf{1.38} & \textbf{1.58} & \textbf{1.4\%} \\
& RoundRobin      & 2.16 & 2.50 & 2.3\% \\
& Hierarchical    & 2.67 & 2.85 & 2.8\% \\
\hline

\multirow{3}{*}{\textbf{VCB}}
& RoundRobin      & \textbf{1.09} & \textbf{1.32} & \textbf{1.8\%} \\
& Hierarchical    & 1.34 & 1.55 & 2.2\% \\
& Ensemble Voting & 1.79 & 1.96 & 2.9\% \\
\hline

\multirow{3}{*}{\textbf{VCS}}
& Ensemble Voting & \textbf{0.68} & \textbf{0.80} & \textbf{1.5\%} \\
& RoundRobin      & 0.72 & 0.81 & 1.5\% \\
& Hierarchical    & 0.75 & 0.82 & 1.6\% \\
\hline

\end{tabular}
}
\end{table}

\subsubsection{Medium-term Performance}
\begin{table}[!h]
\centering
\caption{Medium-term Performance of Prediction Architectures Across Stocks (LLM: Google Gemini 2.0 Flash)}
\resizebox{\columnwidth}{!}{
\begin{tabular}{c c c c c}
\hline
\textbf{Stock} & \textbf{Architecture} & \textbf{MAE} & \textbf{RMSE} & \textbf{MAPE} \\
\hline

\multirow{3}{*}{\textbf{DXG}}
& RoundRobin      & \textbf{1.84} & \textbf{2.16} & \textbf{8.7\%} \\
& Ensemble Voting & 2.06 & 3.19 & 10.0\% \\
& Hierarchical    & 2.99 & 4.53 & 14.6\% \\
\hline

\multirow{3}{*}{\textbf{FPT}}
& Ensemble Voting & \textbf{5.81} & \textbf{6.97} & \textbf{5.7\%} \\
& Hierarchical    & 7.09 & 8.12 & 6.9\% \\
& RoundRobin      & 7.63 & 8.10 & 7.5\% \\
\hline

\multirow{3}{*}{\textbf{VCB}}
& RoundRobin      & \textbf{2.49} & \textbf{2.64} & \textbf{4.0\%} \\
& Ensemble Voting & 2.51 & 2.71 & 4.0\% \\
& Hierarchical    & 2.56 & 2.68 & 4.1\% \\
\hline

\multirow{3}{*}{\textbf{VCS}}
& Hierarchical    & \textbf{0.74} & \textbf{1.19} & \textbf{1.6\%} \\
& Ensemble Voting & 0.75 & 0.94 & 1.6\% \\
& RoundRobin      & 1.38 & 1.98 & 2.8\% \\
\hline

\end{tabular}
}
\end{table}

\subsubsection{Long-term Performance}
\begin{table}[!h]
\centering
\caption{Long-term Performance of Prediction Architectures Across Stocks (LLM: Google Gemini 2.0 Flash)}
\resizebox{\columnwidth}{!}{
\begin{tabular}{c c c c c}
\hline
\textbf{Stock} & \textbf{Architecture} & \textbf{MAE} & \textbf{RMSE} & \textbf{MAPE} \\
\hline

\multirow{3}{*}{\textbf{DXG}}
& Ensemble Voting & \textbf{3.25} & \textbf{4.06} & \textbf{19.1\%} \\
& Hierarchical    & 3.60 & 4.55 & 21.2\% \\
& RoundRobin      & 4.73 & 5.42 & 27.0\% \\
\hline

\multirow{3}{*}{\textbf{FPT}}
& Ensemble Voting & \textbf{4.36} & \textbf{6.03} & \textbf{4.6\%} \\
& Hierarchical    & 5.50 & 7.82 & 5.8\% \\
& RoundRobin      & 5.52 & 7.88 & 5.9\% \\
\hline

\multirow{3}{*}{\textbf{VCB}}
& Hierarchical    & \textbf{5.21} & \textbf{5.81} & \textbf{8.9\%} \\
& Ensemble Voting & 5.66 & 5.74 & 9.7\% \\
& RoundRobin      & 8.83 & 10.10 & 15.0\% \\
\hline

\multirow{3}{*}{\textbf{VCS}}
& RoundRobin      & \textbf{1.55} & \textbf{1.74} & \textbf{3.4\%} \\
& Ensemble Voting & 3.79 & 3.97 & 8.5\% \\
& Hierarchical    & 4.12 & 4.19 & 9.2\% \\
\hline

\end{tabular}
}
\end{table}

\subsection{Conclusion}
Short-term (3-day) forecasts reveal stock-specific strengths: Hierarchical leads in volatile DXG (MAPE 3.1\%), Ensemble Voting in FPT (1.4\%) and VCS (1.5\%), and Round Robin in VCB (1.8\%). Medium-term (14-day) predictions highlight Ensemble Voting's consistent superiority (MAPE 1.6\%–10.0\%), with the lowest error in FPT (5.7\%), joint-best in VCB (4.0\%), near-parity in VCS (1.6\%), and competitive results in DXG (10.0\%). Overall, Ensemble Voting proves the most reliable and balanced architecture for practical forecasting across diverse market conditions.
\subsection{Visualization of Forecast Results}

\begin{figure}[!h]
    \centering
    \begin{subfigure}[b]{0.48\columnwidth}
        \centering
        \includegraphics[width=\linewidth,height=3.5cm]{anh/main/dxg.png}
        \caption{DXG}
        \label{fig:short_dxg}
    \end{subfigure}
    \hfill
    \begin{subfigure}[b]{0.48\columnwidth}
        \centering
        \includegraphics[width=\linewidth,height=3.5cm]{anh/main/fpt.png}
        \caption{FPT}
        \label{fig:short_fpt}
    \end{subfigure}

    \vspace{2mm}

    \begin{subfigure}[b]{0.48\columnwidth}
        \centering
        \includegraphics[width=\linewidth,height=3.5cm]{anh/main/vcb.png}
        \caption{VCB}
        \label{fig:short_vcb}
    \end{subfigure}
    \hfill
    \begin{subfigure}[b]{0.48\columnwidth}
        \centering
        \includegraphics[width=\linewidth,height=3.5cm]{anh/main/vcs.png}
        \caption{VCS}
        \label{fig:short_vcs}
    \end{subfigure}
    
    \caption{Short-term Forecast Comparison Across Stocks}
    \label{fig:short_forecast_all}
\end{figure}

\begin{figure}[!h]
    \centering
    \begin{subfigure}[b]{0.48\columnwidth}
        \centering
        \includegraphics[width=\linewidth,height=3.5cm]{anh/main/dxgmedium.png}
        \caption{DXG}
        \label{fig:medium_dxg}
    \end{subfigure}
    \hfill
    \begin{subfigure}[b]{0.48\columnwidth}
        \centering
        \includegraphics[width=\linewidth,height=3.5cm]{anh/main/fptmedium.png}
        \caption{FPT}
        \label{fig:medium_fpt}
    \end{subfigure}

    \vspace{2mm}

    \begin{subfigure}[b]{0.48\columnwidth}
        \centering
        \includegraphics[width=\linewidth,height=3.5cm]{anh/main/vcbmedium.png}
        \caption{VCB}
        \label{fig:medium_vcb}
    \end{subfigure}
    \hfill
    \begin{subfigure}[b]{0.48\columnwidth}
        \centering
        \includegraphics[width=\linewidth,height=3.5cm]{anh/main/vcsmedium.png}
        \caption{VCS}
        \label{fig:medium_vcs}
    \end{subfigure}
    
    \caption{Medium-term Forecast Comparison Across Stocks}
    \label{fig:medium_forecast_all}
\end{figure}
\newpage
\begin{figure}[!h]
    \centering
    \begin{subfigure}[b]{0.48\columnwidth}
        \centering
        \includegraphics[width=\linewidth,height=3.5cm]{anh/main/dxglong.png}
        \caption{DXG}
        \label{fig:long_dxg}
    \end{subfigure}
    \hfill
    \begin{subfigure}[b]{0.48\columnwidth}
        \centering
        \includegraphics[width=\linewidth,height=3.5cm]{anh/main/fptlong.png}
        \caption{FPT}
        \label{fig:long_fpt}
    \end{subfigure}

    \vspace{2mm}

    \begin{subfigure}[b]{0.48\columnwidth}
        \centering
        \includegraphics[width=\linewidth,height=3.5cm]{anh/main/vcblong.png}
        \caption{VCB}
        \label{fig:long_vcb}
    \end{subfigure}
    \hfill
    \begin{subfigure}[b]{0.48\columnwidth}
        \centering
        \includegraphics[width=\linewidth,height=3.5cm]{anh/main/vcslong.png}
        \caption{VCS}
        \label{fig:long_vcs}
    \end{subfigure}
    
    \caption{Long-term Forecast Comparison Across Stocks}
    \label{fig:long_forecast_all}
\end{figure}

\section{RQ3: Multi-Agent vs. Traditional Baselines}
To address RQ3, we compare multi-agent architectures against single-agent LLM (Core: Gemini 2.0 Flash) and LSTM baselines using medium-term (14-day) forecasts validated with actual closing prices on January 6, 2026 (VCB: 57,300 VND, FPT: 94,000 VND, VCS: 43,000 VND, DXG: 16,650 VND).

\begin{figure}[H]
    \centering
    \begin{subfigure}[b]{0.48\columnwidth}
        \centering
        \includegraphics[width=\linewidth,height=3.2cm]{anh/rq3/dxg.png}
        \caption{DXG}
        \label{fig:single_dxg}
    \end{subfigure}
    \hfill
    \begin{subfigure}[b]{0.48\columnwidth}
        \centering
        \includegraphics[width=\linewidth,height=3.2cm]{anh/rq3/fpt.png}
        \caption{FPT}
        \label{fig:single_fpt}
    \end{subfigure}

    \vspace{2mm}

    \begin{subfigure}[b]{0.48\columnwidth}
        \centering
        \includegraphics[width=\linewidth,height=3.2cm]{anh/rq3/vcb.png}
        \caption{VCB}
        \label{fig:single_vcb}
    \end{subfigure}
    \hfill
    \begin{subfigure}[b]{0.48\columnwidth}
        \centering
        \includegraphics[width=\linewidth,height=3.2cm]{anh/rq3/vcs.png}
        \caption{VCS}
        \label{fig:single_vcs}
    \end{subfigure}

    \caption{Comparative Framework: 14-day Forecasting Results Across Different Approaches}
    \label{fig:single_agent_2x2}
\end{figure}

\begin{table}[htbp]
    \centering
\caption{Detailed Forecast Evaluation Results by Stock (highlighting the VCB exception where Single-Agent outperforms MAS)}
    \label{tab:stock_prediction_results}
    \begin{tabular}{llrr}
        \toprule
        \textbf{Stock (Actual)} & \textbf{Method} & \textbf{MAE} & \textbf{MAPE (\%)} \\
        \midrule
        
        % --- VCB ---
        \multirow{5}{*}{\textbf{VCB} (57.3)} 
         & Single-Agent & \textbf{1.10} & \textbf{1.92} \\
         & Ensemble & 1.47 & 2.57 \\
         & Hierarchical & 1.74 & 3.04 \\
         & RoundRobin & 1.65 & 2.88 \\
         & Base LSTM & 3.09 & 5.39 \\
        \cmidrule{1-4}
        
        % --- FPT ---
        \multirow{5}{*}{\textbf{FPT} (94.0)} 
         & Single-Agent & 1.80 & 1.91 \\
         & Ensemble & 2.82 & 3.00 \\
         & Hierarchical & \textbf{1.37} & \textbf{1.46} \\
         & RoundRobin & 1.94 & 2.06 \\
         & Base LSTM & 2.89 & 3.07 \\
        \cmidrule{1-4}
        
        % --- VCS ---
        \multirow{5}{*}{\textbf{VCS} (43.0)} 
         & Single-Agent & 4.20 & 9.77 \\
         & Ensemble & 1.01 & 2.35 \\
         & Hierarchical & \textbf{0.29} & \textbf{0.67} \\
         & RoundRobin & 2.58 & 6.00 \\
         & Base LSTM & 6.54 & 15.21 \\
        \cmidrule{1-4}
        
        % --- DXG ---
        \multirow{5}{*}{\textbf{DXG} (16.65)} 
         & Single-Agent & 1.80 & 10.81 \\
         & Ensemble & 0.97 & 5.83 \\
         & Hierarchical & \textbf{0.81} & \textbf{4.86} \\
         & RoundRobin & 2.41 & 14.47 \\
         & Base LSTM & 1.28 & 7.69 \\
         
        \bottomrule
    \end{tabular}
\end{table}


\subsection{Empirical Validation Results}

\paragraph{Quantitative Analysis of RQ3 Results by Stock.}

In the relatively stable banking sector (VCB), \textit{Single-Agent} achieved the lowest error (MAE: 1.10; MAPE: 1.92\%), outperforming all MAS variants (\textit{Ensemble}: 2.57\%, \textit{Round Robin}: 2.88\%, \textit{Hierarchical}: 3.04\%) and substantially improving over \textit{Base LSTM} (5.39\%). This ranking indicates that under smoother, policy-anchored price dynamics, a lower-complexity decision pipeline can preserve signal purity more effectively than multi-source aggregation.


In the technology sector (FPT), the \textit{Hierarchical} architecture demonstrated the highest predictive accuracy, achieving a MAPE of 1.46\% and an MAE of 1.37. While the \textit{Single-Agent} (MAPE: 1.91\%) and \textit{Round Robin} (MAPE: 2.06\%) models also outperformed the \textit{Base LSTM} baseline (MAPE: 3.07\%), the Hierarchical structure proved superior in capturing the sustained growth momentum and trend persistence characteristic of technology stocks.


The most pronounced performance disparity was observed in the industrial construction sector (VCS). The \textit{Hierarchical} model delivered near-perfect precision, recording an exceptionally low MAPE of 0.67\% and an MAE of 0.29. In stark contrast, the \textit{Base LSTM} model exhibited catastrophic degradation, with a MAPE of 15.21\% and an MAE of 6.54, while the \textit{Single-Agent} model also struggled, yielding a MAPE of 9.77\%. These findings underscore the Hierarchical architecture’s unique capacity to accommodate complex noise structures and regime-dependent dynamics in industrial stocks, where conventional deep learning models fail to generalize.


The real estate sector (DXG), characterized by elevated volatility and strong sentiment dependence, posed the greatest forecasting challenge. In this environment, the \textit{Hierarchical} (MAPE: 4.86\%) and \textit{Ensemble} (MAPE: 5.83\%) architectures demonstrated superior robustness and risk management capabilities. Conversely, the \textit{Single-Agent} (MAPE: 10.81\%) and \textit{Round Robin} (MAPE: 14.47\%) approaches suffered from substantially higher error rates, indicating that in the absence of hierarchical coordination or ensemble consensus mechanisms, simpler architectures struggle to adapt to abrupt sentiment-driven market fluctuations.

Cross-sector comparison confirms that MAS gains increase with structural noise. In VCS, \textit{Hierarchical} reduces MAPE to 0.67\% versus 9.77\% (\textit{Single-Agent}) and 15.21\% (\textit{Base LSTM}). In DXG, \textit{Hierarchical} and \textit{Ensemble} achieve 4.86\% and 5.83\%, compared with 10.81\% (\textit{Single-Agent}). In FPT, \textit{Hierarchical} remains best at 1.46\% versus 1.91\% (\textit{Single-Agent}). The VCB pattern is therefore a sector-specific boundary condition: when market behavior is more stable and less fragmented, MAS may absorb additional variance from heterogeneous agent signals; when volatility and sentiment noise rise, MAS coordination yields clear error reduction.



\paragraph{Data-Driven Explanation of MAS Effectiveness Across Market Regimes.}
The strong performance of the Multi-Agent framework, particularly the \textit{Hierarchical} architecture, can be attributed to its alignment with key structural characteristics of Vietnam’s emerging stock market. 

First, the market is predominantly \textbf{retail-driven}, with individual investors accounting for approximately 85--90\% of trading volume. The Multi-Agent System (MAS) explicitly integrates sentiment and news-driven agents, enabling it to capture behavioral dynamics that are largely overlooked by traditional LSTM models relying solely on historical price sequences. 

Second, Vietnam’s market exhibits \textbf{high volatility and frequent regime shifts}. While Single-Agent models may perform competitively under relatively stable conditions, as observed in the Banking sector (VCB), coordinated agent interactions within the MAS provide inherent error-correction and consensus mechanisms that are crucial for volatile stocks such as DXG and VCS. 

Third, the market is characterized by pronounced \textbf{information heterogeneity}. The proposed MAS synthesizes diverse information sources, including technical indicators, fundamental metrics, risk profiles, and macroeconomic signals. In contrast, LSTM-based baselines are restricted to modeling temporal dependencies in historical price data, limiting their adaptability to sudden informational shocks. 

Finally, \textbf{sector-specific dynamics} play a critical role in forecasting performance. Specialized agents within the MAS are able to adapt to regulatory constraints in the banking sector, earnings-driven cycles in technology stocks, and credit policy sensitivities in real estate. This specialization enables more consistent performance across sectors where generalist forecasting models tend to fail.

\subsection{Conclusion}
RQ3 evidence supports a nuanced conclusion: Multi-Agent architectures are \textbf{generally superior} to traditional baselines and often outperform Single-Agent settings, but not universally. MAS clearly dominates in high-noise sectors (notably DXG and VCS), and remains competitive in banking; however, VCB is a meaningful exception where \textit{Single-Agent} achieves the best point accuracy. By embedding deliberative reasoning, cross-validation, and consensus-building, MAS improves robustness, stability, and interpretability at the market level, while the VCB case highlights that for structurally stable, policy-anchored banking dynamics, additional agent diversity can introduce mild coordination noise. This exception strengthens the study by showing that architecture selection should be sector-aware rather than one-size-fits-all.

\section{Discussion}
\subsection{Key Findings and Implications}
Empirical evaluation reveals strong horizon-dependent architectural effectiveness. Ensemble Voting demonstrates consistent superiority in medium-term forecasting (14-day) across multiple stocks: FPT (5.7\%), VCB (4.0\%), VCS (1.6\%), DXG (10.0\%), validating ensemble learning theory~\cite{dietterich2000ensemble}. Short-term predictions (3-day) exhibit stock-specific architectural preferences with no single dominant architecture, suggesting idiosyncratic factors require tailored approaches~\cite{campbell1997stock}. Long-term predictions (60-day) show mixed preferences: Hierarchical for VCB (8.9\%), Ensemble for DXG (19.1\%) and FPT (4.6\%), Round Robin for VCS (3.4\%), reflecting increased uncertainty in extended horizons~\cite{guidolin2007asset}.

Integrating Large Language Models with specialized financial agents enables sophisticated interpretation of unstructured information while maintaining domain-specific expertise~\cite{kenton2019bert,tetlock2007giving}.

\subsection{Theoretical Contributions}
This study provides empirical evidence for horizon-dependent coordination effectiveness, demonstrating that optimal architectural choices must consider temporal prediction horizons~\cite{hansen2001neural}. The dynamic weight adjustment mechanism provides a framework for adaptive ensemble learning in non-stationary financial environments~\cite{gama2014survey}. Integrating LLMs with specialized financial agents represents a novel hybrid architecture bridging symbolic AI and connectionist learning paradigms~\cite{bommasani2021foundation,brown2020language}. The modular agent architecture enables transparent attribution of predictions to specific analytical components, supporting regulatory compliance~\cite{arrieta2020explainable}.

\subsection{Research Question Analysis}

\subsubsection{A. RQ1: Horizon-Dependent Architectural Effectiveness}
No single coordination architecture is universally optimal; effectiveness is strictly horizon-dependent.

\paragraph{Medium-Term (14 Days): Ensemble Dominance.}
\textit{Ensemble Voting} consistently demonstrated superior performance: FPT (5.7\%), VCB (4.0\%), VCS (1.6\%), DXG (10.0\%), effectively integrating divergent technical and fundamental signals.

\paragraph{Short-Term (3 Days): Stock-Specific Heterogeneity.}
Hierarchical architecture proved advantageous in high-volatility environments, with centralized coordination mitigating speculative noise within tight settlement windows.

\paragraph{Long-Term (60 Days): Sector-Specific Specialization.}
Preferred architecture varied by sector: Hierarchical for regulated banking, Round Robin for cyclical industrial construction.

\subsubsection{B. RQ2: Impact of Foundation Models (LLMs)}
Optimal LLM selection is highly stock-specific and horizon-dependent. GPT-4o demonstrated superior short-term accuracy for DXG, FPT, and VCB. Medium-term forecasting revealed balanced distribution: Gemini 2.0 Flash for FPT and VCS, GPT-4o for DXG, Llama-3.1-8B for VCB. Long-term forecasting showed no single dominant LLM, suggesting practitioners should adopt a \textit{portfolio approach} to LLM selection.

\subsubsection{C. RQ3: Multi-Agent Systems vs. Traditional Baselines (with a sectoral exception)}
\textit{Multi-Agent Systems} (MAS) significantly outperform LSTM baselines and outperform Single-Agent configurations in most high-volatility contexts, but not in every sector. LSTM models exhibited significant stochastic instability with pronounced amplitude distortion. Single-Agent architectures lacked robust error-correction under high volatility (e.g., DXG and VCS), where MAS acted as \textit{stability anchors} by aggregating heterogeneous perspectives and filtering idiosyncratic noise. However, VCB provides an important exception: in a relatively stable, policy-anchored banking regime, \textit{Single-Agent} achieved the best MAE/MAPE, suggesting that when signal structure is smoother, MAS coordination can introduce additional variance from cross-agent disagreement.





\section{Limitations and Future Work}
\subsection{Current Limitations}
Several limitations must be acknowledged when interpreting these results. First, the evaluation scope encompasses only four stocks from the Vietnamese market over a specific temporal period, which may constrain the generalizability of findings to other emerging markets or broader equity universes. While the selected stocks represent key economic sectors, the limited sample size may not capture the full spectrum of market behaviors and sectoral dynamics present in Vietnam's broader equity landscape~\cite{phan2015stock}.

Second, the study focuses primarily on price prediction accuracy metrics (MAE, RMSE, MAPE) without incorporating transaction costs, liquidity constraints, or market impact considerations that are crucial for practical trading implementation. Vietnam's $T+2.5$ settlement period and relatively constrained liquidity compared to developed markets introduce practical constraints not fully reflected in the evaluation framework. Additionally, the impact of bid-ask spreads, market depth variations, and execution timing on actual trading performance remains unexplored.

\subsection{Temporal Isolation and News Data Leakage Prevention}
A critical methodological consideration in financial forecasting systems that incorporate news-driven agents is the risk of \textit{temporal data leakage}, wherein future information inadvertently influences historical predictions. In our framework, the TickerNews and MarketNews agents process company-specific and macroeconomic news streams to extract sentiment signals and event-driven impact assessments. To ensure temporal validity and prevent look-ahead bias, we implement strict temporal isolation protocols:

\textbf{Timestamp-based filtering:} All news articles, earnings releases, and policy announcements are timestamped at ingestion. During backtesting, each agent receives only news items with publication timestamps strictly preceding the prediction cutoff time $t$. For instance, when generating a 14-day forecast at market close on date $t$, only news published before the close of trading day $t$ is accessible to the agents.

\textbf{Intraday timing constraints:} Vietnamese stock exchanges (HOSE, HNX) operate with specific trading sessions (morning: 9:00--11:30, afternoon: 13:00--15:00). News released during trading hours on day $t$ may already be partially reflected in the closing price of day $t$. To address this, we apply a conservative cutoff: only news published before the opening of day $t$ is used for predictions targeting day $t$ or later, ensuring that sentiment signals do not incorporate information contemporaneous with or posterior to the target price.

\textbf{Embargo period for fundamental data:} Earnings reports, financial statements, and regulatory filings often exhibit publication delays. We enforce a 24-hour embargo period after official release timestamps before incorporating such data into the InvestmentExpert agent's fundamental analysis, mitigating the risk of using information that market participants could not have acted upon in real time.

\textbf{Validation through forward-only backtesting:} The entire evaluation framework employs a strictly forward-rolling window methodology. At each prediction timestamp $t$, the system state (including news corpus, technical indicators, and fundamental metrics) is reconstructed using only data available up to $t-1$, with no access to future information. This walk-forward validation ensures that reported performance metrics reflect realistic out-of-sample forecasting conditions.

Despite these safeguards, residual risks remain. First, the precise timing of information dissemination in Vietnam's market is not always deterministic; informal channels (social media, investor forums) may propagate rumors or unofficial news before formal publication, creating ambiguity in timestamp reliability. Second, the LLM-based reasoning layer may inadvertently encode temporal patterns learned during pretraining on datasets that include future information relative to our evaluation period, though this risk is mitigated by the use of structured prompt templates that explicitly condition on historical context only.

Future work should explore more granular timestamp validation, incorporation of intraday tick-level data to better align news timing with price formation, and adversarial testing frameworks to detect potential leakage pathways in multi-agent coordination layers.

\subsection{Operational and Robustness Limitations}
The framework's reliance on cloud-based LLM services introduces potential operational risks, including latency variability, service availability concerns, and cost scalability issues that could affect real-time trading applications~\cite{qiu2020pre}. Network connectivity disruptions, API rate limitations, and service outages could compromise system reliability in production environments. Furthermore, the cost structure of cloud-based inference may become prohibitive for high-frequency trading applications or for resource-constrained institutional investors.

Third, the framework's performance during extreme market conditions, such as financial crises, high-impact events, or periods of extreme volatility, remains untested. The evaluation period, while including normal market volatility, may not encompass the full range of stress scenarios that could challenge the system's robustness. The coordination mechanisms' behavior during periods of extreme correlation breakdown, liquidity crises, or regulatory interventions requires further investigation to establish operational reliability under adverse conditions~\cite{cont2001empirical}.

\section{Conclusion}

We investigate the problem of stock price forecasting in Vietnam’s emerging and retail-dominated market through the validation of a specialized \textit{Multi-Agent} forecasting framework. Based on a comprehensive empirical evaluation conducted on four representative stocks spanning multiple sectors (VCB, FPT, DXG, and VCS), the study provides clear and structured answers to its guiding research questions.

First, in addressing \textbf{RQ1 (Architectural Effectiveness Across Horizons)}, the results demonstrate that no single coordination architecture is universally optimal across all forecasting horizons. Instead, architectural effectiveness is strongly horizon-dependent. For the medium-term horizon (14 days), \textit{Ensemble Voting} consistently achieves robust performance across multiple stocks: FPT (5.7\% MAPE), VCB (4.0\% MAPE), VCS (1.6\% MAPE), and DXG (10.0\% MAPE), effectively reconciling heterogeneous technical, fundamental, and sentiment-driven signals. In contrast, short-term forecasting (3 days), particularly for highly volatile assets, benefits from a \textit{Hierarchical} coordination mechanism, in which centralized decision-making suppresses speculative noise within the $T+2.5$ settlement cycle. For long-term horizons (60 days), the optimal architecture diverges by sector, with \textit{Hierarchical} coordination proving effective for highly regulated industries such as banking (VCB: 8.9\% MAPE), \textit{Ensemble Voting} for real estate (DXG: 19.1\%) and technology (FPT: 4.6\%), and \textit{Round Robin} for industrial construction (VCS: 3.4\%).

Second, in response to \textbf{RQ2 (Impact of Foundation Models)}, the study establishes that the choice of the underlying Large Language Model (LLM) is highly context-dependent rather than universally applicable. LLM performance varies significantly across forecasting horizons and stock characteristics. For short-term forecasting, GPT-4o demonstrates strong performance for DXG and VCB, while Google Gemini 2.0 Flash proves optimal for FPT and VCS. Medium- and long-term horizons reveal a heterogeneous landscape in which Google Gemini 2.0 Flash, GPT-4o, and Llama-3.1-8B each emerge as optimal under different stock-specific conditions. These findings challenge the assumption of a single dominant LLM for financial forecasting and instead support a portfolio-based approach to LLM selection, in which models are empirically matched to specific assets and forecasting horizons.

Third, in answering \textbf{RQ3 (Superiority of Multi-Agent Systems over Traditional Baselines)}, the empirical results support a conditional rather than absolute claim. \textit{Multi-Agent Systems} consistently outperform traditional LSTM baselines and outperform Single-Agent LLMs in most volatile sectors, but VCB is a clear exception where \textit{Single-Agent} records the lowest error (MAE 1.10; MAPE 1.92\%). LSTM models exhibit pronounced instability under regime shifts, frequently overshooting during bullish phases and undershooting during bearish corrections, reflecting their limited capacity to incorporate market sentiment and structural changes. Single-Agent LLMs, while effective for relatively stable assets, lack robust error-correction under high volatility. In contrast, Multi-Agent architectures—particularly Hierarchical and Ensemble systems—act as stability anchors in sentiment-driven sectors by aggregating diverse perspectives from fundamental, technical, risk, and news-oriented agents. The VCB case therefore serves as a valuable boundary condition: in stable banking dynamics, simpler coordination may preserve cleaner signals than multi-source aggregation.

In summary, this work contributes a \textit{horizon-adaptive Multi-Agent forecasting framework} to the field of financial artificial intelligence. The findings indicate that within emerging markets such as Vietnam, predictive accuracy depends not solely on model complexity but on the strategic alignment between coordination architecture, forecasting horizon, and the underlying cognitive engine. This insight provides practical guidance for both researchers and practitioners seeking robust and adaptive forecasting solutions in volatile, retail-driven financial environments.

\section*{ACKNOWLEDGMENT}
The authors would like to thank the School of Computing at Phenikaa University for providing the computational resources and research environment necessary for this study. We also acknowledge the valuable feedback from reviewers that helped improve the quality of this manuscript.

\bibliographystyle{IEEEtran}
\bibliography{refs}



\begin{IEEEbiography}[{
\includegraphics[width=1in,height=1.25in,clip,keepaspectratio]{anh/tacgia/duong.jpg}}]{Duong Nguyen Minh} is currently pursuing the B.S. degree in Information Technology at the School of Computing, Phenikaa University, Ha Noi, Vietnam. His research interests include artificial intelligence, multi-agent systems, machine learning applications in financial markets, and deep learning for time series forecasting. He has contributed to research projects focusing on the application of AI-driven methodologies to emerging market analysis, particularly in the context of Vietnam's retail-dominated stock market.
\end{IEEEbiography}

\begin{IEEEbiography}[{
\includegraphics[width=1in,height=1.25in,clip,keepaspectratio]{anh/tacgia/dat.jpg}}]{Dat Nguyen Huu} is a Lecturer at the School of Computer Science, Phenikaa University, Hanoi, Vietnam. He received the M.Sc. degree in Applied Mathematics and Informatics from Hanoi University of Science and Technology. His research interests include artificial intelligence, large language models, natural language processing, computer vision, and distributed AI systems.
\end{IEEEbiography}


\EOD
\end{document}




